
%%FILE blueprint/src/chapter/100_percent.tex
\chapter{The 100\% version of PFR}

\begin{definition}[Symmetry group]\label{sym-group-def}
\lean{symmGroup}\leanok
If $X$ is a $G$-valued random variable, then the symmetry group $\mathrm{Sym}[X]$ is the set of all $h \in G$ such that $X +h$ has the same distribution as $X$.
\end{definition}

\begin{lemma}[Symmetry group is a group]\label{sym-group}\uses{sym-group-def}
\lean{symmGroup}\leanok
If $X$ is a $G$-valued random variable, then $\mathrm{Sym}[X]$ is a subgroup of $G$.
\end{lemma}

\begin{proof}\leanok Direct verification of the group axioms.
\end{proof}

\begin{lemma}[Zero Ruzsa distance implies large symmetry group]\label{zero-large}
  \lean{sub_mem_symmGroup}\leanok  If $X$ is a $G$-valued random variable such that
  $d[X ;X]=0$, and $x,y \in G$ are such that $P[X=x], P[X=y]>0$, then $x-y \in \mathrm{Sym}[X]$.
\end{lemma}

\begin{proof}
  \uses{ruz-indep, ruz-copy, relabeled-entropy-cond,vanish-entropy, alternative-mutual, independent-exist,sym-group-def}\leanok
   Let $X_1,X_2$ be independent copies of $X$ (from \Cref{independent-exist}). Let $A$ denote the range of $X$.  From \Cref{ruz-indep} and \Cref{ruz-copy} we have
  $$ \bbH[X_1-X_2] = \bbH[X_1].$$
Observe from \Cref{relabeled-entropy-cond} that
$$ \bbH[X_1-X_2|X_2] = \bbH[X_1|X_2] = \bbH[X_1]$$
and hence by \Cref{alternative-mutual}
$$ \bbI[X_1-X_2 : X_1] = 0.$$
By \Cref{vanish-entropy}, $X_1-X_2$ and $X_1$ are therefore independent, thus the law of $(X_1-X_2|X_1=x)$ does not depend on $x \in A$. The claim follows.
\end{proof}

\begin{lemma}[Translate is uniform on symmetry group]\label{sym-zero}
\lean{isUniform_sub_const_of_rdist_eq_zero}\leanok
  If $X$ is a $G$-valued random variable with $d[X ;X]=0$, and $x_0$ is a point with $P[X=x_0] > 0$, then $X-x_0$ is uniformly distributed on $\mathrm{Sym}[X]$.
\end{lemma}

\begin{proof}\uses{zero-large,sym-group-def,uniform-def}\leanok  The law of $X-x_0$ is invariant under $\mathrm{Sym}[X]$, non-zero at the origin, and supported on  $\mathrm{Sym}[X]$, giving the claim.
\end{proof}


\begin{lemma}[Symmetric 100\% inverse theorem]\label{lem:100pc-self}
  \lean{exists_isUniform_of_rdist_self_eq_zero}\leanok
  Suppose that $X$ is a $G$-valued random variable such that
  $d[X ;X]=0$. Then there exists a subgroup $H \leq G$ such that $d[X ;U_H] = 0$.
\end{lemma}

\begin{proof}\uses{sym-group, sym-zero}\leanok
Take $H$ to be the symmetry group of $X$, which is a group by \Cref{sym-group}. From \Cref{sym-zero}, $X-x_0$ is uniform on $H$, and $d[X ;X-x_0] = d[X ;X] \leq 0$, and the claim follows.
\end{proof}

\begin{corollary}[General 100\% inverse theorem]\label{lem:100pc}
  \lean{exists_isUniform_of_rdist_eq_zero}\leanok
  Suppose that $X_1,X_2$ are $G$-valued random variables such that
  $d[X_1;X_2]=0$. Then there exists a subgroup $H \leq G$ such that $d[X_1;U_H] = d[X_2;U_H] = 0$.
\end{corollary}

\begin{proof}\uses{lem:100pc-self,ruzsa-triangle, ruzsa-nonneg}\leanok
Using \Cref{ruzsa-triangle} and \Cref{ruzsa-nonneg} we have $d[X_1;X_1]=0$, hence by \Cref{lem:100pc-self} $d[X_1;U_H]=0$ for some subgroup $H$.  By \Cref{ruzsa-triangle} and \Cref{ruzsa-nonneg} again we also have $d[X_2;U_H]$ as required.
\end{proof}


%%FILE blueprint/src/chapter/approx_hom_pfr.tex
\chapter{Approximate homomorphism version of PFR}

\begin{definition}[Additive energy]
  \label{energy-def}
  \lean{Finset.addEnergy'}
  \leanok
  If $G$ is a finite additive group and $A$ is a subset of $G$, the \emph{additive energy}
  $E(A)$ of $A$ is the normalized count
  $$ E(A) := \frac{|\{(a_1,a_2,a_3,a_4) \in A^4 : a_1+a_2 = a_3+a_4\}|}{|G|^3}. $$
  The normalization by $|G|^3$ is the convention of \texttt{Finset.addEnergy'}, and every
  statement in this chapter that mentions $E(A)$ uses it.
\end{definition}

\begin{lemma}[Cauchy--Schwarz bound]
  \label{cs-bound}
  \uses{energy-def}
  \lean{Finset.card_sq_le_card_mul_addEnergy'}
  \leanok

  If $G$ is a finite additive group and $A,B$ are subsets of $G$, then
  $$ \left( \frac{|\{ (a,a') \in A \times A: a+a' \in B \}|}{|G|^2} \right)^2
     \leq \frac{|B|}{|G|}\, E(A). $$
  The left-hand side is the square of the density of a subset of $G \times G$, and the
  right-hand side involves the density of a subset of $G$.
\end{lemma}
\begin{proof}
  \leanok

  If $B$ is empty then the claim is trivial (with the Lean convention $0/0$),
  so without loss of generality $B$ is non-empty.  We can rewrite
  $$ |\{ (a,a') \in A \times A: a+a' \in B \}| = \sum_{b \in B} r(b)$$
  where $r: G \to \N$ is the counting function
  $$ r(b) := |\{ (a,a') \in A \times A: a+a' = b \}|.$$
  From double counting we have
  $$ \sum_{b \in G} r(b)^2 = |G|^3 E(A).$$
  The Cauchy--Schwarz inequality then gives
  $$ \Bigl( \sum_{b \in B} r(b) \Bigr)^2 \leq |B| \sum_{b \in B} r(b)^2
     \leq |B| \sum_{b \in G} r(b)^2 = |B|\, |G|^3 E(A),$$
  and the claim follows on dividing both sides by $|G|^4$.
\end{proof}

\begin{lemma}[Balog--Szemer\'edi--Gowers lemma]
  \label{bsg}
  \lean{BSG_self}
  \uses{energy-def}
  \leanok

  Let $G$ be a finite abelian group, and let $A$ be a non-empty subset of $G$
  with $E(A) \geq \frac{1}{K} \left( \frac{|A|}{|G|} \right)^3$ for some $K \geq 1$.
  Then there is a subset $A'$ of $A$ with $|A'| \geq |A| / (C_1 K^{C_2})$ and
  $|A'-A'| \leq C_3 K^{C_4} |A|$, where
  $$C_1 = 2^4, C_2 = 1, C_3 = 2^{10}, C_4 = 5.$$
  The Lean statement gives the two conclusions as inequalities between densities; these are
  equivalent to the ones above, the common factor $|G|$ cancelling.
\end{lemma}
\begin{proof}
  \leanok

  See \url{https://terrytao.files.wordpress.com/2024/01/simplebsg.pdf}.
\end{proof}

\begin{theorem}[Approximate homomorphism form of PFR]\label{approx-hom-pfr}\lean{approx_hom_pfr}\leanok Let $G,G'$ be finite abelian $2$-groups.
  Let $f: G \to G'$ be a function, and suppose that there are at least $|G|^2 / K$ pairs $(x,y) \in G^2$ such that
$$ f(x+y) = f(x) + f(y).$$
Then there exists a homomorphism $\phi: G \to G'$ and a constant $c \in G'$
such that $f(x) = \phi(x)+c$ for at least $|G| / (2 ^ {144} * K ^ {122})$
values of $x \in G$.
\end{theorem}

\begin{proof}\uses{goursat, cs-bound, bsg, pfr-9-aux'}\leanok Consider the graph $A \subset G \times G'$ defined by
$$ A := \{ (x,f(x)): x \in G \}.$$
Clearly, $|A| = |G|$.  By hypothesis, we have $a+a' \in A$ for at least
$|A|^2/K$ pairs $(a,a') \in A^2$. By \Cref{cs-bound}, this implies that
$E(A) \geq \frac{1}{K^2} \left( \frac{|A|}{|G \times G'|} \right)^3$.  Applying \Cref{bsg}, we conclude that there exists a subset
$A' \subset A$ with $|A'| \geq |A|/C_1 K^{2C_2}$ and $|A'+A'| \leq C_1C_3
K^{2(C_2+C_4)} |A'|$. Applying \Cref{pfr-9-aux'}, we may find a subspace $H
\subset G \times G'$ such that $|H| / |A'| \in [L^{-8}, L^{8}]$ and a subset
$c$ of cardinality at most $L^5 |A'|^{1/2} / |H|^{1/2}$ such that $A'
\subseteq c + H$, where $L =  C_1C_3 K^{2(C_2+C_4)}$. If we let $H_0,H_1$ be
as in \Cref{goursat}, this implies on taking projections the projection of
$A'$ to $G$ is covered by at most $|c|$ translates of $H_0$.  This implies
that
$$ |c| |H_0| \geq |A'|;$$
since $|H_0| |H_1| = |H|$, we conclude that
$$ |H_1| \leq |c| |H|/|A'|.$$

By hypothesis, $A'$ is covered by at most $|c|$ translates of $H$, and hence
by at most $|c| |H_1|$ translates of $\{ (x,\phi(x)): x \in G \}$.  As $\phi$
is a homomorphism, each such translate can be written in the form $\{
(x,\phi(x)+c): x \in G \}$ for some $c \in G'$. The number of translates is
bounded by
$$
|c|^2 \frac{|H|}{|A'|} \leq \left(L^5 \frac{|A'|^{1/2}}{|H|^{1/2}}\right)^2 \frac{|H|}{|A'|} = L^{10}.
$$

By the pigeonhole principle, one of these translates must then contain at
least $|A'|/L^{10} \geq |G| / (C_1C_3 K^{2(C_2+C_4)})^{10} (C_1 K^{2C_2})$
elements of $A'$ (and hence of $A$), and the claim follows.
\end{proof}

With a bit more effort, we can remove the constant term $c$, at the cost of reducing the set of agreement slightly.  We need some preliminary lemmas.

\begin{lemma}[Duality]\label{gdual}\lean{card_of_dual}\leanok Let $G$ be a finite abelian $2$-group.  Then the finite abelian $2$-group $\mathrm{Hom}(G,\Z/2\Z)$ of homomorphisms from $G$ to $\Z/2\Z$ has the same order as $G$.
\end{lemma}

\begin{proof}\leanok By the classification of finite abelian groups, $G$ is isomorphic to $(\Z/2\Z)^n$.  Then $\mathrm{Hom}(G,\Z/2\Z)$ is isomorphic to $(\Z/2\Z)^n$ as well, and hence has the same order.
\end{proof}

\begin{lemma}[Counting]\label{gcount}\lean{card_of_dual_constrained}\leanok Let $G$ be a finite abelian $2$-group, and let $x \in G$ be non-zero.  Then there are $|G|/2$ homomorphisms $\phi: G \to \Z/2\Z$ such that $\phi(x) = 1$.
\end{lemma}

\begin{proof}\uses{gdual} \leanok The map $\phi \mapsto \phi(x)$ is a homomorphism from $\mathrm{Hom}(G,\Z/2\Z)$ to $\Z/2\Z$, and by Lemma \ref{gdual} the kernel has order equal to the order of $G / \{0,x\}$, which is $|G|/2$.  Then the preimage of $1$ must also be of order $|G|/2$.
\end{proof}

\begin{lemma}[Slicing]\label{gslice}\lean{card_of_slice}\leanok  Let $G$ be a finite abelian $2$-group, and let $A$ be a subset of $G$.  Then there exists a homomorphism $\phi: G \to \Z/2\Z$ such that $|A \cap \phi^{-1}(1)| \geq (|A|-1)/2$.
\end{lemma}

\begin{proof}\uses{gdual, gcount}\leanok  We have
  $$ \sum_{\phi \in \mathrm{Hom}(G,\Z/2\Z)} |A \cap \phi^{-1}(1)| = \sum_{x \in A} |\{ \phi \in \mathrm{Hom}(G,\Z/2\Z): \phi(x) = 1 \}|$$
  $$ \geq (|A|-1) |G|/2$$
  thanks to Lemma \ref{gcount}. The claim now follows from Lemma \ref{gdual} and the pigeonhole principle.
\end{proof}

\begin{corollary}[Approximate homomorphism form of PFR, no constant term]\label{approx-hom-pfr-no-const}\lean{approx_hom_pfr'}\leanok Let $G,G'$ be finite abelian $2$-groups.
  Let $f: G \to G'$ be a function, and suppose that there are at least $|G|^2 / K$ pairs $(x,y) \in G^2$ such that
$$ f(x+y) = f(x) + f(y).$$
Then there exists a homomorphism $\phi'': G \to G'$
such that $f(x) = \phi''(x)$ for at least $(|G| / (2 ^ {144} * K ^ {122}) - 1)/2$ values of $x \in G$.
\end{corollary}

\begin{proof}\uses{approx-hom-pfr, gslice}\leanok By Theorem \ref{approx-hom-pfr}, there exists a homomorphism $\phi: G \to G'$ and a constant $c \in G'$ such that the set $A := \{ x \in G: f(x) = \phi(x) + c\}$ has cardinality at least $|G| / (2 ^ {144} * K ^ {122})$.  By Lemma \ref{gslice}, there exists a homomorphism $\phi': G \to \Z/2\Z$ such that
$$ |A \cap \phi'^{-1}(1)| \geq (|A|-1)/2 \geq (|G| / (2 ^ {144} * K ^ {122}) - 1)/2.$$
Then the claim follows by taking $\phi'' = \phi + \phi' \bullet c$ (where we view $G'$ as a $\Z/2\Z$-module).
\end{proof}


%%FILE blueprint/src/chapter/distance.tex
\chapter{Entropic Ruzsa calculus}

In this section $G$ will be a finite additive group.  (May eventually want to generalize to infinite $G$.)

\begin{lemma}[Negation preserves entropy]\label{neg-ent}
  \lean{ProbabilityTheory.entropy_neg}\leanok
  If $X$ is $G$-valued, then $\bbH[-X]=\bbH[X]$.
\end{lemma}

\begin{proof}
  \uses{relabeled-entropy}\leanok Immediate from \Cref{relabeled-entropy}.
\end{proof}

\begin{lemma}[Shearing preserves entropy]\label{shear-ent}
  \lean{ProbabilityTheory.condEntropy_add_right, ProbabilityTheory.condEntropy_add_left, ProbabilityTheory.condEntropy_sub_right, ProbabilityTheory.condEntropy_sub_left}\leanok
  If $X,Y$ are $G$-valued, then $\bbH[X \pm Y | Y]=\bbH[X|Y]$ and $\bbH[X \pm Y, Y] = \bbH[X, Y]$.
\end{lemma}

\begin{proof}
  \uses{relabeled-entropy-cond,chain-rule}\leanok Immediate from \Cref{relabeled-entropy-cond} and \Cref{chain-rule}.
\end{proof}

\begin{lemma}[Lower bound of sumset]\label{sumset-lower-gen}
  \lean{ProbabilityTheory.max_entropy_sub_mutualInfo_le_entropy_add, ProbabilityTheory.max_entropy_sub_mutualInfo_le_entropy_sub}\leanok
  If $X,Y$ are $G$-valued random variables on $\Omega$, we have
$$ \max(\bbH[X], \bbH[Y]) -  \bbI[X:Y] \leq \bbH[X \pm Y].$$
\end{lemma}

\begin{proof}\uses{cond-reduce, alternative-mutual, shear-ent, neg-ent} \leanok
    By \Cref{cond-reduce}, \ref{shear-ent}, \ref{alternative-mutual}, \ref{neg-ent} we have
$$
 \bbH[X\pm Y] \geq \bbH[X\pm Y|Y] = \bbH[X|Y]= \bbH[X] -  \bbI[X:Y]
$$
and similarly with the roles of $X,Y$ reversed, giving the claim.
\end{proof}

\begin{corollary}[Conditional lower bound on sumset]\label{sumset-lower-gen-cond}
  \uses{conditional-mutual-def}
  \lean{ProbabilityTheory.max_condEntropy_sub_condMutualInfo_le_condEntropy_mul, ProbabilityTheory.max_condEntropy_sub_condMutualInfo_le_condEntropy_div}\leanok
  If $X,Y$ are $G$-valued random variables on $\Omega$ and $Z$ is another random variable on $\Omega$ then
\[
  \max(\bbH[X|Z], \bbH[Y|Z]) -  \bbI[X:Y|Z] \leq \bbH[X\pm Y|Z],
\]
\end{corollary}

\begin{proof} \uses{sumset-lower-gen}\leanok This follows from \Cref{sumset-lower-gen} by conditioning to $Z = z$ and summing over $z$ (weighted by $\bbP[Z=z]$).
\end{proof}

\begin{corollary}[Independent lower bound on sumset]
  \label{sumset-lower}
  \lean{ProbabilityTheory.max_entropy_le_entropy_add, ProbabilityTheory.max_entropy_le_entropy_sub}
  \leanok

  If $X,Y$ are independent $G$-valued random variables, then
  $$\max(\bbH[X], \bbH[Y]) \leq \bbH[X\pm Y]. $$
\end{corollary}
\begin{proof}
  \uses{sumset-lower-gen, vanish-entropy}
  \leanok

  Combine \Cref{sumset-lower-gen} with \Cref{vanish-entropy}.
\end{proof}

One random variable is said to be a copy of another if they have the same distribution.

\begin{lemma}[Copy preserves entropy]\label{copy-ent}
  \uses{entropy-def}
  \lean{ProbabilityTheory.IdentDistrib.entropy_congr}\leanok
  If $X'$ is a copy of $X$ then $\bbH[X'] = \bbH[X]$.
\end{lemma}

\begin{proof}\leanok Immediate from \Cref{entropy-def}.
\end{proof}

\begin{lemma}[Existence of independent copies]\label{independent-exist}
  \uses{independent-def}
  \lean{ProbabilityTheory.independent_copies, ProbabilityTheory.independent_copies', ProbabilityTheory.independent_copies_two, ProbabilityTheory.independent_copies3_nondep} \leanok
  Let $X_i : \Omega_i \to S_i$ be random variables for $i=1,\dots,k$.  Then if one gives $\prod_{i=1}^k S_i$ the product measure of the laws of $X_i$, the coordinate functions $(x_j)_{j=1}^k \mapsto x_i$ are jointly independent random variables which are copies of the $X_1,\dots,X_k$.
\end{lemma}

\begin{proof} \leanok
  Explicit computation.
\end{proof}

\begin{definition}[Ruzsa distance]\label{ruz-dist-def}
  \uses{entropy-def, independent-exist}
  \lean{rdist_def}\leanok
  Let $X,Y$ be $G$-valued random variables (not necessarily on the same sample space).  The \emph{Ruzsa distance} $d[X ;Y]$ between $X$ and $Y$ is defined to be
$$ d[X ;Y] := \bbH[X' - Y'] - \bbH[X']/2 - \bbH[Y']/2$$
where $X',Y'$ are (the canonical) independent copies of $X,Y$ from \Cref{independent-exist}.
\end{definition}


\begin{lemma}[Distance from zero]\label{dist-zero}
  \uses{ruz-dist-def}\lean{rdist_zero_eq_half_ent}\leanok
If $X$ is a $G$-valued random variable and $0$ is the random variable taking the value $0$ everywhere then
\[d[X;0]=\mathbb{H}(X)/2.\]
\end{lemma}

\begin{proof}\leanok
This is an immediate consequence of the definitions and $X-0\equiv X$ and $\mathbb{H}(0)=0$.
\end{proof}

\begin{lemma}[Copy preserves Ruzsa distance]\label{ruz-copy}
  \uses{ruz-dist-def}
  \lean{ProbabilityTheory.IdentDistrib.rdist_congr}\leanok
  If $X',Y'$ are copies of $X,Y$ respectively then $d[X';Y']=d[X ;Y]$.
\end{lemma}

\begin{proof} \uses{copy-ent}\leanok Immediate from Definitions \ref{ruz-dist-def} and \Cref{copy-ent}.
\end{proof}

\begin{lemma}[Ruzsa distance in independent case]\label{ruz-indep}
  \uses{ruz-dist-def}
  \lean{ProbabilityTheory.IndepFun.rdist_eq}\leanok
  If $X,Y$ are independent $G$-random variables then
  $$ d[X ;Y] := \bbH[X - Y] - \bbH[X]/2 - \bbH[Y]/2.$$
\end{lemma}

\begin{proof} \uses{relabeled-entropy, copy-ent}\leanok Immediate from \Cref{ruz-dist-def} and Lemmas \ref{relabeled-entropy}, \ref{copy-ent}.
\end{proof}

\begin{lemma}[Distance symmetric]\label{ruzsa-symm}
  \uses{ruz-dist-def}
  \lean{rdist_symm}\leanok
  If $X,Y$ are $G$-valued random variables, then
  $$ d[X ;Y] = d[Y;X].$$
\end{lemma}

\begin{proof} \uses{neg-ent}\leanok Immediate from \Cref{neg-ent} and \Cref{ruz-dist-def}.
\end{proof}

\begin{lemma}[Distance controls entropy difference]\label{ruzsa-diff}
  \uses{ruz-dist-def}
  \lean{diff_ent_le_rdist}\leanok
  If $X,Y$ are $G$-valued random variables, then
$$|\bbH[X]-H[Y]| \leq 2 d[X ;Y].$$
\end{lemma}

\begin{proof} \uses{sumset-lower, neg-ent} \leanok Immediate from \Cref{sumset-lower} and \Cref{ruz-dist-def}, and also \Cref{neg-ent}.
\end{proof}

\begin{lemma}[Distance controls entropy growth]\label{ruzsa-growth}
  \uses{ruz-dist-def}
  \lean{diff_ent_le_rdist', diff_ent_le_rdist''}\leanok
  If $X,Y$ are independent $G$-valued random variables, then
$$  \bbH[X-Y] - \bbH[X], \bbH[X-Y] - \bbH[Y] \leq 2d[X ;Y].$$
\end{lemma}

\begin{proof} \uses{sumset-lower, neg-ent} \leanok Immediate from \Cref{sumset-lower} and \Cref{ruz-dist-def}, and also \Cref{neg-ent}.
\end{proof}

\begin{lemma}[Distance nonnegative]\label{ruzsa-nonneg}
  \lean{rdist_nonneg}\leanok
  If $X,Y$ are $G$-valued random variables, then
  $$ d[X ;Y] \geq 0.$$
\end{lemma}

\begin{proof}
  \uses{ruzsa-diff}
  \leanok
  Immediate from \Cref{ruzsa-diff}.
\end{proof}


\begin{lemma}[Projection entropy and distance]\label{dist-projection}\lean{ent_of_proj_le}\leanok
If $G$ is an additive group and $X$ is a $G$-valued random variable and $H\leq G$ is a finite subgroup then, with $\pi:G\to G/H$ the natural homomorphism we have (where $U_H$ is uniform on $H$)
\[\mathbb{H}(\pi(X))\leq 2d[X;U_H].\]
\end{lemma}
\begin{proof}
\uses{independent-exist, ruzsa-diff, chain-rule, shear-ent, submodularity, jensen-bound}\leanok
WLOG, we make $X$, $U_H$ independent (\Cref{independent-exist}).
Now by Lemmas \ref{submodularity}, \ref{shear-ent}, \ref{jensen-bound}
\begin{align*}
&\mathbb{H}(X-U_H|\pi(X)) \geq \mathbb{H}(X-U_H|X) &= \mathbb{H}(U_H) \\
&\mathbb{H}(X-U_H|\pi(X)) \leq \log |H| &= \mathbb{H}(U_H)
\end{align*}
By \Cref{chain-rule}
\[\mathbb{H}(X-U_H)=\mathbb{H}(\pi(X))+\mathbb{H}(X-U_H|\pi(X))=\mathbb{H}(\pi(X))+\mathbb{H}(U_H)\]
and therefore
\[d[X;U_H]=\mathbb{H}(\pi(X))+\frac{\mathbb{H}(U_H)-\mathbb{H}(X)}{2}.\]
Furthermore by \Cref{ruzsa-diff}
\[d[X;U_H]\geq \frac{\lvert \mathbb{H}(X)-\mathbb{H}(U_H)\rvert}{2}.\]
Adding these inequalities gives the result.
\end{proof}

\begin{lemma}[Improved Ruzsa triangle inequality]
  \label{ruzsa-triangle-improved}
  \lean{ent_of_diff_le}
  \leanok
  If $X,Y,Z$ are $G$-valued random variables on $\Omega$ with $(X,Y)$ independent of $Z$, then
  \begin{equation}\label{submod-explicit} \bbH[X - Y] \leq \bbH[X-Z] + \bbH[Z-Y] - \bbH[Z]\end{equation}
\end{lemma}

This is an improvement over the usual Ruzsa triangle inequality because $X,Y$ are not assumed to be independent.  However we will not utilize this improvement here.

\begin{proof}
  \uses{alt-submodularity,subadditive, relabeled-entropy, add-entropy}
  \leanok
  Apply \Cref{alt-submodularity} to obtain
  \[\bbH[X - Z, X - Y] + \bbH[Y, X - Y] \geq \bbH[X - Z, Y, X - Y] + \bbH[X - Y].\]
  Using
  \[\bbH[X - Z, X - Y] \leq \bbH[X - Z] + \bbH[Y - Z]\]
  (from \Cref{relabeled-entropy}, \Cref{subadditive}),
  \[\bbH[Y, X - Y] = \bbH[X, Y] \]
  (from \Cref{relabeled-entropy}), and
  \[\bbH[X - Z, Y, X - Y] = \bbH[X, Y, Z] = \bbH[X, Y] + \bbH[Z]\]
  (from \Cref{relabeled-entropy} and \Cref{add-entropy}) and rearranging, we indeed obtain~\eqref{submod-explicit}.
\end{proof}

\begin{lemma}[Ruzsa triangle inequality]
  \label{ruzsa-triangle}
  \lean{rdist_triangle}
  \leanok
  If $X,Y,Z$ are $G$-valued random variables, then
$$ d[X ;Y] \leq d[X ;Z] + d[Z;Y].$$
\end{lemma}

\begin{proof}\uses{ruz-copy, ruz-indep, independent-exist, ruzsa-triangle-improved}\leanok By \Cref{ruz-copy} and Lemmas \ref{independent-exist}, \ref{ruz-indep}, it suffices to prove this inequality assuming that $X,Y,Z$ are defined on the same space and are independent.  But then the claim follows from \Cref{ruzsa-triangle-improved} and \Cref{ruz-dist-def}.
\end{proof}

\begin{definition}[Conditioned Ruzsa distance]\label{cond-dist-def}
  \uses{ruz-dist-def}
  \lean{condRuzsaDist}\leanok
If $(X, Z)$ and $(Y, W)$ are random variables (where $X$ and $Y$ are $G$-valued) we define
$$ d[X  | Z; Y | W] := \sum_{z,w}  \bbP[Z=z]  \bbP[W=w] d[(X|Z=z); (Y|(W=w))].$$
similarly
$$ d[X ; Y | W] := \sum_{w}  \bbP[W=w] d[X ; (Y|(W=w))].$$
\end{definition}

\begin{lemma}[Alternate form of distance]\label{cond-dist-alt}
\lean{condRuzsaDist_of_copy, condRuzsaDist'_of_copy, condRuzsaDist_of_indep, condRuzsaDist'_of_indep}\leanok
  The expression $d[X  | Z;Y | W]$ is unchanged if $(X,Z)$ or $(Y,W)$ is replaced by a copy.  Furthermore, if $(X,Z)$ and $(Y,W)$ are independent, then
$$  d[X  | Z;Y | W] = \bbH[X-Y|Z,W] - \bbH[X|Z]/2 - \bbH[Y|W]/2$$
and similarly
$$  d[X ;Y | W] = \bbH[X-Y|W] - \bbH[X]/2 - \bbH[Y|W]/2.$$
\end{lemma}

\begin{proof}\uses{copy-ent, ruz-copy, ruz-indep, cond-dist-def, conditional-entropy-def}\leanok Straightforward thanks to \Cref{copy-ent}, \Cref{ruz-copy}, \Cref{ruz-indep}, \Cref{cond-dist-def}, \Cref{conditional-entropy-def}.
\end{proof}

\begin{lemma}[Kaimanovich-Vershik-Madiman inequality]\label{kv}
  \lean{kaimanovich_vershik}\leanok
Suppose that $X, Y, Z$ are independent $G$-valued random variables. Then
\[
  \bbH[X + Y + Z] - \bbH[X + Y] \leq \bbH[Y+ Z] - \bbH[Y].
\]
\end{lemma}

\begin{proof}\uses{submodularity, add-entropy, relabeled-entropy}\leanok
From \Cref{submodularity} we have
$$ \bbH[X, X + Y+ Z] + \bbH[Z, X + Y+ Z] \geq \bbH[X, Z, X + Y+ Z] + \bbH[X + Y+ Z].$$
However, using Lemmas \ref{add-entropy}, \ref{relabeled-entropy} repeatedly we have $\bbH[X, X + Y+ Z] = \bbH[X, Y+ Z] = \bbH[X] + \bbH[Y+ Z]$, $\bbH[Z, X + Y + Z] = \bbH[Z, X + Y] = \bbH[Z] + \bbH[X + Y]$ and $\bbH[X, Z, X + Y+ Z] = \bbH[X, Y, Z] = \bbH[X] + \bbH[Y] + \bbH[Z]$.  The claim then follows from a calculation.
\end{proof}

\begin{lemma}[Existence of conditional independent trials]\label{cond-indep-exist}
  \uses{conditional-independent-def}
  \lean{ProbabilityTheory.condIndep_copies}\leanok
  For $X,Y$ random variables, there exist random variables $X_1,X_2,Y'$ on a common probability space with $(X_1, Y'), (X_2, Y')$ both having the distribution of $(X,Y)$, and $X_1, X_2$ conditionally independent over $Y'$ in the sense of \Cref{conditional-independent-def}.
\end{lemma}

\begin{proof}\leanok Explicit construction.
\end{proof}


\begin{lemma}[Balog-Szemer\'edi-Gowers]\label{entropic-bsg}
  \lean{ent_bsg}\leanok
  Let $A,B$ be $G$-valued random variables on $\Omega$, and set $Z := A+B$.
Then
\begin{equation}\label{2-bsg-takeaway} \sum_{z}  \bbP[Z=z] d[(A | Z = z); (B | Z = z)] \leq 3  \bbI[A:B] + 2 \bbH[Z] - \bbH[A] - \bbH[B]. \end{equation}
\end{lemma}

\begin{proof}
  \uses{cond-indep-exist, cond-trial-ent, conditional-entropy-def,submodularity, copy-ent, relabeled-entropy, add-entropy, ruz-indep}
  \leanok
Let $(A_1, B_1)$ and $(A_2, B_2)$ (and $Z'$, which by abuse of notation we call $Z$) be conditionally independent trials of $(A,B)$ relative to $Z$ as produced by \Cref{cond-indep-exist}, thus $(A_1,B_1)$ and $(A_2,B_2)$ are coupled through the random variable $A_1 + B_1 = A_2 + B_2$, which by abuse of notation we shall also call $Z$.

Observe from \Cref{ruz-indep} that the left-hand side of~\eqref{2-bsg-takeaway} is
\begin{equation}\label{lhs-to-bound}
\bbH[A_1 - B_2| Z] - \bbH[A_1 | Z]/2 - \bbH[B_2 | Z]/2.
\end{equation}
since, crucially, $(A_1 | Z=z)$ and $(B_2 | Z=z)$ are independent for all $z$.

Applying submodularity (\Cref{alt-submodularity}) gives
\begin{equation}\label{bsg-31} \begin{split}
&\bbH[A_1 - B_2] + \bbH[A_1 - B_2, A_1, B_1] \\
&\qquad \leq \bbH[A_1 - B_2, A_1] + \bbH[A_1 - B_2,B_1].
\end{split}\end{equation}
We estimate the second, third and fourth terms appearing here.
First note that, by \Cref{cond-trial-ent} and \Cref{relabeled-entropy} (noting that the tuple $(A_1 - B_2, A_1, B_1)$  determines the tuple $(A_1, A_2, B_1, B_2)$ since $A_1+B_1=A_2+B_2$)
\begin{equation}\label{bsg-24} \bbH[A_1 - B_2, A_1, B_1] = \bbH[A_1, B_1, A_2, B_2,Z] = 2\bbH[A,B] - \bbH[Z].\end{equation}
Next observe that
\begin{equation}\label{bsg-23} \bbH[A_1 - B_2, A_1] = \bbH[A_1, B_2] \leq \bbH[A] + \bbH[B].
\end{equation}
Finally, we have
\begin{equation}\label{bsg-25} \bbH[A_1 - B_2, B_1] = \bbH[A_2 - B_1, B_1] = \bbH[A_2, B_1] \leq \bbH[A] + \bbH[B].\end{equation}
Substituting~\eqref{bsg-24},~\eqref{bsg-23} and~\eqref{bsg-25} into~\eqref{bsg-31} yields
\[\bbH[A_1 - B_2] \leq 2 \bbI[A:B] + \bbH[Z]\] and so by \Cref{cond-reduce}
\[\bbH[A_1 - B_2 | Z]  \leq 2 \bbI[A:B] + \bbH[Z].\]
Since
\begin{align*} \bbH[A_1 | Z] & = \bbH[A_1, A_1 + B_1] - \bbH[Z] \\ & = \bbH[A,B] - \bbH[Z] \\ & = \bbH[Z] -  \bbI[A:B] - 2 \bbH[Z]- \bbH[A]-\bbH[B]\end{align*}
and similarly for $\bbH[B_2 | Z]$, we see that~\eqref{lhs-to-bound} is bounded by
$3 \bbI[A:B] + 2\bbH[Z]-\bbH[A]-\bbH[B]$ as claimed.
\end{proof}


\begin{lemma}[Upper bound on conditioned Ruzsa distance]\label{cond-dist-fact}
  \uses{cond-dist-def, information-def}
  \lean{condRuzsaDist_le, condRuzsaDist_le'}\leanok
  Suppose that $(X, Z)$ and $(Y, W)$ are random variables, where $X, Y$ take values in an abelian group. Then
  \[   d[X  | Z;Y | W] \leq d[X ; Y] + \tfrac{1}{2}  \bbI[X : Z] + \tfrac{1}{2}  \bbI[Y : W].\]
  In particular,
  \[   d[X ;Y | W] \leq d[X ; Y] + \tfrac{1}{2}  \bbI[Y : W].\]
\end{lemma}

\begin{proof}
\uses{cond-dist-alt, independent-exist, cond-reduce}\leanok
Using \Cref{cond-dist-alt} and \Cref{independent-exist}, if $(X',Z'), (Y',W')$ are independent copies of the variables $(X,Z)$, $(Y,W)$, we have
\begin{align*}
  d[X  | Z; Y | W]&= \bbH[X'-Y'|Z',W'] - \tfrac{1}{2} \bbH[X'|Z'] - \tfrac{1}{2}H[Y'|W'] \\
                       &\le \bbH[X'-Y']- \tfrac{1}{2} \bbH[X'|Z'] - \tfrac{1}{2}H[Y'|W'] \\
                       &= d[X';Y'] + \tfrac{1}{2}  \bbI[X' : Z'] + \tfrac{1}{2}  \bbI[Y' : W'].
\end{align*}
Here, in the middle step we used \Cref{cond-reduce}, and in the last step we used \Cref{ruz-dist-def} and \Cref{information-def}.
\end{proof}

\begin{lemma}[Comparison of Ruzsa distances, I]\label{first-useful}
  \lean{condRuzsaDist_diff_le, condRuzsaDist_diff_le', condRuzsaDist_diff_le'', condRuzsaDist_diff_le'''}\leanok
  Let $X, Y, Z$ be random variables taking values in some abelian group of characteristic $2$, and with $Y, Z$ independent. Then we have
  \begin{align}\nonumber d[X ; Y + Z] -d[X ; Y] &  \leq \tfrac{1}{2} (\bbH[Y+ Z] - \bbH[Y]) \\ & = \tfrac{1}{2} d[Y; Z] + \tfrac{1}{4} \bbH[Z] - \tfrac{1}{4} \bbH[Y]. \label{lem51-a} \end{align}
  and
  \begin{align}\nonumber
  d[X ;Y|Y+ Z] - d[X ;Y] & \leq \tfrac{1}{2} \bigl(\bbH[Y+ Z] - \bbH[Z]\bigr) \\ & = \tfrac{1}{2} d[Y;Z] + \tfrac{1}{4} \bbH[Y] - \tfrac{1}{4} \bbH[Z].
    \label{ruzsa-3}
  \end{align}
  \end{lemma}

  \begin{proof}
    \uses{ruz-copy, independent-exist, kv, ruz-indep, relabeled-entropy, cond-dist-fact}\leanok
  We first prove~\eqref{lem51-a}. We may assume (taking an independent copy, using \Cref{independent-exist} and \Cref{ruz-copy}, \ref{ruz-indep}) that $X$ is independent of $Y, Z$. Then we have
  \begin{align*}  d[X ;Y+ Z] & - d[X ;Y] \\ & = \bbH[X + Y + Z] - \bbH[X + Y] - \tfrac{1}{2}\bbH[Y + Z] + \tfrac{1}{2} \bbH[Y].\end{align*}
  Combining this with \Cref{kv} gives the required bound. The second form of the result is immediate \Cref{ruz-indep}.

  Turning to~\eqref{ruzsa-3}, we have from \Cref{information-def} and \Cref{relabeled-entropy}
  \begin{align*}  \bbI[Y : Y+ Z] & = \bbH[Y] + \bbH[Y + Z] - \bbH[Y, Y + Z] \\ & = \bbH[Y] + \bbH[Y + Z] - \bbH[Y, Z]  = \bbH[Y + Z] - \bbH[Z],\end{align*}
  and so~\eqref{ruzsa-3} is a consequence of \Cref{cond-dist-fact}. Once again the second form of the result is immediate from \Cref{ruz-indep}.
\end{proof}

\begin{lemma}[Comparison of Ruzsa distances, II]\label{second-useful}
  \lean{condRuzsaDist_diff_ofsum_le}\leanok
  Let $X, Y, Z, Z'$ be random variables taking values in some abelian group, and with $Y, Z, Z'$ independent. Then we have
  \begin{align}\nonumber
  & d[X ;Y + Z | Y + Z + Z'] - d[X ;Y] \\ & \qquad \leq \tfrac{1}{2} ( \bbH[Y + Z + Z'] + \bbH[Y + Z] - \bbH[Y] - \bbH[Z']).\label{7111}
  \end{align}
  \end{lemma}

  \begin{proof}
    \uses{first-useful}\leanok
  By \Cref{first-useful} (with a change of variables) we have
  \[d[X ; Y + Z | Y + Z + Z'] - d[X ; Y + Z] \leq \tfrac{1}{2}( \bbH[Y + Z + Z'] - \bbH[Z']).\]
  Adding this to~\eqref{lem51-a} gives the result.
  \end{proof}


%%FILE blueprint/src/chapter/entropy.tex
\chapter{Shannon entropy inequalities}

Random variables in this paper are measurable maps $X : \Omega \to S$ from a probability space $\Omega$ to a measurable space $S$, and called $S$-valued random variables. In many cases we will assume that singletons in $S$ are measurable.  Often we will restrict further to the case when $S$ is finite with the discrete $\sigma$-algebra, which of course implies that $S$ has measurable singletons.

\begin{definition}[Entropy]
  \label{entropy-def}
  \lean{ProbabilityTheory.entropy}
  \leanok
  If $X$ is an $S$-valued random variable, the entropy $\bbH[X]$ of $X$ is defined
  $$ \bbH[X] := \sum_{s \in S} \bbP[X=x] \log \frac{1}{\bbP[X=x]}$$
  with the convention that $0 \log \frac{1}{0} = 0$.
\end{definition}

\begin{lemma}[Entropy and relabeling]\label{relabeled-entropy}\lean{ProbabilityTheory.entropy_comp_of_injective, ProbabilityTheory.entropy_of_comp_eq_of_comp} \uses{entropy-def}\leanok
  \begin{itemize}
\item[(i)]   If $X: \Omega \to S$ and $Y: \Omega \to T$ are random variables, and $Y = f(X)$ for some injection $f: S \to T$, then $\bbH[X] = \bbH[Y]$.
\item[(ii)]   If $X: \Omega \to S$ and $Y: \Omega \to T$ are random variables, and $Y = f(X)$ and $X = g(Y)$ for some functions $f: S \to T$, $g: T \to S$, then $\bbH[X] = \bbH[Y]$.
\end{itemize}
\end{lemma}

\begin{proof} \leanok Expand out both entropies and rearrange.
\end{proof}

\begin{lemma}[Jensen bound]\label{jensen-bound}
  \uses{entropy-def}
  \lean{ProbabilityTheory.entropy_le_log_card, ProbabilityTheory.entropy_le_log_card_of_mem}
  \leanok
  If $X$ is an $S$-valued random variable, then $\bbH[X] \leq \log |S|$.
\end{lemma}

\begin{proof}
  \uses{concave}
  \leanok

  This is a direct consequence of \Cref{concave} and Jensen's inequality.
\end{proof}

\begin{definition}[Uniform distribution]\label{uniform-def}\leanok
  \lean{ProbabilityTheory.IsUniform} If $H$ is a subset of $S$, an $S$-random variable $X$ is said to be uniformly distributed on $H$ if $\bbP[X = s] = 1/|H|$ for $s \in X$ and $\bbP[X=s] = 0$ otherwise.
\end{definition}

\begin{lemma}[Uniform distributions exist]\label{unif-exist}
  \uses{uniform-def}
  \lean{ProbabilityTheory.exists_isUniform, ProbabilityTheory.exists_isUniform_measureSpace}\leanok
  Given a finite non-empty subset $H$ of a set $S$, there exists a random variable $X$ (on some probability space) that is uniformly distributed on $H$.
\end{lemma}

\begin{proof}\leanok Direct construction.
\end{proof}


\begin{lemma}[Entropy of uniform random variable]\label{uniform-entropy}
  \uses{entropy-def}
  \lean{ProbabilityTheory.IsUniform.entropy_eq}\leanok
  If $X$ is $S$-valued random variable, then $\bbH[X] = \log |S|$ if and only if $X$ is uniformly distributed on $S$.
\end{lemma}

\begin{proof}
  \uses{concave}
  \leanok

  Direct computation in one direction.
  Converse direction needs the strict Jensen inequality and \Cref{concave}.
\end{proof}

\begin{lemma}[Entropy of uniform random variable, II]\label{uniform-entropy-II}
  \uses{entropy-def, uniform-def}
  \lean{ProbabilityTheory.IsUniform.entropy_eq}\leanok
  If $X$ is uniformly distributed on $H$, then, then $\bbH[X] = \log |H|$.
\end{lemma}

\begin{proof}\leanok Direct computation.
\end{proof}

\begin{lemma}[Bounded entropy implies concentration]\label{bound-conc}
  \uses{entropy-def}
  \lean{ProbabilityTheory.prob_ge_exp_neg_entropy}\leanok
  If $X$ is an $S$-valued random variable, then there exists $s \in S$ such that $\bbP[X=s] \geq \exp(-\bbH[X])$.
\end{lemma}

\begin{proof}
  \leanok
  We have
  $$ \bbH[X] = \sum_{s \in S} \bbP[X=s] \log \frac{1}{\bbP[X=s]} \geq \min_{s \in S} \log \frac{1}{\bbP[X=s]}$$
  and the claim follows.
\end{proof}

We use $X,Y$ to denote the pair $\omega \mapsto (X(\omega),Y(\omega)$).

\begin{lemma}[Commutativity and associativity of joint entropy]
  \label{entropy-comm}
  \lean{ProbabilityTheory.entropy_comm, ProbabilityTheory.entropy_assoc}
  \leanok
  If $X: \Omega \to S$, $Y: \Omega \to T$, and $Z: \Omega \to U$ are random variables, then $\bbH[X, Y] = \bbH[Y, X]$ and $\bbH[X, (Y,Z)] = \bbH[(X,Y), Z]$.
\end{lemma}
\begin{proof}
  \leanok
  \uses{relabeled-entropy}
  Set up an injection from $(X,Y)$ to $(Y,X)$ and use \Cref{relabeled-entropy} for the first claim. Similarly for the second claim.
\end{proof}


\begin{definition}[Conditioned event]
  \label{condition-event-def}
  \lean{ProbabilityTheory.cond}
  \leanok
  If $X: \Omega \to S$ is an $S$-valued random variable and $E$ is an event in $\Omega$, then the conditioned event $(X|E)$ is defined to be the same random variable as $X$, but now the ambient probability measure has been conditioned to $E$.
\end{definition}

Note: it may happen that $E$ has zero measure.  In which case, the ambient probability measure should be replaced with a zero measure.  (In our formalization we achieve this by working with arbitrary measures, and normalizing them to be probability measures if possible, else using the zero measure.  Conditioning is also formalized using existing Mathlib definitions.)

\begin{definition}[Conditional entropy]
  \label{conditional-entropy-def}
  \uses{entropy-def,condition-event-def}
  \lean{ProbabilityTheory.condEntropy}
  \leanok
  If $X: \Omega \to S$ and $Y: \Omega \to T$ are random variables, the conditional entropy $\bbH[X|Y]$ is defined as
  $$ \bbH[X|Y] := \sum_{y \in Y} \bbP[Y = y] \bbH[(X | Y=y)].$$
\end{definition}

\begin{lemma}[Conditional entropy and relabeling]\label{relabeled-entropy-cond}
  \uses{conditional-entropy-def}
  \lean{ProbabilityTheory.condEntropy_comp_of_injective, ProbabilityTheory.condEntropy_of_injective'}\leanok
  If $X: \Omega \to S$, $Y: \Omega \to T$, and $Z: \Omega \to U$ are random variables, and $Y = f(X,Z)$ almost surely for some map $f: S \times U \to T$ that is injective for each fixed $U$, then $\bbH[X|Z] = \bbH[Y|Z]$.

  Similarly, if $g: T \to U$ is injective, then $\bbH[X|g(Y)] = \bbH[X|Y]$.
\end{lemma}

\begin{proof}
  \uses{relabeled-entropy}\leanok
  For the first part, use \Cref{conditional-entropy-def} and then \Cref{relabeled-entropy}.  The second part is a direct computation.
\end{proof}


\begin{lemma}[Chain rule]\label{chain-rule}
  \uses{conditional-entropy-def}
  \lean{ProbabilityTheory.chain_rule, ProbabilityTheory.chain_rule'}
  \leanok
  If $X: \Omega \to S$ and $Y: \Omega \to T$ are random variables, then
  $$ \bbH[X, Y] = \bbH[Y] + \bbH[X|Y].$$
  \end{lemma}
  \begin{proof}
  \leanok
  Direct computation.
\end{proof}

\begin{lemma}[Conditional chain rule]\label{conditional-chain-rule}
  \lean{ProbabilityTheory.cond_chain_rule, ProbabilityTheory.cond_chain_rule'}\leanok
  If $X: \Omega \to S$, $Y: \Omega \to T$, $Z: \Omega \to U$ are random variables, then
$$ \bbH[X, Y | Z] = \bbH[Y | Z] + \bbH[X|Y, Z].$$
\end{lemma}

\begin{proof}  \uses{chain-rule}\leanok
   For each $z \in U$, we can apply \Cref{chain-rule} to the random variables $(X|Z=z)$ and $(Y|Z=z)$ to obtain
  $$ \bbH[(X|Z=z),(Y|Z=z)] = \bbH[Y|Z=z] + \bbH[(X|Z=z)|(Y|Z=z)].$$
  Now multiply by $\bbP[Z=z]$ and sum.  Some helper lemmas may be needed to get to the form above.  This sort of ``average over conditioning'' argument to get conditional entropy inequalities from unconditional ones is commonly used in this paper.
\end{proof}

\begin{definition}[Mutual information]
  \label{information-def}
  \uses{entropy-def}
  \lean{ProbabilityTheory.mutualInfo, ProbabilityTheory.mutualInfo_def}
  \leanok
  If $X: \Omega \to S$, $Y: \Omega \to T$ are random variables, then
  $$\bbI[X : Y] := \bbH[X] + \bbH[Y] - \bbH[X,Y].$$
\end{definition}

\begin{lemma}[Alternative formulae for mutual information]
  \label{alternative-mutual}
  \uses{information-def}
  \lean{ProbabilityTheory.mutualInfo_comm, ProbabilityTheory.mutualInfo_eq_entropy_sub_condEntropy}
  \leanok
  With notation as above, we have
  $$  \bbI[X : Y] = \bbI[Y:X]$$
  $$  \bbI[X : Y] = \bbH[X] - \bbH[X|Y]$$
  $$  \bbI[X : Y] = \bbH[Y] - \bbH[Y|X]$$
\end{lemma}

\begin{proof}
  \uses{entropy-comm,chain-rule}
  \leanok
  Immediate from Lemmas \ref{entropy-comm}, \ref{chain-rule}.
\end{proof}

\begin{lemma}[Nonnegativity of mutual information]
  \label{mutual-nonneg}
  \lean{ProbabilityTheory.mutualInfo_nonneg}
  \leanok
  We have $\bbI[X:Y] \geq 0$.
\end{lemma}

\begin{proof}
  \uses{concave, alternative-mutual}
  \leanok

  An application of jensen's inequality and \Cref{concave,alternative-mutual}.
\end{proof}

\begin{corollary}[Subadditivity]
  \label{subadditive}
  \lean{ProbabilityTheory.entropy_pair_le_add}
  \leanok
  With notation as above, we have $\bbH[X,Y] \leq \bbH[X] + \bbH[Y]$.
\end{corollary}

\begin{proof}
  \uses{mutual-nonneg, alternative-mutual}
  \leanok
  Use \Cref{mutual-nonneg}.
\end{proof}

\begin{corollary}[Conditioning reduces entropy]
  \label{cond-reduce}
  \lean{ProbabilityTheory.condEntropy_le_entropy}
  \leanok
  With notation as above, we have $\bbH[X|Y] \leq \bbH[X]$.
\end{corollary}
\begin{proof}
  \uses{mutual-nonneg, alternative-mutual}
  \leanok
  Combine \Cref{mutual-nonneg} with \Cref{alternative-mutual}.
\end{proof}

\begin{corollary}[Submodularity]\label{submodularity}
  \lean{ProbabilityTheory.entropy_submodular, ProbabilityTheory.condEntropy_comp_ge}\leanok
  With three random variables $X,Y,Z$, one has $\bbH[X|Y,Z] \leq \bbH[X|Z]$.
\end{corollary}

\begin{proof} \uses{cond-reduce}\leanok Apply the ``averaging over conditioning'' argument to \Cref{cond-reduce}.
\end{proof}

\begin{corollary}[Alternate form of submodularity]\label{alt-submodularity}
  \lean{ProbabilityTheory.entropy_triple_add_entropy_le}\leanok
  With three random variables $X,Y,Z$, one has
  $$ \bbH[X,Y,Z] + \bbH[Z] \leq \bbH[X,Z] + \bbH[Y,Z].$$
\end{corollary}

\begin{proof}  \uses{submodularity,chain-rule} \leanok Apply \Cref{submodularity} and \Cref{chain-rule}.
\end{proof}

\begin{definition}[Independent random variables]
  \label{independent-def}
  \lean{ProbabilityTheory.IndepFun}
  \leanok
  Two random variables $X: \Omega \to S$ and $Y: \Omega \to T$ are independent if the law of $(X,Y)$ is the product of the law of $X$ and the law of $Y$.  Similarly for more than two variables.
\end{definition}

\begin{lemma}[Vanishing of mutual information]
  \label{vanish-entropy}
  \uses{information-def, independent-def}
  \lean{ProbabilityTheory.mutualInfo_eq_zero}
  \leanok
  If $X,Y$ are random variables, then $\bbI[X:Y] = 0$ if and only if $X,Y$ are independent.
  \end{lemma}

  \begin{proof}
    \uses{concave}
    \leanok

    An application of the equality case of Jensen's inequality and \Cref{concave}.
  \end{proof}

\begin{corollary}[Additivity of entropy]\label{add-entropy}
  \lean{ProbabilityTheory.entropy_pair_eq_add}
  \leanok
  If $X,Y$ are random variables, then $\bbH[X,Y] = \bbH[X] + \bbH[Y]$ if and only if $X,Y$ are independent.
\end{corollary}

\begin{proof} \uses{vanish-entropy}\leanok
  Direct from \Cref{vanish-entropy}.
\end{proof}

\begin{definition}[Conditional mutual information]
\label{conditional-mutual-def}
\uses{information-def,condition-event-def}
\lean{ProbabilityTheory.condMutualInfo}
\leanok
If $X,Y,Z$ are random variables, with $Z$ $U$-valued, then
$$ \bbI[X:Y|Z] := \sum_{z \in U} P[Z=z] \bbI[(X|Z=z): (Y|Z=z)].$$
\end{definition}

\begin{lemma}[Alternate formula for conditional mutual information]
  \label{conditional-mutual-alt}
  \uses{conditional-mutual-def}
  \lean{ProbabilityTheory.condMutualInfo_eq, ProbabilityTheory.condMutualInfo_eq'}\leanok
We have
  $$ \bbI[X:Y|Z] := \bbH[X|Z] + \bbH[Y|Z] - \bbH[X,Y|Z]$$
and
  $$ \bbI[X:Y|Z] := \bbH[X|Z] - \bbH[X|Y,Z].$$
\end{lemma}

\begin{proof}\leanok Routine computation.
\end{proof}

\begin{lemma}[Nonnegativity of conditional mutual information]
\label{conditional-nonneg}
\uses{conditional-mutual-def}
\lean{ProbabilityTheory.condMutualInfo_nonneg}\leanok
If $X,Y,Z$ are random variables, then $\bbI[X:Y|Z] \ge 0$.
\end{lemma}
\begin{proof}
\uses{submodularity}
\leanok
Use \Cref{conditional-mutual-def} and \Cref{submodularity}.
\end{proof}

\begin{definition}[Conditionally independent random variables]
\label{conditional-independent-def}
\lean{ProbabilityTheory.CondIndepFun}\leanok
  Two random variables $X: \Omega \to S$ and $Y: \Omega \to T$ are conditionally independent relative to another random variable $Z: \Omega \to U$ if $P[X = s \wedge Y = t| Z=u] = P[X=s|Z=u] P[Y=t|Z=u]$ for all $s \in S, t \in T, u \in U$.  (We won't need conditional independence for more variables than this.)
\end{definition}

\begin{lemma}[Vanishing conditional mutual information]\label{conditional-vanish}
  \lean{ProbabilityTheory.condMutualInfo_eq_zero}
  \leanok
  If $X,Y,Z$ are random variables, then $\bbI[X:Y|Z] = 0$ iff $X,Y$ are conditionally independent over $Z$.
\end{lemma}

\begin{proof} \uses{vanish-entropy, conditional-independent-def}\leanok Immediate from \Cref{vanish-entropy} and \Cref{conditional-independent-def}.
\end{proof}

\begin{corollary}[Entropy of conditionally independent variables]\label{cond-trial-ent}
  \lean{ProbabilityTheory.ent_of_cond_indep}\leanok
  If $X, Y$ are conditionally independent over $Z$, then
  $$ \bbH[X,Y,Z] =\bbH[X,Z] + \bbH[Y,Z] - \bbH[Z].$$
\end{corollary}

\begin{proof} \uses{conditional-vanish, conditional-mutual-alt}\leanok Immediate from \Cref{conditional-vanish} and \Cref{conditional-mutual-alt}.
\end{proof}


%%FILE blueprint/src/chapter/entropy_pfr.tex
\chapter{Entropy version of PFR}

\begin{definition}\label{eta-def}
 \leanok
  $\eta := 1/9$.
\end{definition}

Throughout this chapter,  $G = \F_2^n$, and $X^0_1, X^0_2$ are $G$-valued random variables.

\begin{definition}[$\tau$ functional]\label{tau-def}
  \uses{eta-def, ruz-dist-def}
  \lean{tau}\leanok
If $X_1,X_2$ are two $G$-valued random variables, then
$$  \tau[X_1; X_2] := d[X_1; X_2] + \eta  d[X^0_1; X_1] + \eta d[X^0_2; X_2].$$
\end{definition}

\begin{lemma}[$\tau$ depends only on distribution]\label{tau-copy}\leanok
  \uses{tau-def}
  \lean{ProbabilityTheory.IdentDistrib.tau_eq}  If $X'_1, X'_2$ are copies of $X_1,X_2$, then $\tau[X'_1;X'_2] = \tau[X_1;X_2]$.
\end{lemma}


\begin{proof}\uses{copy-ent}\leanok Immediate from \Cref{copy-ent}.
\end{proof}

\begin{definition}[$\tau$-minimizer]\label{tau-min-def}
\uses{tau-def}
\lean{TauMinimizes}\leanok
A pair of $G$-valued random variables $X_1, X_2$ are said to be a $\tau$-minimizer if one has
  $$\tau[X_1;X_2] \leq \tau[X'_1;X'_2]
  $$
for all $G$-valued random variables $X'_1, X'_2$.
\end{definition}


\begin{proposition}[$\tau$ has minimum]\label{tau-min}\uses{tau-min-def}
  \lean{tau_minimizer_exists}\leanok
A pair $X_1, X_2$ of $\tau$-minimizers exist.
\end{proposition}

\begin{proof}\uses{tau-copy}\leanok By \Cref{tau-copy}, $\tau$ only depends on the probability distributions of $X_1, X_2$. This ranges over a compact space, and $\tau$ is continuous.  So $\tau$ has a minimum.
\end{proof}

\section{Basic facts about minimizers}

In this section we assume that $X_1,X_2$ are $\tau$-minimizers. We also write $k := d[X_1;X_2]$.

\begin{lemma}[Distance lower bound]\label{distance-lower}
  \uses{tau-min-def}\leanok
  \lean{distance_ge_of_min}
  For any $G$-valued random variables $X'_1,X'_2$, one has
$$ d[X'_1;X'_2] \geq k - \eta (d[X^0_1;X'_1] - d[X^0_1;X_1] ) - \eta (d[X^0_2;X'_2] - d[X^0_2;X_2] ).$$
\end{lemma}

\begin{proof}
  \uses{tau-def, tau-min}\leanok
  Immediate from \Cref{tau-def} and \Cref{tau-min}.
\end{proof}

\begin{lemma}[Conditional distance lower bound]\label{cond-distance-lower}
  \uses{tau-min-def, cond-dist-def}
  \lean{condRuzsaDistance_ge_of_min}\leanok
  For any $G$-valued random variables $X'_1,X'_2$ and random variables $Z,W$, one has
$$ d[X'_1|Z;X'_2|W] \geq k - \eta (d[X^0_1;X'_1|Z] - d[X^0_1;X_1] ) - \eta (d[X^0_2;X'_2|W] - d[X^0_2;X_2] ).$$
\end{lemma}

\begin{proof}\uses{distance-lower}\leanok  Apply \Cref{distance-lower} to conditioned random variables and then average.
\end{proof}

\section{First estimate}

We continue the assumptions from the preceding section.

Let $X_1, X_2, \tilde X_1, \tilde X_2$ be independent random variables, with $X_1,\tilde X_1$ copies of $X_1$ and $X_2,\tilde X_2$ copies of $X_2$.  (This is possible thanks to \Cref{independent-exist}.)

We also define the quantity
$$ I_1 :=  I [X_1+X_2 : \tilde X_1 + X_2 | X_1+X_2+\tilde X_1+\tilde X_2].$$

\begin{lemma}[Fibring identity for first estimate]\label{first-fibre}
  \lean{rdist_add_rdist_add_condMutual_eq}\leanok
  We have
  \begin{align*}
    &   d[X_1+\tilde X_2;X_2+\tilde X_1] + d[X_1|X_1+\tilde X_2; X_2|X_2+\tilde X_1] \\
    &\quad + \bbI[X_1+ X_2 : \tilde X_1 + X_2 \,|\, X_1 + X_2 + \tilde X_1 + \tilde X_2] = 2k.
  \end{align*}
\end{lemma}

\begin{proof}\uses{cor-fibre} \leanok  Immediate from \Cref{cor-fibre}.
\end{proof}

\begin{lemma}[Lower bound on distances]\label{first-dist-sum}
  \lean{rdist_of_sums_ge}\leanok
  We have
  \begin{align*}
    d[X_1+\tilde X_2; X_2+\tilde X_1] \geq k &- \eta (d[X^0_1; X_1+\tilde X_2] - d[X^0_1; X_1]) \\& \qquad- \eta (d[X^0_2; X_2+\tilde X_1] - d[X^0_2; X_2])
  \end{align*}
\end{lemma}


\begin{proof}\uses{distance-lower}\leanok Immediate from \Cref{distance-lower}.
\end{proof}

\begin{lemma}[Lower bound on conditional distances]\label{first-cond}
  \lean{condRuzsaDist_of_sums_ge}\leanok
  We have
  \begin{align*}
    & d[X_1|X_1+\tilde X_2; X_2|X_2+\tilde X_1]  \\  & \qquad\quad \geq k - \eta (d[X^0_1; X_1 | X_1 + \tilde X_2] - d[X^0_1; X_1]) \\
    & \qquad\qquad\qquad\qquad  - \eta(d[X^0_2; X_2 | X_2 + \tilde X_1] - d[X^0_2; X_2]).
  \end{align*}
\end{lemma}

\begin{proof}\uses{cond-distance-lower}\leanok Immediate from \Cref{cond-distance-lower}.
\end{proof}

\begin{lemma}[Upper bound on distance differences]\label{first-upper}\leanok
  \lean{diff_rdist_le_1, diff_rdist_le_2, diff_rdist_le_3, diff_rdist_le_4}
  We have
  \begin{align*}
    d[X^0_1; X_1+\tilde X_2] - d[X^0_1; X_1] &\leq \tfrac{1}{2} k + \tfrac{1}{4} \bbH[X_2] - \tfrac{1}{4} \bbH[X_1]\\
    d[X^0_2;X_2+\tilde X_1] - d[X^0_2; X_2] &\leq \tfrac{1}{2} k + \tfrac{1}{4} \bbH[X_1] - \tfrac{1}{4} \bbH[X_2], \\
    d[X_1^0;X_1|X_1+\tilde X_2] - d[X_1^0;X_1] &\leq \tfrac{1}{2} k + \tfrac{1}{4} \bbH[X_1] - \tfrac{1}{4} \bbH[X_2]  \\
    d[X_2^0; X_2|X_2+\tilde X_1] - d[X_2^0; X_2] &\leq \tfrac{1}{2}k + \tfrac{1}{4} \bbH[X_2] - \tfrac{1}{4} \bbH[X_1].
  \end{align*}
\end{lemma}

\begin{proof}\uses{first-useful} \leanok  Immediate from \Cref{first-useful} (and recalling that $k$ is defined to be $d[X_1;X_2]$).
\end{proof}

\begin{lemma}[First estimate]\label{first-estimate}
  \lean{first_estimate}\leanok We have $I_1 \leq 2 \eta k$.
\end{lemma}

\begin{proof}\uses{first-fibre, first-dist-sum, first-cond, first-upper}\leanok  Take a suitable linear combination of \Cref{first-fibre}, \Cref{first-dist-sum}, \Cref{first-cond}, and \Cref{first-upper}.
\end{proof}

One can also extract the following useful inequality from the proof of the above lemma.

\begin{lemma}[Entropy bound on quadruple sum]\label{foursum-bound}
  \lean{ent_ofsum_le}\leanok
  With the same notation, we have
  \begin{equation}
    \label{HS-bound}
    \bbH[X_1+X_2+\tilde X_1+\tilde X_2] \le \tfrac{1}{2} \bbH[X_1]+\tfrac{1}{2} \bbH[X_2] + (2 + \eta) k - I_1.
  \end{equation}
\end{lemma}

\begin{proof}\uses{first-cond, first-fibre, first-upper, ruz-indep}\leanok
  Subtracting \Cref{first-cond} from \Cref{first-fibre}, and combining the resulting inequality with \Cref{first-upper} gives the bound
\[
  d[X_1+\tilde X_2;X_2+\tilde X_1] \le (1 + \eta) k - I_1,
\]
and the claim follows from \Cref{ruz-indep} and the definition of $k$.
\end{proof}


\section{Second estimate}

We continue the assumptions from the preceding section.
We introduce the quantity
$$ I_2 := \bbI[X_1+X_2 : X_1 + \tilde X_1 | X_1+X_2+\tilde X_1+\tilde X_2].$$

\begin{lemma}[Distance between sums]\label{dist-sums}
  \lean{rdist_of_sums_ge'}\leanok
We have
$$ d[X_1+\tilde X_1; X_2+\tilde X_2] \geq k - \frac{\eta}{2} ( d[X_1; X_1] + d[X_2;X_2] ).$$
\end{lemma}

\begin{proof}\uses{distance-lower, first-useful}\leanok From \Cref{distance-lower} one has
\begin{align*}
    d[X_1+\tilde X_1; X_2+\tilde X_2] \geq k &- \eta(d[X^0_1;X_1] - d[X^0_1;X_1+\tilde X_1]) \\
    &- \eta(d[X^0_2;X_2] - d[X^0_2;X_2+\tilde X_2]).
\end{align*}
Now \Cref{first-useful} gives
$$
d[X^0_1;X_1+\tilde X_1] - d[X^0_1;X_1] \leq \tfrac{1}{2} d[X_1;X_1]$$
and
$$
  d[X^0_2;X_2+\tilde X_2] - d[X^0_2;X_2] \leq \tfrac{1}{2} d[X_2;X_2],
$$
and the claim follows.
\end{proof}

\begin{lemma}\label{second-estimate-aux}\lean{second_estimate_aux}\leanok
We have
\[d[X_1;X_1] + d[X_2;X_2] \leq 2 k + \frac{2(2 \eta k - I_1)}{1-\eta}. \]
\end{lemma}
\begin{proof}
\uses{ruz-indep, foursum-bound, dist-sums}\leanok
We may use \Cref{ruz-indep} to expand
  \begin{align*}
   & d[X_1+\tilde X_1;X_2+\tilde X_2] \\ &= \bbH[X_1+\tilde X_1 + X_2 + \tilde X_2]  - \tfrac{1}{2} \bbH[X_1+\tilde X_1] - \tfrac{1}{2} \bbH[X_2+\tilde X_2] \\
    &= \bbH[X_1+\tilde X_1 + X_2 + \tilde X_2]  - \tfrac{1}{2} \bbH[X_1] - \tfrac{1}{2} \bbH[X_2]  \\ & \qquad\qquad\qquad   - \tfrac{1}{2} \left( d[X_1;X_1] + d[X_2; X_2] \right),
  \end{align*}
  and hence by \Cref{foursum-bound}
  \[
    d[X_1+\tilde X_1; X_2+\tilde X_2] \leq (2+\eta) k - \tfrac{1}{2} \left( d[X_1;X_1] + d[X_2;X_2] \right) - I_1.
  \]
  Combining this bound with \Cref{dist-sums} we obtain
  the result.
\end{proof}

\begin{lemma}[Second estimate]\label{second-estimate}
  \lean{second_estimate}\leanok
We have
$$ I_2 \leq 2 \eta k + \frac{2 \eta (2 \eta k - I_1)}{1 - \eta}.$$
\end{lemma}

\begin{proof}
  \uses{cor-fibre,distance-lower,first-useful,second-estimate-aux}\leanok
  We apply \Cref{cor-fibre}, but now with the choice
  \[
    (Y_1,Y_2,Y_3,Y_4) := (X_2, X_1, \tilde X_2, \tilde X_1).
  \]
  Now \Cref{cor-fibre} can be rewritten as
  \begin{align*}
    &d[X_1+\tilde X_1;X_2+\tilde X_2] + d[X_1|X_1+\tilde X_1; X_2|X_2+\tilde X_2] \\
    &\quad + \bbI[X_1+X_2 : X_1 + \tilde X_1 \,|\, X_1+X_2+\tilde X_1+\tilde X_2] = 2k,
  \end{align*}
  recalling once again that $k := d[X_1;X_2]$.  From \Cref{cond-distance-lower} one has
  \begin{align*}
    d[X_1|X_1+\tilde X_1; X_2|X_2+\tilde X_2]   \geq k &- \eta (d[X^0_1;X_1] - d[X^0_1;X_1|X_1+\tilde X_1]) \\& - \eta (d[X^0_2;X_2] - d[X^0_2;X_2|X_2+\tilde X_2]) .
  \end{align*}
  while from \Cref{first-useful} we have
  \[
    d[X^0_1;X_1|X_1+\tilde X_1] -  d[X^0_1;X_1] \leq  \tfrac{1}{2} d[X_1;X_1],
  \]
  and
  \[
    d[X^0_2;X_2|X_2+\tilde X_2] -  d[X^0_2;X_2] \leq \tfrac{1}{2} d[X_1;X_2].
    \]
  Combining all these inequalities with \Cref{dist-sums}, we have
  \begin{equation}\label{combined}
  \bbI[X_1+X_2 : X_1 + \tilde X_1 | X_1+X_2+\tilde X_1+\tilde X_2] \leq \eta ( d[X_1; X_1] + d[X_2; X_2] ).
  \end{equation}
  Together with \Cref{second-estimate-aux}, this gives the conclusion.
\end{proof}


\section{Endgame}

Let $X_1,X_2,\tilde X_1,\tilde X_2$ be as before, and introduce the random variables
\[U := X_1 + X_2, \qquad V := \tilde X_1 + X_2, \qquad W := X_1 + \tilde X_1\] and
\[  S := X_1 + X_2 + \tilde X_1 + \tilde X_2.\]

\begin{lemma}[Symmetry identity]\label{symm-lemma}
  \lean{I₃_eq}\leanok
We have
$$ I(U:W | S) = I(V:W | S).$$
\end{lemma}

\begin{proof}\uses{copy-ent, conditional-mutual-alt, chain-rule}\leanok
This should follow from \Cref{copy-ent}, \Cref{conditional-mutual-alt}, and \Cref{chain-rule}.
\end{proof}

\begin{lemma}[Bound on conditional mutual informations]\label{uvw-s}
\uses{conditional-mutual-def}
\lean{I₃_eq,sum_condMutual_le}\leanok
We have
$$
I(U : V \, | \, S) + I(V : W \, | \,S) + I(W : U \, | \, S) \leq 6 \eta k - \frac{1 - 5 \eta}{1-\eta} (2 \eta k - I_1).
$$
\end{lemma}

\begin{proof}
  \uses{second-estimate,symm-lemma}\leanok
From the definitions of $I_1,I_2$ and \Cref{symm-lemma}, we see that
\[
  I_1 = I(U : V \, | \, S), \qquad I_2 = I(W : U \, | \, S), \qquad I_2 = I(V : W \, | \,S).
\]
Applying \Cref{first-estimate} and \Cref{second-estimate} we have the inequalities
\[  I_2 \leq 2 \eta k + \frac{2\eta(2 \eta k - I_1)}{1-\eta} .
\]
We conclude that
$$
   I_1 + I_2 + I_2 \leq I_1+4\eta k+ \frac{4\eta(2 \eta k - I_1)}{1-\eta}
$$
and the claim follows from some calculation.
\end{proof}

\begin{lemma}[Bound on distance increments]\label{total-dist}\leanok
  \lean{sum_dist_diff_le}
   We have
\begin{align*}
  \sum_{i=1}^2 \sum_{A\in\{U,V,W\}} \big(d[X^0_i;A|S] & - d[X^0_i;X_i]\big) \\
    &\leq (6 - 3\eta) k + 3(2 \eta k - I_1).
\end{align*}
\end{lemma}

\begin{proof}
  \uses{second-useful, foursum-bound}\leanok
  By \Cref{second-useful} (taking $X = X_1^0$, $Y = X_1$, $Z = X_2$ and $Z' = \tilde X_1 + \tilde X_2$, so that $Y + Z = U$ and $Y + Z + Z' = S$) we have, noting that $\bbH[Y+ Z] = \bbH[Z']$,
  \[
    d[X^0_1;U|S] - d[X^0_1;X_1] \leq \tfrac{1}{2} (\bbH[S] -  \bbH[X_1]).
  \]
  Further applications of \Cref{second-useful} give
  \begin{align*}
  d[X^0_2;U|S] - d[X^0_2; X_2] &\leq \tfrac{1}{2} (\bbH[S] -  \bbH[X_2]) \\
  d[X^0_1;V|S] - d[X^0_1;X_1] &\leq \tfrac{1}{2} (\bbH[S] -  \bbH[X_1])\\
  d[X^0_2;V|S] - d[X^0_2;X_2] &\leq \tfrac{1}{2} (\bbH[S] -  \bbH[X_2])
  \end{align*}
  and
  \[d[X^0_1;W|S] - d[X^0_1;X_1] \leq \tfrac{1}{2} (\bbH[S] + \bbH[W] - \bbH[X_1] - \bbH[W']),\] where $W' := X_2 + \tilde X_2$.
  To treat $d[X^0_2;W|S]$, first note that this equals $d[X^0_2;W'|S]$, since for a fixed choice $s$ of $S$ we have $W' = W + s$ (here we need some helper lemma about Ruzsa distance). Now we may apply \Cref{second-useful} to obtain
  \[d[X^0_2;W'|S] - d[X^0_2;X_2] \leq \tfrac{1}{2} (\bbH[S] + \bbH[W'] - \bbH[X_2] - \bbH[W]).\]
  Summing these six estimates and using \Cref{foursum-bound}, we conclude that
  \begin{align*}
    \sum_{i=1}^2 \sum_{A\in\{U,V,W\}} \big(d[X^0_i;A|S] & - d[X^0_i;X_i]\big) \\
      &\leq 3 \bbH[S] - \tfrac{3}{2} \bbH[X_1] - \tfrac{3}{2} \bbH[X_2]\\
      &\leq (6 - 3\eta) k + 3(2 \eta k - I_1)
  \end{align*}
  as required.
\end{proof}

\begin{lemma}[Key identity]\label{key-ident}
\lean{sum_uvw_eq_zero}\leanok
We have $U+V+W=0$.
\end{lemma}

\begin{proof} \leanok Obvious because we are in characteristic two.
\end{proof}

For the next two lemmas, let $(T_1,T_2,T_3)$ be a $G^3$-valued random variable such that $T_1+T_2+T_3=0$ holds identically. Set
\begin{equation}\label{delta-t1t2t3-def}
  \delta := \sum_{1 \leq i < j \leq 3} \bbI[T_i;T_j].
\end{equation}

\begin{lemma}[Constructing good variables, I]\label{construct-good-prelim}
  \lean{construct_good_prelim}\leanok
  One has
  \begin{align*}  k \leq
    \delta + \eta (& d[X^0_1;T_1]-d[X^0_1;X_1])
      + \eta (d[X^0_2;T_2]-d[X^0_2;X_2]) \\ & + \tfrac12 \eta \bbI[T_1:T_3] + \tfrac12 \eta \bbI[T_2:T_3].
  \end{align*}
\end{lemma}

(Note: in the paper, this lemma was phrased in a more intuitive formulation that is basically the contrapositive of the one here. Similarly for the next two lemmas.)

\begin{proof}\uses{entropic-bsg, cond-dist-fact,distance-lower}\leanok
  We apply \Cref{entropic-bsg} with $(A,B) = (T_1, T_2)$ there.
  Since $T_1 + T_2 = T_3$, the conclusion is that
  \begin{align} \nonumber \sum_{t_3} \bbP[T_3 = t_3] & d[(T_1 | T_3 = t_3); (T_2 | T_3 = t_3)] \\ & \leq 3 \bbI[T_1 : T_2] + 2 \bbH[T_3] - \bbH[T_1] - \bbH[T_2].\label{bsg-t1t2}
  \end{align}
  The right-hand side in~\eqref{bsg-t1t2} can be rearranged as
  \begin{align*} & 2( \bbH[T_1] + \bbH[T_2] + \bbH[T_3]) - 3 \bbH[T_1,T_2] \\ & = 2(\bbH[T_1] + \bbH[T_2] + \bbH[T_3]) - \bbH[T_1,T_2] - \bbH[T_2,T_3] - \bbH[T_1, T_3] = \delta,\end{align*}
  using the fact (from \Cref{relabeled-entropy}) that all three terms $\bbH[T_i,T_j]$ are equal to $\bbH[T_1,T_2,T_3]$ and hence to each other.
  We also have
  \begin{align*}
  &  \sum_{t_3} P[T_3 = t_3] \bigl(d[X^0_1; (T_1 | T_3=t_3)] - d[X^0_1;X_1]\bigr) \\
  &\quad = d[X^0_1; T_1 | T_3] - d[X^0_1;X_1] \leq d[X^0_1;T_1] - d[X^0_1;X_1] + \tfrac{1}{2} \bbI[T_1 : T_3]
  \end{align*}
  by \Cref{cond-dist-fact}, and similarly
  \begin{align*}
  &  \sum_{t_3} \bbP[T_3 = t_3] (d[X^0_2;(T_2 | T_3=t_3)] - d[X^0_2; X_2]) \\
  &\quad\quad\quad\quad\quad\quad \leq d[X^0_2;T_2] - d[X^0_2;X_2] + \tfrac{1}{2} \bbI[T_2 : T_3].
  \end{align*}
  Putting the above observations together, we have
  \begin{align*}
   \sum_{t_3} \bbP[T_3=t_3] \psi[(T_1 | T_3=t_3); (T_2 | T_3=t_3)] \leq \delta + \eta (d[X^0_1;T_1]-d[X^0_1;X_1]) \\
     + \eta (d[X^0_2;T_2]-d[X^0_2;X_2]) + \tfrac12 \eta \bbI[T_1:T_3] + \tfrac12 \eta \bbI[T_2:T_3]
   \end{align*}
where we introduce the notation
\[\psi[Y_1; Y_2] := d[Y_1;Y_2] +  \eta (d[X_1^0;Y_1] - d[X_1^0;X_1]) + \eta(d[X_2^0;Y_2] - d[X_2^0;X_2]).\]
On the other hand, from \Cref{distance-lower} we have $k \leq \psi[Y_1;Y_2]$, and the claim follows.
  \end{proof}


\begin{lemma}[Constructing good variables, II]\label{construct-good}
  \lean{construct_good}\leanok
One has
\begin{align*}  k & \leq  \delta + \frac{\eta}{3} \biggl( \delta + \sum_{i=1}^2 \sum_{j = 1}^3 (d[X^0_i;T_j] - d[X^0_i; X_i]) \biggr).
  \end{align*}
\end{lemma}

\begin{proof}
\uses{construct-good-prelim}\leanok
Average \Cref{construct-good-prelim} over all six permutations of $T_1,T_2,T_3$.
\end{proof}




\begin{theorem}[$\tau$-decrement]\label{de-prop}
  \lean{tau_strictly_decreases}\leanok
  Let $X_1, X_2$ be tau-minimizers.  Then $d[X_1;X_2] = 0$.
\end{theorem}

\begin{proof}
  \uses{construct-good, key-ident, uvw-s, total-dist, first-estimate, eta-def}\leanok
  Set $k := d[X_1;X_2]$.  Applying \Cref{construct-good} with any random variables $(T_1,T_2,T_3)$ such that $T_1+T_2+T_3=0$ holds identically, we deduce that
\[
  k \leq \delta + \frac{\eta}{3} \biggl( \delta + \sum_{i=1}^2 \sum_{j = 1}^3 (d[X^0_1;T_j] - d[X^0_i;X_i]) \biggr).
\]
Note that $\delta$ is still defined by~\eqref{delta-t1t2t3-def} and thus depends on $T_1,T_2,T_3$.
In particular we may apply this for
\[
  T_1 = (U | S = s),\qquad T_2 = (V | S = s), \qquad T_3 = (W | S = s)
\]
for $s$ in the range of $S$ (which is a valid choice by \Cref{key-ident}) and then average over $s$ with weights $p_S(s)$, to obtain
\[k \leq \tilde \delta + \frac{\eta}{3} \biggl( \tilde \delta + \sum_{i=1}^2 \sum_{A\in\{U,V,W\}} \bigl(  d[X^0_i;A|S] - d[X^0_i;X_i]\bigr) \biggr),\]
where
\[
  \tilde \delta :=  \bbI[U : V | S] + \bbI[V : W | S] + \bbI[W : U | S].
\]
Putting this together with \Cref{uvw-s} and \Cref{total-dist}, we conclude that
\begin{align*}
  k &\leq \Bigl(1+\frac{\eta}{3}\Bigr)\Bigl(6\eta k-\frac{1-5\eta}{1-\eta}(2\eta k-I_1)\Bigr)+\frac{\eta}{3}\Bigl((6-3\eta)k+3(2\eta k-I_1)\Bigr)\\
  &= (8\eta + \eta^2) k - \biggl( \frac{1 - 5 \eta}{1-\eta}\Bigl(1 + \frac{\eta}{3}\Bigr) -  \eta \biggr)(2 \eta k - I_1)\\
  &\leq (8 \eta + \eta^2) k
 \end{align*}
since the quantity $2 \eta k - I_1$ is non-negative (by \Cref{first-estimate}), and its coefficient in the above expression is non-positive provided that $\eta(2\eta + 17) \le 3$, which is certainly the case with \Cref{eta-def}.
Moreover, from \Cref{eta-def} we have $8 \eta + \eta^2 < 1$. It follows that $k=0$, as desired.
\end{proof}


\section{Conclusion}

\begin{theorem}[Entropy version of PFR]\label{entropy-pfr}
  \lean{entropic_PFR_conjecture}\leanok
  Let $G = \F_2^n$, and suppose that $X^0_1, X^0_2$ are $G$-valued random variables.
  Then there is some subgroup $H \leq G$ such that
  \[
    d[X^0_1;U_H] + d[X^0_2;U_H] \le 11 d[X^0_1;X^0_2],
  \]
  where $U_H$ is uniformly distributed on $H$.
  Furthermore, both $d[X^0_1;U_H]$ and $d[X^0_2;U_H]$ are at most $6 d[X^0_1;X^0_2]$.
\end{theorem}

\begin{proof} \uses{de-prop, tau-min, lem:100pc, ruzsa-triangle} \leanok  Let $X_1, X_2$ be the $\tau$-minimizer from \Cref{tau-min}.  From \Cref{de-prop}, $d[X_1;X_2]=0$.  From \Cref{lem:100pc}, $d[X_1;U_H] = d[X_2; U_H] = 0$.  Also from $\tau$-minimization we have $\tau[X_1;X_2] \leq \tau[X^0_2;X^0_1]$.  Using this and the Ruzsa triangle inequality we can conclude.
\end{proof}

Note: a ``stretch goal'' for this project would be to obtain a `decidable` analogue of this result (see the remark at the end of Section 2 for some related discussion).


%%FILE blueprint/src/chapter/fibring.tex
\chapter{The Fibring lemma}

\begin{proposition}[General fibring identity]\label{fibring-ident}
  \lean{rdist_of_indep_eq_sum_fibre, rdist_le_sum_fibre}\leanok
  Let $\pi : H \to H'$ be a homomorphism additive groups, and let $Z_1,Z_2$ be $H$-valued random variables. Then we have
  \[
    d[Z_1; Z_2] \geq d[\pi(Z_1);\pi(Z_2)] + d[Z_1|\pi(Z_1); Z_2 |\pi(Z_2)].
  \]
  Moreover, if $Z_1,Z_2$ are taken to be independent, then the difference between the two sides is
$$I( Z_1 - Z_2 : (\pi(Z_1), \pi(Z_2))  |  \pi(Z_1 - Z_2) ).$$
\end{proposition}


\begin{proof}\uses{ruz-copy, independent-exist, submodularity,conditional-mutual-alt,chain-rule,relabeled-entropy, cond-dist-alt}\leanok
  Let $Z_1,Z_2$ be independent throughout (this is possible by \Cref{ruz-copy} and \Cref{independent-exist}).  By \Cref{cond-dist-alt}, We have
  \begin{align*}
    & d[Z_1  |\pi(Z_1); Z_2 |\pi(Z_2)] \\
    & = \bbH[Z_1 - Z_2 | \pi(Z_1),\pi(Z_2)] - \tfrac{1}{2} \bbH[Z_1 | \pi(Z_1)] - \tfrac{1}{2} \bbH[Z_2 | \pi(Z_2)] \\
    & \leq  \bbH[Z_1 - Z_2 | \pi(Z_1+Z_2)]  - \tfrac{1}{2} \bbH[Z_1 | \pi(Z_1)] - \tfrac{1}{2}H[Z_2 | \pi(Z_2)] \\
    & = d[Z_1;Z_2] - d[\pi(Z_1);\pi(Z_2)].
  \end{align*}
  In the middle step, we used \Cref{submodularity}, and in the last step we used the fact that
  \[\bbH[Z_1 - Z_2 |  \pi(Z_1-Z_2)] = \bbH[Z_1 - Z_2] - \bbH[\pi(Z_1-Z_2)]\]
  (thanks to \Cref{chain-rule} and \Cref{relabeled-entropy}) and that
  \[\bbH[Z_i| \pi(Z_i)] = \bbH[Z_i] - \bbH[\pi(Z_i)]\] (since $Z_i$ determines $\pi(Z_i)$).
  This gives the claimed inequality. The difference between the two sides is precisely
  \[\bbH[Z_1 - Z_2  | \pi(Z_1 - Z_2)] - \bbH[Z_1 - Z_2  | \pi(Z_1),\pi(Z_2)].\]
  To rewrite this in terms of (conditional) mutual information, we use the identity
  \[\bbH[A|B] - \bbH[A | B,C] = \bbI[A : C | B],\]
  (which follows \Cref{conditional-mutual-alt})
  taking
  $A := Z_1 - Z_2$, $B := \pi(Z_1 - Z_2)$ and $C := (\pi(Z_1),\pi(Z_{2}))$, and noting that in this case $\bbH[A | B,C] = \bbH[A | C]$ since $C$ uniquely determines $B$ (this may require another helper lemma about entropy).
  This completes the proof.
\end{proof}


\begin{corollary}\label{fibring-ineq}
 \lean{rdist_of_hom_le} \leanok
If $\pi:G\to H$ is a homomorphism of additive groups and $X,Y$ are $G$-valued random variables then
\[d[X;Y]\geq d[\pi(X);\pi(Y)].\]
\end{corollary}

\begin{proof}
\uses{fibring-ident, ruzsa-nonneg}\leanok
By \Cref{fibring-ident} and the nonnegativity of conditional Ruzsa distance (from \Cref{ruzsa-nonneg}) we have
\[d[X;Y]\geq d[\pi(X);\pi(Y)]+d[X\mid \pi(X);Y\mid \pi(Y)].\]
The inequality follows from $d[X\mid \pi(X);Y\mid \pi(Y)]\geq 0$ (\Cref{ruzsa-nonneg}).
\end{proof}


\begin{corollary}[Specific fibring identity]\label{cor-fibre}
  \lean{sum_of_rdist_eq}\leanok
  Let $Y_1,Y_2,Y_3$ and $Y_4$ be independent $G$-valued random variables.
  Then
\begin{align*}
  & d[Y_1+Y_3; Y_2+Y_4] + d[Y_1|Y_1+Y_3; Y_2|Y_2+Y_4] \\
  &\qquad + \bbI[Y_1+Y_2 : Y_2 + Y_4 | Y_1+Y_2+Y_3+Y_4] = d[Y_1; Y_2] + d[Y_3; Y_4].
\end{align*}
\end{corollary}

\begin{proof}
  \uses{fibring-ident, add-entropy, relabeled-entropy-cond}\leanok
  We apply \Cref{fibring-ident} with $H := G \times G$, $H' := G$, $\pi$ the addition homomorphism $\pi(x,y) := x+y$, and with the random variables $Z_1 := (Y_1,Y_3)$ and $Z_2 := (Y_2,Y_4)$.  Then by independence (\Cref{add-entropy})
\[
  d[Z_1; Z_2] = d[Y_1; Y_2] + d[Y_3; Y_4]
\]
while by definition
\[
  d[\pi(Z_1); \pi(Z_2)] = d[Y_1+Y_3; Y_2+Y_4].
\]
Furthermore,
\[
  d[Z_1|\pi(Z_1); Z_2|\pi(Z_2)] = d[Y_1|Y_1+Y_3;Y_2|Y_2+Y_4],
\]
since $Z_1=(Y_1,Y_3)$ and $Y_1$ are linked by an invertible affine transformation once $\pi(Z_1)=Y_1+Y_3$ is fixed, and similarly for $Z_2$ and $Y_2$.  (This has to do with \Cref{relabeled-entropy-cond})
Finally, we have
\begin{align*}
  &\bbI[Z_1 + Z_2 : (\pi(Z_1),\pi(Z_2)) \,|\, \pi(Z_1) + \pi(Z_2)] \\
  &\ = \bbI[(Y_1+Y_2, Y_3+Y_4) : (Y_1+Y_3, Y_2+Y_4) \,|\, Y_1+Y_2+Y_3+Y_4] \\
  &\ = \bbI[Y_1+Y_2 : Y_2+Y_4 \,|\, Y_1+Y_2+Y_3+Y_4]
\end{align*}
where in the last line we used the fact that $(Y_1+Y_2, Y_1+Y_2+Y_3+Y_4)$ uniquely determine $Y_3+Y_4$ and similarly
$(Y_2+Y_4, Y_1+Y_2+Y_3+Y_4)$ uniquely determine $Y_1+Y_3$.  (This requires another helper lemma about entropy.)
\end{proof}


%%FILE blueprint/src/chapter/further_improvement.tex
\chapter{Further improvement to exponent}

\section{Kullback--Leibler divergence}

In the definitions below, $G$ is a set.

\begin{definition}[Kullback--Leibler divergence]\label{kl-div}\lean{KLDiv}\leanok If $X,Y$ are two $G$-valued random variables, the Kullback--Leibler divergence is defined as
  $$ D_{KL}(X\Vert Y) := \sum_x \mathbf{P}(X=x) \log \frac{\mathbf{P}(X=x)}{\mathbf{P}(Y=x)}.$$
\end{definition}

\begin{lemma}[Kullback--Leibler divergence of copy]
  \label{kl-div-copy}
  \lean{ProbabilityTheory.IdentDistrib.KLDiv_eq}
  \leanok

  If $X'$ is a copy of $X$, and $Y'$ is a copy of $Y$, then $D_{KL}(X'\Vert Y') = D_{KL}(X\Vert Y)$.
\end{lemma}

\begin{proof}\uses{kl-div}\leanok  Clear from definition.
\end{proof}

\begin{lemma}[Gibbs inequality]\label{Gibbs}\uses{kl-div}\lean{KLDiv_nonneg}\leanok  $D_{KL}(X\Vert Y) \geq 0$.
\end{lemma}

\begin{proof}\leanok
  \uses{log-sum}
  Apply \Cref{log-sum} on the definition.
\end{proof}

\begin{lemma}[Converse Gibbs inequality]\label{Gibbs-converse}\lean{KLDiv_eq_zero_iff_identDistrib}\leanok  If $D_{KL}(X\Vert Y) = 0$, then $Y$ is a copy of $X$.
\end{lemma}

\begin{proof}\leanok
  \uses{converse-log-sum}
  Apply \Cref{converse-log-sum}.
\end{proof}

\begin{lemma}[Convexity of Kullback--Leibler]\label{kl-div-convex}\lean{KLDiv_of_convex}\leanok  If $S$ is a finite set, $\sum_{s \in S} w_s = 1$ for some non-negative $w_s$, and ${\bf P}(X=x) = \sum_{s\in S} w_s  {\bf P}(X_s=x)$, ${\bf P}(Y=x) = \sum_{s\in S} w_s  {\bf P}(Y_s=x)$ for all $x$, then
$$D_{KL}(X\Vert Y) \le \sum_{s\in S} w_s D_{KL}(X_s\Vert Y_s).$$
\end{lemma}

\begin{proof}\leanok
  \uses{kl-div,log-sum}
For each $x$, replace $\log \frac{\mathbf{P}(X_s=x)}{\mathbf{P}(Y_s=x)}$ in the definition with $\log \frac{w_s\mathbf{P}(X_s=x)}{w_s\mathbf{P}(Y_s=x)}$ for each $s$, and apply \Cref{log-sum}.
\end{proof}


\begin{lemma}[Kullback--Leibler and injections]\label{kl-div-inj}\lean{KLDiv_of_comp_inj}\leanok  If $f:G \to H$ is an injection, then $D_{KL}(f(X)\Vert f(Y)) = D_{KL}(X\Vert Y)$.
\end{lemma}

\begin{proof}\leanok\uses{kl-div} Clear from definition.
\end{proof}

Now let $G$ be an additive group.

\begin{lemma}[Kullback--Leibler and sums]\label{kl-sums}\lean{KLDiv_add_le_KLDiv_of_indep}\leanok If $X, Y, Z$ are independent $G$-valued random variables, then
  $$D_{KL}(X+Z\Vert Y+Z) \leq D_{KL}(X\Vert Y).$$
\end{lemma}

\begin{proof}\leanok
  \uses{kl-div-inj,kl-div-convex}
For each $z$, $D_{KL}(X+z\Vert Y+z)=D_{KL}(X\Vert Y)$ by \Cref{kl-div-inj}. Then apply \Cref{kl-div-convex} with $w_z=\mathbf{P}(Z=z)$.
\end{proof}

\begin{definition}[Conditional Kullback--Leibler divergence]\label{ckl-div}  \uses{kl-div}\leanok
  If $X,Y,Z$ are random variables, with $X,Z$ defined on the same sample space, we define
$$ D_{KL}(X|Z \Vert  Y) := \sum_z \mathbf{P}(Z=z) D_{KL}( (X|Z=z) \Vert  Y).$$
\end{definition}

\begin{lemma}[Kullback--Leibler and conditioning]\label{kl-cond}\lean{condKLDiv_eq}\leanok If $X, Y$ are independent $G$-valued random variables, and $Z$ is another random variable defined on the same sample space as $X$, then
  $$D_{KL}((X|Z)\Vert Y) = D_{KL}(X\Vert Y) + \bbH[X] - \bbH[X|Z].$$
\end{lemma}

\begin{proof}\leanok
  \uses{ckl-div} Compare the terms correspond to each $x\in G$ on both sides.
\end{proof}

\begin{lemma}[Conditional Gibbs inequality]\label{Conditional-Gibbs}\lean{condKLDiv_nonneg}\leanok  $D_{KL}((X|W)\Vert Y) \geq 0$.
\end{lemma}

\begin{proof}\leanok \uses{Gibbs, ckl-div}  Clear from Definition \ref{ckl-div} and Lemma \ref{Gibbs}.
\end{proof}

\section{Rho functionals}

Let $G$ be an additive group, and let $A$ be a non-empty subset of $G$.

\begin{definition}[Rho minus]\label{rhominus-def}\lean{rhoMinus}\leanok  For any $G$-valued random variable $X$, we define $\rho^-(X)$ to be the infimum of $D_{KL}(X \Vert  U_A + T)$, where $U_A$ is uniform on $A$ and $T$ ranges over $G$-valued random variables independent of $U_A$.
\end{definition}

\begin{definition}[Rho plus]\label{rhoplus-def}\lean{rhoPlus}\uses{rhominus-def}\leanok  For any $G$-valued random variable $X$, we define $\rho^+(X) := \rho^-(X) + \bbH(X) - \bbH(U_A)$.
\end{definition}

\begin{lemma}[Rho minus non-negative]\label{rhominus-nonneg}\label{rhoMinus_nonneg}\uses{rhominus-def}\leanok  We have $\rho^-(X) \geq 0$.
\end{lemma}

\begin{proof}\leanok \uses{Conditional-Gibbs} Clear from Lemma \ref{Conditional-Gibbs}.
\end{proof}

\begin{lemma}[Rho minus of subgroup]\label{rhominus-subgroup}\lean{rhoMinus_of_subgroup}\uses{rhominus-def}\leanok If $H$ is a finite subgroup of $G$, then $\rho^-(U_H) = \log |A| - \log \max_t |A \cap (H+t)|$.
\end{lemma}

\begin{proof}\leanok
  \uses{log-sum}
  For every $G$-valued random variable $T$ that is independent of $Y$,
  $$D_{KL}(U_H \Vert U_A+T) = \sum_{h\in H} \frac{1}{|H|}\log\frac{1/|H|}{\mathbf{P}[U_A+T=h]}\ge -\log(\mathbf{P}[U_A+T\in H]),$$
  by \Cref{log-sum}. Then observe that $$-\log(\mathbf{P}[U_A+T\in H])=-\log(\mathbf{P}[U_A\in H-T])\ge -\log(\max_{t\in G} \mathbf{P}[U_A\in H+t]).$$ This proves $\ge$.

  To get the equality, let $t^*:=\arg\max_t |A \cap (H+t)|$ and observe that $$\rho^-(U_H)\le D_{KL}(U_H \Vert  U_A+(U_H-t^*))= \log |A| - \log \max_t|A \cap (H+t)|.$$
\end{proof}

\begin{corollary}[Rho plus of subgroup]\label{rhoplus-subgroup}\lean{rhoPlus_of_subgroup}\uses{rhoplus-def}\leanok If $H$ is a finite subgroup of $G$, then $\rho^+(U_H) = \log |H| - \log \max_t |A \cap (H+t)|$.
\end{corollary}

\begin{proof}\leanok \uses{rhominus-subgroup} Straightforward by definition and \Cref{rhominus-subgroup}.
\end{proof}

\begin{definition}[Rho functional]\label{rho-def}\lean{rho}\uses{rhoplus-def}\leanok  We define $\rho(X) := (\rho^+(X) + \rho^-(X))/2$.
\end{definition}

\begin{lemma}\label{rho-init}\label{rho_of_uniform}\uses{rho-def}\leanok  We have $\rho(U_A) = 0$.
\end{lemma}

\begin{proof}\leanok \uses{rhominus-nonneg} $\rho^-(U_A)\le 0$ by the choice $T=0$. The claim then follows from \Cref{rhominus-nonneg}.
\end{proof}

\begin{lemma}[Rho of subgroup]\label{rho-subgroup}\uses{rho-def}\lean{rho_of_subgroup}\leanok
If $H$ is a finite subgroup of $G$, and $\rho(U_H) \leq r$, then there exists
$t$ such that $|A \cap (H+t)| \geq e^{-r} \sqrt{|A||H|}$, and
$|H|/|A|\in[e^{-2r},e^{2r}]$.
\end{lemma}

\begin{proof}\leanok\uses{rhominus-subgroup, rhoplus-subgroup} The first claim is a direct corollary of \Cref{rhominus-subgroup} and \Cref{rhoplus-subgroup}. To see the second claim, observe that \Cref{rhominus-nonneg} and \Cref{rhoplus-subgroup} imply $\rho^-(U_H),\rho^+(U_H)\ge 0$. Therefore $$|H(U_A)-H(U_H)|=|\rho^+(U_H)-\rho^-(U_H)|\le \rho^-(U_H)+\rho^+(U_H)= 2\rho(U_H)\le 2r,$$
  which implies the second claim.
\end{proof}

\begin{lemma}[Rho invariant]\label{rho-invariant}\lean{rho_of_translate}\leanok\uses{rho-def}  For any $s \in G$, $\rho(X+s) = \rho(X)$.
\end{lemma}

\begin{proof}\leanok\uses{kl-div-inj} Observe that by \Cref{kl-div-inj}, $$\inf_T D_{KL}(X\Vert U_A+T)=\inf_T D_{KL}(X+s\Vert U_A+T+s)=\inf_{T'} D_{KL}(X+s\Vert U_A+T').$$
\end{proof}

\begin{lemma}[Rho continuous]\label{rho-cts}\lean{rho_continuous}\leanok\uses{rho-def} $\rho(X)$ depends continuously on the distribution of $X$.
\end{lemma}

\begin{proof}  \leanok Clear from definition.
\end{proof}

\begin{lemma}[Rho and sums]\label{rho-sums}\lean{rhoMinus_of_sum, rhoPlus_of_sum, rho_of_sum}\leanok  If $X,Y$ are independent, one has
  $$ \rho^-(X+Y) \leq \rho^-(X)$$
  $$ \rho^+(X+Y) \leq \rho^+(X) + \bbH[X+Y] - \bbH[X]$$
and
  $$ \rho(X+Y) \leq \rho(X) + \frac{1}{2}( \bbH[X+Y] - \bbH[X] ).$$
\end{lemma}

\begin{proof}\leanok
  \uses{kl-sums}
The first inequality follows from \Cref{kl-sums}. The second and third inequalities are direct corollaries of the first.
\end{proof}


\begin{definition}[Conditional Rho functional]\label{rho-cond-def}\lean{condRho}\leanok We define $\rho(X|Y) := \sum_y {\bf P}(Y=y) \rho(X|Y=y)$.
\end{definition}

\begin{lemma}[Conditional rho and translation]\label{rho-cond-invariant}\lean{condRho_of_translate}\leanok
  For any $s\in G$, $\rho(X+s|Y)=\rho(X|Y)$.
\end{lemma}
\begin{proof}
  \uses{rho-cond-def,rho-invariant}\leanok
  Direct corollary of \Cref{rho-invariant}.
\end{proof}

\begin{lemma}[Conditional rho and relabeling]\label{rho-cond-relabeled}\lean{condRho_of_injective}\leanok
  If $f$ is injective, then $\rho(X|f(Y))=\rho(X|Y)$.
\end{lemma}
\begin{proof}\leanok
  \uses{rho-cond-def}
  Clear from the definition.
\end{proof}

\begin{lemma}[Rho and conditioning]\label{rho-cond}\lean{condRhoMinus_le, condRhoPlus_le, condRho_le}\leanok  If $X,Z$ are defined on the same space, one has
  $$ \rho^-(X|Z) \leq \rho^-(X) + \bbH[X] - \bbH[X|Z]$$
  $$ \rho^+(X|Z) \leq \rho^+(X)$$
and
  $$ \rho(X|Z) \leq \rho(X) + \frac{1}{2}( \bbH[X] - \bbH[X|Z] ).$$
\end{lemma}

\begin{proof}\leanok
  \uses{kl-cond}
  The first inequality follows from \Cref{kl-cond}. The second and third inequalities are direct corollaries of the first.
\end{proof}

The following lemmas hold for $G=\mathbb{F}_2^n$.

\begin{lemma}[Rho and sums, symmetrized]\label{rho-sums-sym}\lean{rho_of_sum_le}\leanok  If $X,Y$ are independent, then
  $$ \rho(X+Y) \leq \frac{1}{2}(\rho(X)+\rho(Y) + d[X;Y]).$$
\end{lemma}
\begin{proof}\leanok
  \uses{rho-sums}
Apply \Cref{rho-sums} for $(X,Y)$ and $(Y,X)$ and take their average.
\end{proof}

\begin{lemma}[Rho and conditioning, symmetrized]\label{rho-cond-sym}\lean{condRho_of_sum_le}\leanok
  If $X,Y$ are independent, then
  $$ \rho(X | X+Y) \leq \frac{1}{2}(\rho(X)+\rho(Y) + d[X;Y]).$$
\end{lemma}
\begin{proof}\leanok
  \uses{rho-invariant,rho-cond}
  First apply \Cref{rho-cond} to get $\rho(X|X+Y)\le \rho(X) + \frac{1}{2}(\bbH[X+Y]-\bbH[Y])$, and $\rho(Y|X+Y)\le \rho(Y)+\frac{1}{2}(\bbH[X+Y]-\bbH[X])$. Then apply \Cref{rho-invariant} to get $\rho(Y|X+Y)=\rho(X|X+Y)$ and take the average of the two inequalities.
\end{proof}

\section{Studying a minimizer}

Set $\eta < 1/8$. In this section, consider $G=\mathbb{F}_2^n$.

\begin{definition}\label{phi-min-def}\lean{phiMinimizes}\leanok  Given $G$-valued random variables $X,Y$, define
$$ \phi[X;Y] := d[X;Y] + \eta(\rho(X) + \rho(Y))$$
and define a \emph{$\phi$-minimizer} to be a pair of random variables $X,Y$ which minimizes $\phi[X;Y]$.
\end{definition}

\begin{lemma}[$\phi$-minimizers exist]\label{phi-min-exist}\lean{phi_min_exists}\leanok  There exists a $\phi$-minimizer.
\end{lemma}

\begin{proof}\leanok\uses{rho-cts} Clear from compactness.
\end{proof}

Let $(X_1, X_2)$ be a $\phi$-minimizer, and $\tilde X_1, \tilde X_2$ be independent copies of $X_1,X_2$ respectively.
Similar to the original proof we define
$$ I_1 :=  I [X_1+X_2 : \tilde X_1 + X_2 | X_1+X_2+\tilde X_1+\tilde X_2], I_2 := \bbI[X_1+X_2 : X_1 + \tilde X_1 | X_1+X_2+\tilde X_1+\tilde X_2].$$
First we need the $\phi$-minimizer variants of \Cref{first-estimate} and \Cref{second-estimate}.

\begin{lemma}\label{phi-first-estimate}\lean{I_one_le}\leanok
$I_1\le 2\eta d[X_1;X_2]$
\end{lemma}

\begin{proof}\leanok
  \uses{phi-min-def,first-fibre}
Similar to \Cref{first-estimate}: get upper bounds for $d[X_1;X_2]$ by $\phi[X_1;X_2]\le \phi[X_1+X_2;\tilde X_1+\tilde X_2]$ and $\phi[X_1;X_2]\le \phi[X_1|X_1+X_2;\tilde X_2|\tilde X_1+\tilde X_2]$, and then apply \Cref{first-fibre} to get an upper bound for $I_1$.
\end{proof}

\begin{lemma}\label{I1-I2-diff}\lean{rdist_add_rdist_eq}\leanok
  $d[X_1;X_1]+d[X_2;X_2]= 2d[X_1;X_2]+(I_2-I_1)$.
\end{lemma}
\begin{proof}\leanok
  \uses{first-fibre,cor-fibre}
Compare \Cref{first-fibre} with the identity obtained from applying \Cref{cor-fibre} on $(X_1,\tilde X_1, X_2, \tilde X_2)$.
\end{proof}

\begin{lemma}\label{phi-second-estimate}\lean{I_two_le}\leanok
$I_2\le 2\eta d[X_1;X_2] + \frac{\eta}{1-\eta}(2\eta d[X_1;X_2]-I_1)$.
\end{lemma}

\begin{proof}\leanok
  \uses{phi-min-def,cor-fibre,I1-I2-diff}
First of all, by $\phi[X_1;X_2]\le \phi[X_1+\tilde X_1;X_2+\tilde X_2]$, $\phi[X_1;X_2]\le \phi[X_1|X_1+\tilde X_1;X_2|X_2+\tilde X_2]$, and the fibring identity obtained by applying \Cref{cor-fibre} on $(X_1,X_2,\tilde X_1,\tilde X_2)$,
we have $I_2\le \eta (d[X_1;X_1]+d[X_2;X_2])$. Then apply \Cref{I1-I2-diff} to get $I_2\le 2\eta d[X_1;X_2] +\eta(I_2-I_1)$, and rearrange.
\end{proof}


Next we need some inequalities for the endgame.

\begin{lemma}\label{rho-BSG-triplet}\lean{dist_le_of_sum_zero}\leanok
  If $G$-valued random variables $T_1,T_2,T_3$ satisfy $T_1+T_2+T_3=0$, then
  $$d[X_1;X_2]\le 3\bbI[T_1:T_2] + (2\bbH[T_3]-\bbH[T_1]-\bbH[T_2])+ \eta(\rho(T_1|T_3)+\rho(T_2|T_3)-\rho(X_1)-\rho(X_2)).$$
\end{lemma}

\begin{proof}\leanok\uses{entropic-bsg,phi-min-def}
  Conditioned on every $T_3=t$, $d[X_1;X_2]\le d[T_1|T_3=t;T_2|T_3=t]+\eta(\rho(T_1|T_3=t)+\rho(T_2|T_3=t)-\rho(X_1)-\rho(X_2))$ by \Cref{phi-min-def}. Then take the weighted average with weight $\mathbf{P}(T_3=t)$ and then apply \Cref{entropic-bsg} to bound the RHS.
\end{proof}

\begin{lemma}\label{rho-BSG-triplet-symmetrized}\lean{dist_le_of_sum_zero'}\leanok
  If $G$-valued random variables $T_1,T_2,T_3$ satisfy $T_1+T_2+T_3=0$, then
  $$d[X_1;X_2] \leq  \sum_{1 \leq i<j \leq 3} \bbI[T_i:T_j] + \frac{\eta}{3}   \sum_{1 \leq i<j \leq 3} (\rho(T_i|T_j) + \rho(T_j|T_i) -\rho(X_1)-\rho(X_2))$$
\end{lemma}

\begin{proof}\leanok\uses{rho-BSG-triplet}
  Take the average of \Cref{rho-BSG-triplet} over all $6$ permutations of $T_1,T_2,T_3$.
\end{proof}


\begin{lemma}\label{rho-increase}\lean{condRho_sum_le}\leanok
  For independent random variables $Y_1,Y_2,Y_3,Y_4$ over $G$, define $S:=Y_1+Y_2+Y_3+Y_4$, $T_1:=Y_1+Y_2$, $T_2:=Y_1+Y_3$. Then
  $$\rho(T_1|T_2,S)+\rho(T_2|T_1,S) - \frac{1}{2}\sum_{i} \rho(Y_i)\le \frac{1}{2}(d[Y_1;Y_2]+d[Y_3;Y_4]+d[Y_1;Y_3]+d[Y_2;Y_4]).$$
\end{lemma}

\begin{proof}\leanok\uses{rho-sums-sym, rho-cond, rho-cond-sym, rho-cond-relabeled, cor-fibre}
  Let $T_1':=Y_3+Y_4$, $T_2':=Y_2+Y_4$.
  First note that
  \begin{align*}
    \rho(T_1|T_2,S)
    &\le \rho(T_1|S) + \frac{1}{2}\bbI(T_1:T_2\mid S)  \\
    &\le \frac{1}{2}(\rho(T_1)+\rho(T_1'))+\frac{1}{2}(d[T_1;T_1']+\bbI(T_1:T_2\mid S)) \\
    &\le \frac{1}{4} \sum_{i} \rho(Y_i) +\frac{1}{4}(d[Y_1;Y_2]+d[Y_3;Y_4]) + \frac{1}{2}(d[T_1;T_1']+\bbI(T_1:T_2\mid S)).
  \end{align*}
  by \Cref{rho-cond}, \Cref{rho-cond-sym}, \Cref{rho-sums-sym} respectively.
  On the other hand, observe that
  \begin{align*}
    \rho(T_1|T_2,S)
    &=\rho(Y_1+Y_2|T_2,T_2') \\
    &\le \frac{1}{2}(\rho(Y_1|T_2)+\rho(Y_2|T_2'))+\frac{1}{2}(d[Y_1|T_2;Y_2|T_2'])  \\
    &\le \frac{1}{4} \sum_{i} \rho(Y_i) +\frac{1}{4}(d[Y_1;Y_3]+d[Y_2;Y_4]) + \frac{1}{2}(d[Y_1|T_2;Y_2|T_2']).
  \end{align*}
  by \Cref{rho-cond-relabeled}, \Cref{rho-sums-sym}, \Cref{rho-cond-sym} respectively.
  By replacing $(Y_1,Y_2,Y_3,Y_4)$ with $(Y_1,Y_3,Y_2,Y_4)$ in the above inequalities, one has
  $$\rho(T_2|T_1,S) \le \frac{1}{4} \sum_{i} \rho(Y_i) +\frac{1}{4}(d[Y_1;Y_3]+d[Y_2;Y_4]) + \frac{1}{2}(d[T_2;T_2']+\bbI(T_1:T_2\mid S))$$
  and
  $$\rho(T_2|T_1,S) \le \frac{1}{4} \sum_{i} \rho(Y_i) +\frac{1}{4}(d[Y_1;Y_2]+d[Y_3;Y_4]) + \frac{1}{2}(d[Y_1|T_1;Y_3|T_1']).$$
  Finally, take the sum of all four inequalities, apply \Cref{cor-fibre} on $(Y_1,Y_2,Y_3,Y_4)$ and $(Y_1,Y_3,Y_2,Y_4)$ to rewrite the sum of last terms in the four inequalities, and divide the result by $2$.
\end{proof}

\begin{lemma}\label{rho-increase-symmetrized}\lean{condRho_sum_le'}\leanok
  For independent random variables $Y_1,Y_2,Y_3,Y_4$ over $G$, define $T_1:=Y_1+Y_2,T_2:=Y_1+Y_3,T_3:=Y_2+Y_3$ and $S:=Y_1+Y_2+Y_3+Y_4$. Then
  $$\sum_{1 \leq i<j \leq 3} (\rho(T_i|T_j,S) + \rho(T_j|T_i,S) -  \frac{1}{2}\sum_{i} \rho(Y_i))\le \sum_{1\leq i < j \leq 4}d[Y_i;Y_j]$$
\end{lemma}

\begin{proof}\uses{rho-increase}\leanok
  Apply Lemma \ref{rho-increase} on $(Y_i,Y_j,Y_k,Y_4)$ for $(i,j,k)=(1,2,3),(2,3,1),(1,3,2)$, and take the sum.
\end{proof}

\begin{proposition}\label{phi-minimizer-zero-distance}\lean{dist_of_min_eq_zero}\leanok  If $X_1,X_2$ is a $\phi$-minimizer, then $d[X_1;X_2] = 0$.
\end{proposition}

\begin{proof}\leanok
  \uses{rho-BSG-triplet-symmetrized,rho-increase-symmetrized,I1-I2-diff,phi-first-estimate,phi-second-estimate}
Consider $T_1:=X_1+X_2,T_2:=X_1+\tilde X_1, T_3:=\tilde X_1 + X_2$, and $S=X_1+X_2+\tilde X_1 + \tilde X_2$. Note that $T_1+T_2+T_3=0$.
First apply \Cref{rho-BSG-triplet-symmetrized} on $(T_1,T_2,T_3)$ when conditioned on $S$ to get
\begin{align}  \label{eq:further-bsg}
d[X_1;X_2] &\leq  \sum_{1 \leq i<j \leq 3} \bbI[T_i:T_j\mid S] + \frac{\eta}{3}   \sum_{1 \leq i<j \leq 3} (\rho(T_i|T_j,S) + \rho(T_j|T_i,S) -\rho(X_1)-\rho(X_2))\nonumber\\
&= (I_1+2I_2) + \frac{\eta}{3}   \sum_{1 \leq i<j \leq 3} (\rho(T_i|T_j,S) + \rho(T_j|T_i,S) -\rho(X_1)-\rho(X_2)).
\end{align}
Then apply \Cref{rho-increase-symmetrized} on $(X_1,X_2,\tilde X_1,\tilde X_2)$ and get
$$\sum_{1 \leq i<j \leq 3} (\rho(T_i|T_j,S) + \rho(T_j|T_i,S) -  \rho(X_1)-\rho(X_2))\le (4d[X_1;X_2]+d[X_1;X_2]+d[X_2;X_2])= 6 d[X_1;X_2]+(I_2-I_1)$$
by \Cref{I1-I2-diff}. Plug in the inequality above to (\ref{eq:further-bsg}), we get
$$d[X_1;X_2] \le (I_1+2I_2)+2\eta d[X_1;X_2]+\frac{\eta}{3}(I_2-I_1).$$
By \Cref{phi-second-estimate} we can conclude that
$$d[X_1;X_2] \le 8\eta d[X_1;X_2]-\frac{3-10\eta}{3-3\eta} (2\eta d[X_1;X_2]-I_1).$$
Finally by \Cref{phi-first-estimate} and $\eta<1$ we get that the second term is $\le 0$, and thus $d[X_1;X_2] \le 8\eta d[X_1;X_2]$. By the choice $\eta<1/8$ and the non-negativity of $d$ we have $d[X_1;X_2]=0$.
\end{proof}


\begin{proposition} \label{pfr-rho}\lean{rho_PFR_conjecture}\leanok For any random variables $Y_1,Y_2$, there exist a subgroup $H$ such that
  $$ 2\rho(U_H) \leq \rho(Y_1) + \rho(Y_2) + 8 d[Y_1;Y_2].$$
\end{proposition}

\begin{proof}\leanok
  \uses{phi-min-exist, phi-minimizer-zero-distance,phi-min-def,rho-invariant,sym-zero}
Let $X_1,X_2$ be a $\phi$-minimizer. By \Cref{phi-minimizer-zero-distance} $d[X_1;X_2]=0$, which by \Cref{phi-min-def} implies $\rho(X_1)+\rho(X_2)\le \rho(Y_1) + \rho(Y_2) + \frac{1}{\eta} d[Y_1;Y_2]$ for every $\eta<1/8$. Take the limit at $\eta=1/8$ to get $\rho(X_1)+\rho(X_2)\le \rho(Y_1) + \rho(Y_2) + 8 d[Y_1;Y_2]$.
By \Cref{ruzsa-triangle} and \Cref{ruzsa-nonneg} we have $d[X_1;X_1]=d[X_2;X_2]=0$, and by \Cref{sym-zero} there are $H_1:=\mathrm{Sym}[X_1],H_2:=\mathrm{Sym}[X_2]$ such that $X_1=U_{H_1}+x_1$ and $X_2=U_{H_2}+x_2$ for some $x_2$.
By \Cref{rho-invariant} we get $\rho(U_{H_1})+\rho(U_{H_2})\le \rho(Y_1) + \rho(Y_2) + 8 d[Y_1;Y_2]$, and thus the claim holds for $H=H_1$ or $H=H_2$.
\end{proof}

\begin{corollary}\label{pfr-9-aux}\lean{better_PFR_conjecture_aux0}\leanok
If $|A+A| \leq K|A|$, then there exists a subgroup $H$ and $t\in G$ such that
$|A \cap (H+t)| \geq K^{-4} \sqrt{|A||H|}$, and $|H|/|A|\in[K^{-8},K^8]$.
\end{corollary}

\begin{proof}\leanok
  \uses{pfr-rho,rho-init,rho-subgroup}
  Apply \Cref{pfr-rho} on $U_A,U_A$ to get a subspace such that $2\rho(U_H)\le 2\rho(U_A)+8d[U_A;U_A]$. Recall that $d[U_A;U_A]\le \log K$ as proved in \Cref{pfr_aux}, and $\rho(U_A)=0$ by \Cref{rho-init}. Therefore $\rho(U_H)\le 4\log(K)$. The claim then follows from \Cref{rho-subgroup}.
\end{proof}


\begin{corollary}\label{pfr-9-aux'}\lean{better_PFR_conjecture_aux}\leanok
If $|A+A| \leq K|A|$, then there exist a subgroup $H$ and a subset $c$ of $G$
with $A \subseteq c + H$, such that $|c| \leq K^{5} |A|^{1/2}/|H|^{1/2}$ and
$|H|/|A|\in[K^{-8},K^8]$.
\end{corollary}

\begin{proof}\leanok
  \uses{pfr-9-aux, ruz-cov}
Apply \Cref{pfr-9-aux} and \Cref{ruz-cov} to get the result, as in the proof
of \Cref{pfr_aux}.
\end{proof}


\begin{theorem}[PFR with \texorpdfstring{$C=9$}{C=9}]\label{pfr-9}\lean{better_PFR_conjecture}\leanok  If $A \subset {\bf F}_2^n$ is finite non-empty with $|A+A| \leq K|A|$, then there exists a subgroup $H$ of ${\bf F}_2^n$ with $|H| \leq |A|$ such that $A$ can be covered by at most $2K^9$ translates of $H$.
\end{theorem}

\begin{proof}\leanok
  \uses{pfr-9-aux',ruz-cov}
Given \Cref{pfr-9-aux'}, the proof is the same as that of \Cref{pfr}.
\end{proof}


%%FILE blueprint/src/chapter/hom_pfr.tex
\chapter{Homomorphism version of PFR}

In this section, $G, G'$ are finite abelian $2$-groups.

\begin{lemma}[Hahn-Banach type theorem]\label{hb-thm}\lean{hahn_banach}\leanok  Let $H_0$ be a subgroup of $G$.  Then every homomorphism $\phi: H_0 \to G'$ can be extended to a homomorphism $\tilde \phi: G \to G'$.
\end{lemma}

\begin{proof}\leanok  By induction it suffices to treat the case where $H_0$ has index $2$ in $G$, but then the extension can be constructed by hand.
\end{proof}

\begin{lemma}[Goursat type theorem]\label{goursat}\lean{goursat}\leanok  Let $H$ be a subgroup of $G \times G'$.  Then there exists a subgroup $H_0$ of $G$, a subgroup $H_1$ of $G'$, and a homomorphism $\phi: G \to G'$ such that
$$ H := \{ (x, \phi(x) + y): x \in H_0, y \in H_1 \}.$$
In particular, $|H| = |H_0| |H_1|$.
\end{lemma}

\begin{proof}\uses{hb-thm}\leanok We can take $H_0$ to be the projection of $H$ to $G$, and $H_1$ to be the slice $H_1 := \{ y: (0,y) \in H \}$.  One can construct $\phi$ on $H_0$ one generator at a time by the greedy algorithm, and then extend to $G$ by \Cref{hb-thm}.  The cardinality bound is clear from direct counting.
\end{proof}

\begin{theorem}[Homomorphism form of PFR]\label{hom-pfr}\lean{homomorphism_pfr}\leanok  Let $f: G \to G'$ be a function, and let $S$ denote the set
$$ S := \{ f(x+y)-f(x)-f(y): x,y \in G \}.$$
Then there exists a homomorphism $\phi: G \to G'$ such that
$$ |\{ f(x) - \phi(x): x \in G \}| \leq |S|^{10}.$$
\end{theorem}

\begin{proof}\uses{goursat, pfr-9-aux'}\leanok
Consider the graph $A \subset G \times G'$ defined by
$$ A := \{ (x,f(x)): x \in G \}.$$
Clearly, $|A| = |G|$.  By hypothesis, we have
$$ A+A \subset \{ (x,f(x)+s): x \in G, s \in S\}$$
and hence $|A+A| \leq |S| |A|$.  Applying \Cref{pfr-9-aux'}, we may find a
subspace $H \subset G \times G'$ such that $|H|/ |A| \in [|S|^{-8}, |S|^{8}]$
and $A$ is covered by $c + H$ with $|c| \le |S|^5|A|^{1/2} / |H|^{1/2}$. If
we let $H_0, H_1$ be as in \Cref{goursat}, this implies on taking projections
that $G$ is covered by at most $|c|$ translates of $H_0$. This implies that
$$ |c| |H_0| \geq |G|;$$
since $|H_0| |H_1| = |H|$, we conclude that
$$ |H_1| \leq |c| |H|/|G| = |c| |H|/|A|.$$
By hypothesis, $A$ is covered by at most $|c|$ translates of $H$, and hence
by at most $|c| |H_1|$ translates of $\{ (x,\phi(x)): x \in G \}$. As $\phi$
is a homomorphism, each such translate can be written in the form $\{
(x,\phi(x)+d): x \in G \}$ for some $d \in G'$. Since
$$
|c| |H_1| \le |c|^2 \frac{|H|}{|A|} \le \left(|S|^5 \frac{|A|^{1/2}}{|H|^{1/2}}\right)^2 \frac{|H|}{|A|}
= |S|^{10},
$$
the result follows.
\end{proof}


%%FILE blueprint/src/chapter/improved_exponent.tex
\chapter{Improving the exponents}

The arguments here are due to Jyun-Jie Liao.

\begin{definition}[New definition of $\eta$]\label{eta-def-new}\leanok  $\eta$ is a real parameter with $\eta > 0$.
\end{definition}

Previously in \Cref{eta-def} we had set $\eta=1/9$.  To implement
this chapter, one should refactor the previous arguments so that $\eta$ is
now free to be a positive number, though the specific hypothesis $\eta=1/9$
would now need to be added to \Cref{de-prop}.

Let $X^0_1, X^0_2$ be $G$-valued random variables, and let $X_1, X_2$ be $\tau$-minimizers as defined in \Cref{tau-min-def}.

For the next two lemmas, let $(T_1,T_2,T_3)$ be a $G^3$-valued random variable such that $T_1+T_2+T_3=0$ holds identically. Let $\delta$ be the quantity in \eqref{delta-t1t2t3-def}.

We have the following variant of \Cref{construct-good-prelim}:

\begin{lemma}[Constructing good variables, I']\label{construct-good-prelim-improv}\lean{construct_good_prelim'}\leanok
  One has
  \begin{align*}  k \leq
    \delta + \eta (& d[X^0_1;T_1|T_3]-d[X^0_1;X_1])
      + \eta (d[X^0_2;T_2|T_3]-d[X^0_2;X_2]).
  \end{align*}
\end{lemma}

\begin{proof} \uses{entropic-bsg,distance-lower}\leanok
  We apply \Cref{entropic-bsg} with $(A,B) = (T_1, T_2)$ there.
  Since $T_1 + T_2 = T_3$, the conclusion is that
  \begin{align} \nonumber \sum_{t_3} \bbP[T_3 = t_3] & d[(T_1 | T_3 = t_3); (T_2 | T_3 = t_3)] \\ & \leq 3 \bbI[T_1 : T_2] + 2 \bbH[T_3] - \bbH[T_1] - \bbH[T_2].\label{bsg-t1t2'}
  \end{align}
  The right-hand side in~\eqref{bsg-t1t2'} can be rearranged as
  \begin{align*} & 2( \bbH[T_1] + \bbH[T_2] + \bbH[T_3]) - 3 \bbH[T_1,T_2] \\ & = 2(\bbH[T_1] + \bbH[T_2] + \bbH[T_3]) - \bbH[T_1,T_2] - \bbH[T_2,T_3] - \bbH[T_1, T_3] = \delta,\end{align*}
  using the fact (from \Cref{relabeled-entropy}) that all three terms $\bbH[T_i,T_j]$ are equal to $\bbH[T_1,T_2,T_3]$ and hence to each other.
  We also have
  \begin{align*}
  &  \sum_{t_3} P[T_3 = t_3] \bigl(d[X^0_1; (T_1 | T_3=t_3)] - d[X^0_1;X_1]\bigr) \\
  &\quad = d[X^0_1; T_1 | T_3] - d[X^0_1;X_1]
  \end{align*}
  and similarly
  \begin{align*}
  &  \sum_{t_3} \bbP[T_3 = t_3] (d[X^0_2;(T_2 | T_3=t_3)] - d[X^0_2; X_2]) \\
  &\quad\quad\quad\quad\quad\quad \leq d[X^0_2;T_2|T_3] - d[X^0_2;X_2].
  \end{align*}
  Putting the above observations together, we have
  \begin{align*}
   \sum_{t_3} \bbP[T_3=t_3] \psi[(T_1 | T_3=t_3); (T_2 | T_3=t_3)] \leq \delta + \eta (d[X^0_1;T_1|T_3]-d[X^0_1;X_1]) \\
     + \eta (d[X^0_2;T_2|T_3]-d[X^0_2;X_2])
   \end{align*}
where we introduce the notation
\[\psi[Y_1; Y_2] := d[Y_1;Y_2] +  \eta (d[X_1^0;Y_1] - d[X_1^0;X_1]) + \eta(d[X_2^0;Y_2] - d[X_2^0;X_2]).\]
On the other hand, from \Cref{distance-lower} we have $k \leq \psi[Y_1;Y_2]$, and the claim follows.
\end{proof}

(One could in fact refactor \Cref{construct-good-prelim} to follow from \Cref{construct-good-prelim-improv} and \Cref{cond-dist-fact}).

\begin{lemma}[Constructing good variables, II']\label{construct-good-improv}\lean{construct_good_improved'}\leanok
One has
\begin{align*}  k & \leq  \delta + \frac{\eta}{6}  \sum_{i=1}^2 \sum_{1 \leq j,l \leq 3; j \neq l} (d[X^0_i;T_j|T_l] - d[X^0_i; X_i])
  \end{align*}
\end{lemma}

\begin{proof}
\uses{construct-good-prelim-improv}\leanok
Average \Cref{construct-good-prelim-improv} over all six permutations of $T_1,T_2,T_3$.
\end{proof}

Now let $X_1, X_2, \tilde X_1, \tilde X_2$ be independent copies of $X_1, X_2, X_1, X_2$, and set

\[U := X_1 + X_2, \qquad V := \tilde X_1 + X_2, \qquad W := X_1 + \tilde X_1\] and
\[  S := X_1 + X_2 + \tilde X_1 + \tilde X_2\]
and introduce the quantities
$$ k = d[X_1;X_2]$$
and
$$ I_1 = I(U : V \, | \, S).$$

\begin{lemma}[Constructing good variables, III']\label{averaged-construct-good}\lean{averaged_construct_good}\leanok
  One has
  \begin{align*}  k & \leq I(U : V \, | \, S) + I(V : W \, | \,S) + I(W : U \, | \, S) + \frac{\eta}{6}  \sum_{i=1}^2 \sum_{A,B \in \{U,V,W\}: A \neq B} (d[X^0_i;A|B,S] - d[X^0_i; X_i]).
    \end{align*}
\end{lemma}

\begin{proof}\uses{construct-good-improv, key-ident}\leanok  For each $s$ in the range of $S$, apply \Cref{construct-good-improv} with $T_1,T_2,T_3$ equal to $(U|S=s)$, $(V|S=s)$, $(W|S=s)$ respectively (which works thanks to \Cref{key-ident}), multiply by $\bbP[S=s]$, and sum in $s$ to conclude.
\end{proof}

To control the expressions in the right-hand side of this lemma we need a general entropy inequality.

\begin{lemma}[General inequality]\label{gen-ineq}\lean{gen_ineq_00, gen_ineq_01, gen_ineq_10} \leanok
Let $X_1, X_2, X_3, X_4$ be independent $G$-valued random variables, and let $Y$ be another $G$-valued random variable.  Set $S := X_1+X_2+X_3+X_4$. Then
  \begin{align*}
    &  d[Y; X_1+X_2|X_1 + X_3, S] - d[Y; X_1] \\
    &\quad \leq \tfrac{1}{4} (d[X_1;X_2] + 2d[X_1;X_3] + d[X_2;X_4])\\
    &\qquad \qquad + \tfrac{1}{4} (d[X_1|X_1+X_3;X_2|X_2+X_4] - d[X_3|X_3+X_4; X_1|X_1+X_2])\\
    &\qquad \qquad + \tfrac{1}{8} (\bbH[X_1+X_2] - \bbH[X_3+X_4] + \bbH[X_2] - \bbH[X_3]\\
    &\qquad \qquad \qquad + \bbH[X_2|X_2+X_4] - \bbH[X_1|X_1+X_3]).
\end{align*}
\end{lemma}

\begin{proof}\uses{cond-dist-fact, first-useful, cor-fibre}\leanok
  On the one hand, by \Cref{cond-dist-fact} and two applications of \Cref{first-useful} we have
  \begin{align*}
  &d[Y;X_1+X_2|X_1 + X_3, S] \\
  &\quad \leq d[Y;X_1+X_2|S] + \tfrac{1}{2} \bbI[X_1 + X_2 : X_1 + X_3|S] \\
  &\quad \leq d[Y;X_1+X_2]\\
  &\qquad + \tfrac{1}{2} (d[X_1+X_2;X_3+X_4] + \bbI[X_1 + X_2 : X_1 + X_3|S])\\
  &\qquad + \tfrac{1}{4} (\bbH[X_1+X_2] - \bbH[X_3+X_4])\\
  &\quad \leq d[Y;X_1] \\
  &\qquad + \tfrac{1}{2} (d[X_1;X_2] + d[X_1+X_2;X_3+X_4] + \bbI[X_1 + X_2 : X_1 + X_3|S])\\
  &\qquad + \tfrac{1}{4} (\bbH[X_1+X_2] - \bbH[X_3+X_4] + \bbH[X_2] - \bbH[X_1]).
  \end{align*}
  From \Cref{cor-fibre} (with $Y_1,Y_2,Y_3,Y_4$ set equal to $X_3, X_1, X_4, X_2$ respectively) one has
  $$ d[X_3+X_4; X_1+X_2] + d[X_3|X_3+X_4; X_1|X_1+X_2] $$
  $$ + \bbI[X_3 + X_1 : X_1 + X_2|S] = d[X_3;X_1] + d[X_4;X_2].$$
Rearranging the mutual information and Ruzsa distances slightly, we conclude that
\begin{align*}
  &d[Y;X_1+X_2|X_1 + X_3, S] \\
  &\quad \leq d[Y;X_1] \\
  &\qquad + \tfrac{1}{2} (d[X_1;X_2] + d[X_1;X_3] + d[X_2;X_4] - d[X_3|X_3+X_4; X_1|X_1+X_2])\\
  &\qquad + \tfrac{1}{4} (\bbH[X_1+X_2] - \bbH[X_3+X_4] + \bbH[X_2] - \bbH[X_1]).
  \end{align*}



  On the other hand, $(X_1+X_2|X_1 + X_3, S)$ has an identical distribution to the independent sum of $(X_1|X_1+X_3)$ and $(X_2|X_2+X_4)$.  We may therefore apply \Cref{first-useful} to conditioned variables $(X_1|X_1+X_3=s)$ and $(X_2|X_2+X_4=t)$ and average in $s,t$ to
  obtain the alternative bound
  \begin{align*}
  &  d[Y;X_1+X_2|X_1 + X_3, S] \\
  &\quad \leq d[Y;X_1|X_1+X_3] + \tfrac{1}{2} d[X_1|X_1+X_3; X_2|X_2+X_4] \\
  &\qquad + \tfrac{1}{4} (\bbH[X_2|X_2+X_4] - \bbH[X_1|X_1+X_3]) \\
  &\quad  \leq d[Y;X_1] \\
  &\qquad + \tfrac{1}{2} (d[X_1;X_3] + d[X_1|X_1+X_3;X_2|X_2+X_4])\\
  &\qquad + \tfrac{1}{4} (\bbH[X_2|X_2+X_4] - \bbH[X_1|X_1+X_3] + \bbH[X_1] - \bbH[X_3]).
  \end{align*}
  If one takes the arithmetic mean of these two bounds and simplifies using \Cref{cor-fibre}, one obtains the claim.
\end{proof}

Returning to our specific situation, we now have

\begin{lemma}[Bound on distance differences]\label{dist-diff-bound}\lean{dist_diff_bound_1, dist_diff_bound_2}\leanok We have
\begin{align*}  &\sum_{i=1}^2 \sum_{A,B \in \{U,V,W\}: A \neq B}  d[X_i^0;A|B, S] - d[X_i^0;X_i]\\
&\qquad \leq 12 k + \frac{4(2 \eta k - I_1)}{1-\eta}.
\end{align*}
\end{lemma}

\begin{proof}\uses{gen-ineq, relabeled-entropy-cond,second-estimate-aux}\leanok
If we apply \Cref{gen-ineq} with $X_1:=X_1$, $Y:=X_1^0$ and $(X_2,X_3,X_4)$ equal to the $3!$ permutations of $(X_2,\tilde X_1,\tilde X_2)$, and sums (using the symmetry $\bbH[X|X+Y] = \bbH[Y|X+Y]$, which follows from \Cref{relabeled-entropy-cond}), we can bound
$$ \sum_{A,B \in \{U,V,W\}: A \neq B}  d[X_1^0;A|B, S] - d[X_1^0;X_1]$$
by
\begin{align*}
  &\quad \tfrac{1}{4} (6d[X_1;X_2] + 6d[X_1;\tilde X_2]\\
  &\qquad + 6d[X_1;\tilde X_1] + 2d[\tilde X_1;\tilde X_2] + 2 d[\tilde X_1;X_2] + 2d[X_2;\tilde X_2])\\
  &\quad + \tfrac{1}{8} (2\bbH[X_1+X_2] + 2\bbH[X_1+\tilde X_1] + 2 \bbH[X_1+\tilde X_2] \\
  &\qquad - 2\bbH[\tilde X_1+X_2] - 2\bbH[X_2+\tilde X_2] - 2\bbH[\tilde X_1+\tilde X_2])\\
  &\qquad \qquad + \tfrac{1}{4} (\bbH[X_2|X_2+\tilde X_2] + \bbH[\tilde X_1|\tilde X_1+\tilde X_2] + \bbH[\tilde X_1|X_1+\tilde X_2] \\
  &\qquad \qquad \qquad - \bbH[X_1|X_1+\tilde X_1] - \bbH[X_1|X_1+X_2] - \bbH[X_1|X_1+\tilde X_2]),
\end{align*}
which simplifies to
\begin{align*}
  &\quad \tfrac{1}{4} (16k + 6d[X_1;X_1] + 2d[X_2;X_2])\\
  &\qquad \qquad + \tfrac{1}{4} (H[X_1+\tilde X_1] - H[X_2+\tilde X_2] + d[X_2|X_2+\tilde X_2] - d[X_1|X_1+\tilde X_1]).
\end{align*}
A symmetric argument also bounds
$$ \sum_{A,B \in \{U,V,W\}: A \neq B}  d[X_2^0;A|B, S] - d[X_2^0;X_2]$$
by
\begin{align*}
  &\quad \tfrac{1}{4} (16k + 6d[X_2;X_2] + 2d[X_1;X_1])\\
  &\qquad \qquad + \tfrac{1}{4} (H[X_2+\tilde X_2] - H[X_1+\tilde X_1] + d[X_1|X_1+\tilde X_1] - d[X_2|X_2+\tilde X_2]).
\end{align*}
On the other hand, from \Cref{second-estimate-aux} one has
$$ d[X_1;X_1] + d[X_2;X_2] \leq 2 k + \frac{2(2 \eta k - I_1)}{1-\eta}.$$
Summing the previous three estimates, we obtain the claim.
\end{proof}

\begin{theorem}[Improved $\tau$-decrement]\label{de-prop-improv}\lean{tau_strictly_decreases'}\leanok
Suppose $0 < \eta < 1/8$.  Let $X_1, X_2$ be tau-minimizers.  Then $d[X_1;X_2] = 0$.
\end{theorem}

\begin{proof}\uses{averaged-construct-good, dist-diff-bound, uvw-s}\leanok
From \Cref{averaged-construct-good}, \Cref{dist-diff-bound}, and \Cref{uvw-s} one has
  \[ k \leq 8\eta k - \frac{(1 -5 \eta - \frac{4}{6} \eta)(2 \eta k - I_1)}{(1-\eta)}.\]
  For any $\eta < 1/8$, we see from \Cref{first-estimate} that the expression $\frac{(1 -5 \eta - \frac{4}{6} \eta)(2 \eta k - I_1)}{(1-\eta)}$ is nonnegative, and hence $k = 0$ as required.
\end{proof}


\begin{theorem}[Limiting improved $\tau$-decrement]\label{de-prop-lim-improv}\lean{tau_minimizer_exists_rdist_eq_zero}\leanok
For $\eta = 1/8$, there exist tau-minimizers $X_1, X_2$ satisfying $d[X_1;X_2] = 0$.
\end{theorem}
\begin{proof}\uses{de-prop-improv, tau-min}\leanok
For each $\eta<1/8$, consider minimizers $X_1^\eta$ and $X_2^\eta$ from \Cref{tau-min}.
By \Cref{de-prop-improv},
they satisfy $d[X_1^\eta; X_2^\eta]=0$. By compactness of the space of probability measures
on $G$, one may extract a converging subsequence of the distributions of $X_1^\eta$ and $X_2^\eta$
as $\eta \to 1/8$. By continuity of all the involved quantities, the limit is a pair of
tau-minimizers for $1/8$ satisfying additionally $d[X_1;X_2] = 0$.
\end{proof}

\begin{theorem}[Improved entropy version of PFR]\label{entropy-pfr-improv}\lean{entropic_PFR_conjecture_improv}\leanok
  Let $G = \F_2^n$, and suppose that $X^0_1, X^0_2$ are $G$-valued random variables.
  Then there is some subgroup $H \leq G$ such that
  \[
    d[X^0_1;U_H] + d[X^0_2;U_H] \le 10 d[X^0_1;X^0_2],
  \]
  where $U_H$ is uniformly distributed on $H$.
  Furthermore, both $d[X^0_1;U_H]$ and $d[X^0_2;U_H]$ are at most $\tfrac{11}{2} d[X^0_1;X^0_2]$.
\end{theorem}

\begin{proof} \uses{de-prop-lim-improv, lem:100pc, ruzsa-triangle}\leanok
Let $X_1, X_2$ be the good $\tau$-minimizer from \Cref{de-prop-lim-improv}. By construction, $d[X_1;X_2]=0$.
From \Cref{lem:100pc}, $d[X_1;U_H] = d[X_2; U_H] = 0$.  Also from $\tau$-minimization we have $\tau[X_1;X_2] \leq \tau[X^0_2;X^0_1]$.  Using this and the Ruzsa triangle inequality we can conclude.
\end{proof}

One can then replace \Cref{pfr_aux} with

\begin{lemma}\label{pfr_aux-improv}\lean{PFR_conjecture_improv_aux}\leanok
If $A \subset {\bf F}_2^n$ is non-empty and
$|A+A| \leq K|A|$, then $A$ can be covered by at most $K^6 |A|^{1/2}/|H|^{1/2}$ translates of a subspace $H$ of ${\bf F}_2^n$ with
$$
  |H|/|A| \in [K^{-10}, K^{10}].
$$
\end{lemma}

\begin{proof}\uses{entropy-pfr-improv, unif-exist, uniform-entropy-II, jensen-bound,ruz-dist-def,ruzsa-diff,bound-conc,ruz-cov}\leanok
By repeating the proof of \Cref{pfr_aux} and using \Cref{entropy-pfr-improv} one can obtain the claim with $13/2$
replaced with $6$ and $11$ replaced by $10$.
\end{proof}

This implies the following improved version of Theorem~\ref{pfr}:

\begin{theorem}[Improved PFR]\label{pfr-improv}\lean{PFR_conjecture_improv}\leanok
  If $A \subset {\bf F}_2^n$ is non-empty and $|A+A| \leq K|A|$, then $A$ can be covered by most $2K^{11}$ translates of a subspace $H$ of ${\bf F}_2^n$ with $|H| \leq |A|$.
\end{theorem}

\begin{proof}\uses{pfr_aux-improv}\leanok
By repeating the proof of \Cref{pfr} and using \Cref{pfr_aux-improv} one obtains $2K^{11}$ in place of $2K^{12}$.
\end{proof}

Of course, by replacing \Cref{pfr} with \Cref{pfr-improv} we may also improve constants in downstream theorems in a straightforward manner.


%%FILE blueprint/src/chapter/jensen.tex
\chapter{Applications of Jensen's inequality}

In this chapter, $h$ denotes the function $h(x) := x \log \frac{1}{x}$ for $x \in [0,1]$.

\begin{lemma}[Concavity]\label{concave}
  \lean{Real.strictConcaveOn_negMulLog}\leanok
  $h$ is strictly concave on $[0,\infty)$.
\end{lemma}

\begin{proof} \leanok Check that $h'$ is strictly monotone decreasing.
\end{proof}


\begin{lemma}[log sum inequality]
  \label{log-sum}\lean{Real.sum_mul_log_div_leq}\leanok
  If $S$ is a finite set, and $a_s,b_s$ are non-negative for $s\in S$, then
  $$\sum_{s\in S} a_s \log\frac{a_s}{b_s}\ge \left(\sum_{s\in S}a_s\right)\log\frac{\sum_{s\in S} a_s}{\sum_{s\in S} b_s},$$
  with the convention $0\log\frac{0}{b}=0$ for any $b\ge 0$ and $0\log\frac{a}{0}=\infty$ for any $a>0$.
\end{lemma}
\begin{proof}\leanok
  Let $B:=\sum_{s\in S} b_s$. Apply Jensen and \Cref{concave} to show that $\sum_{s\in S} \frac{b_s}{B} h(\frac{a_s}{b_s})\ge h(\frac{\sum_{s\in S} a_s}{B})$.
\end{proof}

\begin{lemma}[converse log sum]
  \label{converse-log-sum}\lean{Real.sum_mul_log_div_eq_iff}\leanok
If equality holds in \Cref{log-sum}, then $a_s=r\cdot b_s$ for every $s\in S$, for some constant $r\in \mathbb{R}$.
\end{lemma}

\begin{proof}\leanok
\uses{concave} By the fact that $h$ is strictly concave and the equality condition of Jensen.
\end{proof}


%%FILE blueprint/src/chapter/main.tex
% This is the main point of entry to the blueprint.
% Add chapters of the blueprint here.
% This file is not meant to be built. Build src/web.tex or src/print.text instead.

\input{chapter/jensen}
\input{chapter/entropy}
\input{chapter/distance}
\input{chapter/100_percent}
\input{chapter/fibring}
\input{chapter/entropy_pfr}
\input{chapter/pfr}
\input{chapter/improved_exponent}
\input{chapter/hom_pfr}
\input{chapter/approx_hom_pfr}
\input{chapter/weak_pfr}
\input{chapter/torsion}
\input{chapter/further_improvement}


%%FILE blueprint/src/chapter/pfr.tex
\chapter{Proof of PFR}

\begin{lemma}[Ruzsa covering lemma]
\label{ruz-cov}
\lean{Finset.ruzsa_covering_add}
\leanok
If $A,B$ are finite non-empty subsets of a group $G$, then $A$ can be covered by at most $|A+B|/|B|$ translates of $B-B$.
\end{lemma}
\begin{proof}
\leanok
Cover $A$ greedily by disjoint translates of $B$.
\end{proof}

\begin{lemma}\label{pfr_aux}
  \lean{PFR_conjecture_aux}\leanok If $A \subset {\bf F}_2^n$ is non-empty and
  $|A+A| \leq K|A|$, then $A$ can be covered by at most $K ^
  {13/2}|A|^{1/2}/|H|^{1/2}$ translates of a subspace $H$ of ${\bf F}_2^n$ with
  \begin{equation}
    \label{ah}
    |H|/|A| \in [K^{-11}, K^{11}].
  \end{equation}
  \end{lemma}

\begin{proof}
\uses{entropy-pfr, unif-exist, uniform-entropy-II, jensen-bound,ruz-dist-def,ruzsa-diff,bound-conc,ruz-cov}\leanok
  Let $U_A$ be the uniform distribution on $A$ (which exists by \Cref{unif-exist}), thus $\bbH[U_A] = \log |A|$ by \Cref{uniform-entropy-II}. By \Cref{jensen-bound} and the fact that $U_A + U_A$ is supported on $A + A$, $\bbH[U_A + U_A] \leq \log|A+A|$. By \Cref{ruz-dist-def}, the doubling condition $|A+A| \leq K|A|$ therefore gives
  \[d[U_A;U_A] \leq \log K.\]
  By \Cref{entropy-pfr}, we may thus find a subspace $H$ of $\F_2^n$ such that
  \begin{equation}\label{uauh} d[U_A;U_H] \leq \tfrac{1}{2} C' \log K\end{equation}
  with $C' = 11$.
  By \Cref{ruzsa-diff} we conclude that
  \begin{equation*}
    |\log |H| - \log |A|| \leq C' \log K,
  \end{equation*}
  proving~\eqref{ah}.
  From \Cref{ruz-dist-def},~\eqref{uauh} is equivalent to
  \[\bbH[U_A - U_H] \leq \log( |A|^{1/2} |H|^{1/2}) + \tfrac{1}{2} C' \log K.\]
  By \Cref{bound-conc} we conclude the existence of a point $x_0 \in \F_p^n$ such that
  \[p_{U_A-U_H}(x_0) \geq |A|^{-1/2} |H|^{-1/2} K^{-C'/2},\]
  or equivalently
  \[|A \cap (H + x_0)| \geq K^{-C'/2} |A|^{1/2} |H|^{1/2}.\]
  Applying \Cref{ruz-cov}, we may thus cover $A$ by at most
  \[\frac{|A + (A \cap (H+x_0))|}{|A \cap (H + x_0)|} \leq \frac{K|A|}{K^{-C'/2} |A|^{1/2} |H|^{1/2}} = K^{C'/2+1} \frac{|A|^{1/2}}{|H|^{1/2}}\]
  translates of
  \[\bigl(A \cap (H + x_0)\bigr) - \bigl(A \cap (H + x_0)\bigr) \subseteq H.\]
  This proves the claim.
\end{proof}


\begin{theorem}[PFR]\label{pfr}
\lean{PFR_conjecture}\leanok
If $A \subset {\bf F}_2^n$ is non-empty and $|A+A| \leq K|A|$, then $A$ can be covered by most $2K^{12}$ translates of a subspace $H$ of ${\bf F}_2^n$ with $|H| \leq |A|$.
\end{theorem}

\begin{proof}
  \uses{pfr_aux}\leanok
  Let $H$ be given by \Cref{pfr_aux}.
  If $|H| \leq |A|$ then we are already done thanks to~\eqref{ah}.  If $|H| > |A|$ then we can cover $H$ by at most $2 |H|/|A|$ translates of a subspace $H'$ of $H$ with $|H'| \leq |A|$.  We can thus cover $A$ by at most
  \[2K^{13/2} \frac{|H|^{1/2}}{|A|^{1/2}}\]
  translates of $H'$, and the claim again follows from~\eqref{ah}.
\end{proof}

\begin{corollary}[PFR in infinite groups]\label{pfr-cor}
  \lean{PFR_conjecture}\leanok
  If $G$ is an abelian $2$-torsion group, $A \subset G$ is non-empty finite, and $|A+A| \leq K|A|
  $, then $A$ can be covered by most $2K^{12}$ translates of a finite group  $H$ of $G$ with $|H| \leq |A|$.
  \end{corollary}

\begin{proof}\uses{pfr}\leanok  Apply \Cref{pfr} to the group generated by $A$, which is isomorphic to $\F_2^n$ for some $n$.
\end{proof}


%%FILE blueprint/src/chapter/torsion.tex
\chapter{The \texorpdfstring{$m$}{m}-torsion case}

\section{Data processing inequality}

\begin{lemma}[Data processing for a single variable]\label{data-process-single}\lean{ProbabilityTheory.entropy_comp_le}\leanok  Let $X$ be a random variable.  Then for any function $f$ on the range of $X$, one has $\bbH[f(X)] \leq \bbH[X]$.
\end{lemma}

\begin{proof}\uses{relabeled-entropy, chain-rule}\leanok
We have
$$ \bbH[X] = \bbH[X,f(X)] = \bbH[f(X)] + \bbH[X|f(X)]$$
thanks to \Cref{relabeled-entropy} and \Cref{chain-rule}, giving the claim.
\end{proof}

\begin{lemma}[One-sided unconditional data processing inequality]\label{data-process-unc-one}\lean{ProbabilityTheory.mutual_comp_le}\leanok Let $X,Y$ be random variables. For any function $f, g$ on the range of $X$, we have $\bbI[f(X) : Y] \leq \bbI[X:Y]$.
\end{lemma}

\begin{proof}\uses{alternative-mutual, submodularity, relabeled-entropy}\leanok
  By \Cref{alternative-mutual} it suffices to show that $\bbH[Y|X] \leq \bbH[Y|f(X)]$. But this follows from \Cref{submodularity} (and \Cref{relabeled-entropy}).
\end{proof}

\begin{lemma}[Unconditional data processing inequality]\label{data-process-unc}\lean{ProbabilityTheory.mutual_comp_comp_le}\leanok Let $X,Y$ be random variables. For any functions $f, g$ on the ranges of $X, Y$ respectively, we have $\bbI[f(X) : g(Y )] \leq \bbI[X : Y]$.
\end{lemma}

\begin{proof}\uses{data-process-unc-one, entropy-comm}\leanok
   From \Cref{data-process-unc-one}, \Cref{entropy-comm} we have $\bbI[f(X) : Y] \leq \bbI[X:Y]$ and $\bbI[f(X): g(Y)] \leq \bbI[f(X):Y]$, giving the claim.
\end{proof}

\begin{lemma}[Data processing inequality]\label{data-process}\lean{ProbabilityTheory.condMutual_comp_comp_le}\leanok Let $X,Y,Z$. For any functions $f, g$
on the ranges of $X, Y$ respectively, we have $\bbI[f(X) : g(Y )|Z] \leq \bbI[X :Y |Z]$.
\end{lemma}

\begin{proof}\uses{data-process-unc, conditional-mutual-def}\leanok  Apply \Cref{data-process-unc} to $X,Y$ conditioned to the event $Z=z$, multiply by ${\bf P}[Z=z]$, and sum using \Cref{conditional-mutual-def}.
\end{proof}

\section{More Ruzsa distance estimates}

Let $G$ be an additive group.

\begin{lemma}[Flipping a sign]\label{sign-flip}\lean{rdist_of_neg_le}\leanok
  If $X,Y$ are $G$-valued, then
  $$  d[X ; -Y]  \leq 3 d[X;Y].$$
\end{lemma}

\begin{proof}\uses{ruz-copy, independent-exist, cond-indep-exist, alt-submodularity, ruz-indep, neg-ent, relabeled-entropy, add-entropy, subadditive, data-process-single}\leanok
Without loss of generality (using \Cref{ruz-copy} and \Cref{independent-exist}) we may take $X,Y$ to be independent.
By $(X_1,Y_1)$, $(X_2,Y_2)$ be copies of $(X,Y)$ that are conditionally independent over $X_1-Y_1=X_2-Y_2$ (this exists thanks to \Cref{cond-indep-exist}).  By \Cref{independent-exist}, we can also find another copy $(X_3,Y_3)$ of $(X,Y)$ that is independent of $X_1,Y_1,X_2,Y_2$.  From \Cref{alt-submodularity}, one has
$$ \bbH[X_3-Y_2, X_1-Y_3, X_2, Y_1, X_3, Y_3, X_3+Y_3] + \bbH[X_3+Y_3] \leq \bbH[X_3-Y_2, X_1-Y_3, X_2, Y_1, X_3+Y_3] + \bbH[X_3, Y_3, X_3+Y_3].$$
From \Cref{ruz-indep}, \Cref{neg-ent}, \Cref{ruz-copy} we have
$$ \bbH[X_3+Y_3] = \frac{1}{2} \bbH[X_3] + \frac{1}{2} \bbH[-Y_3] + d[X_3;-Y_3] = \frac{1}{2} \bbH[X] + \frac{1}{2} \bbH[Y] + d[X;-Y].$$
Since $X_3+Y_3$ is a function of $X_3,Y_3$, we see from \Cref{relabeled-entropy} and \Cref{add-entropy} that
$$ \bbH[X_3,Y_3,X_3+Y_3] = \bbH[X_3,Y_3] = \bbH[X,Y] = \bbH[X]+\bbH[Y].$$
Because $X_1-Y_1=X_2-Y_2$, we have
$$ X_3+Y_3 = (X_3-Y_2) - (X_1-Y_3) + (X_2+Y_1)$$
and thus by \Cref{relabeled-entropy}
$$ \bbH[X_3-Y_2, X_1-Y_3, X_2, Y_1, X_3+Y_3] = \bbH[X_3-Y_2, X_1-Y_3, X_2, Y_1]$$
and hence by \Cref{subadditive}
$$ \bbH[X_3-Y_2, X_1-Y_3, X_2, Y_1, X_3+Y_3] \leq \bbH[X_3-Y_2] + \bbH[X_1-Y_3] + \bbH[X_2] + \bbH[Y_1].$$
Since $X_3,Y_2$ are independent, we see from \Cref{ruz-indep}, \Cref{ruz-copy} that
$$\bbH[X_3-Y_2] = \frac{1}{2} \bbH[X] + \frac{1}{2} \bbH[Y] + d[X; Y].$$
Similarly
$$ \bbH[X_1-Y_3] = \frac{1}{2} \bbH[X] + \frac{1}{2} \bbH[Y] + d[X; Y].$$
We conclude that
$$ \bbH[X_3-Y_2, X_1-Y_3, X_2, Y_1, X_3+Y_3] \leq 2\bbH[X] + 2\bbH[Y] + 2d[X; Y].$$
Finally, from \Cref{data-process-single} we have
$$ \bbH[X_1,Y_1,X_2,Y_2,X_3,Y_3] \leq \bbH[X_3-Y_2, X_1-Y_3, X_2, Y_1, X_3, Y_3, X_3+Y_3].$$
From \Cref{add-entropy} followed by \Cref{cond-trial-ent}, we have
$$\bbH[X_1,Y_1,X_2,Y_2,X_3,Y_3] = \bbH[X_1,Y_1,X_1-Y_1] + \bbH[X_2,Y_2,X_2-Y_2] - \bbH[X_1-Y_1] + \bbH[X_3,Y_3]$$
and thus by \Cref{ruz-indep}, \Cref{ruz-copy}, \Cref{relabeled-entropy}, \Cref{add-entropy}
$$\bbH[X_1,Y_1,X_2,Y_2,X_3,Y_3] = \bbH[X] + \bbH[Y] + \bbH[X] + \bbH[Y] -\left(\frac{1}{2}\bbH[X] + \frac{1}{2}\bbH[Y] + d[X; Y]\right) + \bbH[X] + \bbH[Y].$$
Applying all of these estimates, the claim now follows from linear arithmetic.
\end{proof}

\begin{lemma}[Kaimonovich--Vershik--Madiman inequality]\label{klm-1}\lean{kvm_ineq_I}\leanok  If $n \geq 0$ and $X, Y_1, \dots, Y_n$ are jointly independent $G$-valued random variables, then
  $$\bbH\left[X + \sum_{i=1}^n Y_i\right] - \bbH[X] \leq \sum_{i=1}^n \left(\bbH[X+Y_i] - \bbH[X]\right).$$
\end{lemma}

\begin{proof}\uses{kv}\leanok  This is trivial for $n=0,1$, while the $n=2$ case is \Cref{kv}.  Now suppose inductively that $n > 2$, and the claim was already proven for $n-1$.  By a further application of \Cref{kv} one has
$$  \bbH\left[X + \sum_{i=1}^n Y_i\right] -  \bbH\left[X + \sum_{i=1}^{n-1} Y_i\right] \leq \bbH[X+Y_n] - \bbH[X].$$
By induction hypothesis one has
$$ \bbH\left[X + \sum_{i=1}^{n-1} Y_i\right] - \bbH[X] \leq \sum_{i=1}^{n-1} \bbH[X+Y_i] - \bbH[X].$$
Summing the two inequalities, we obtain the claim.
\end{proof}

\begin{lemma}[Kaimonovich--Vershik--Madiman inequality, II]\label{klm-2}\lean{kvm_ineq_II}\leanok  If $n \geq 1$ and $X, Y_1, \dots, Y_n$ are jointly independent $G$-valued random variables, then
  $$ d[X; \sum_{i=1}^n Y_i] \leq 2 \sum_{i=1}^n d[X; Y_i].$$
\end{lemma}

\begin{proof}\uses{klm-1, neg-ent, ruz-indep, sumset-lower, ruzsa-diff}\leanok
  Applying \Cref{klm-1} with all the $Y_i$ replaced by $-Y_i$, and using \Cref{neg-ent} and \Cref{ruz-indep}, we obtain after some rearranging
$$ d[X; \sum_{i=1}^n Y_i] + \frac{1}{2} (\bbH[\sum_{i=1}^n Y_i] - \bbH[X]) \leq \sum_{i=1}^n \left(d[X;Y_i] + \frac{1}{2} (\bbH[Y_i] - \bbH[X])\right).$$
From \Cref{sumset-lower} we have
$$ \bbH[\sum_{i=1}^n Y_i] \geq \bbH[Y_i]$$
for all $i$; subtracting $\bbH[X]$ and averaging, we conclude that
$$ \bbH[\sum_{i=1}^n Y_i] - \bbH[X] \geq \frac{1}{n} \sum_{i=1}^n \bbH[Y_i] - \bbH[X]$$
and thus
$$ d[X; \sum_{i=1}^n Y_i] \leq \sum_{i=1}^n d[X;Y_i] + \frac{n-1}{2n} (\bbH[Y_i] - \bbH[X]).$$
From \Cref{ruzsa-diff} we have
$$ \bbH[Y_i] - \bbH[X] \leq 2 d[X;Y_i].$$
Since $0 \leq \frac{n-1}{2n} \leq \frac{1}{2}$, the claim follows.
\end{proof}

\begin{lemma}[Kaimonovich--Vershik--Madiman inequality, III]\label{klm-3}\lean{kvm_ineq_III}\leanok  If $n \geq 1$ and $X, Y_1, \dots, Y_n$ are jointly independent $G$-valued random variables, then
  $$d\left[X; \sum_{i=1}^n Y_i\right] \leq d\left[X; Y_1\right] + \frac{1}{2}\left(\bbH\left[ \sum_{i=1}^n Y_i\right] - \bbH[Y_1]\right).$$
\end{lemma}

\begin{proof}\uses{kv, ruz-indep}\leanok
  From \Cref{kv} one has
  $$ \bbH\left[-X + \sum_{i=1}^n Y_i\right] \leq \bbH[ - X + Y_1 ] + \bbH\left[ \sum_{i=1}^n Y_i \right] - \bbH[Y_1].$$
  The claim then follows from \Cref{ruz-indep} and some elementary algebra.
\end{proof}

\begin{lemma}[Comparing sums]\label{compare-sums}\lean{ent_of_sum_le_ent_of_sum}\leanok  Let $(X_i)_{1 \leq i \leq m}$ and $(Y_j)_{1 \leq j \leq l}$ be tuples of jointly independent random variables (so the $X$'s and $Y$'s are also independent of each other), and let $f: \{1,\dots,l\} \to \{1,\dots,m\}$ be a function, then
  $$ \bbH[\sum_{j=1}^l Y_j] \leq \bbH[ \sum_{i=1}^m X_i ] + \sum_{j=1}^l (\bbH[ Y_j - X_{f(j)}] - \bbH[X_{f(j)}]).$$
\end{lemma}

\begin{proof}\uses{klm-1, kv, neg-ent, sumset-lower}\leanok  Write $W := \sum_{i=1}^m X_i$.  From \Cref{sumset-lower} we have
$$ \bbH[\sum_{j=1}^l Y_j]  \leq \bbH[-W + \sum_{j=1}^l Y_j]$$
while from \Cref{klm-1} one has
$$ \bbH[-W + \sum_{j=1}^l Y_j] \leq \bbH[-W] + \sum_{j=1}^l \bbH[-W + Y_j] - \bbH[-W].$$
From \Cref{kv} one has
$$ \bbH[-W + Y_j] - \bbH[-W] \leq \bbH[-X_{f(j)} + Y_j] - \bbH[-X_{f(j)}].$$
The claim now follows from \Cref{neg-ent} and some elementary algebra.
\end{proof}

\begin{lemma}[Sums of dilates I]
  \label{sum-dilate-I}
  \lean{ent_sub_nsmul_le, ent_of_sub_smul'}
  \leanok

  Let $X,Y,X'$ be independent $G$-valued random variables, with $X'$ a copy of $X$,
  and let $a$ be an integer.  Then
  $$\bbH[X-(a+1)Y] \leq \bbH[X-aY] + \bbH[X-Y-X'] - \bbH[X]$$
  and
  $$\bbH[X-(a-1)Y] \leq \bbH[X-aY] + \bbH[X-Y-X'] - \bbH[X].$$
\end{lemma}
\begin{proof}
  \uses{ruzsa-triangle-improved, neg-ent}
  \leanok

  From \Cref{ruzsa-triangle-improved} we have
  $$ \bbH[(X-Y)-aY] \leq \bbH[(X-Y) - X'] + \bbH[X'-aY] - \bbH[X']$$
  which gives the first inequality.  Similarly from \Cref{ruzsa-triangle-improved} we have
  $$ \bbH[(X+Y)-aY] \leq \bbH[(X+Y) - X'] + \bbH[X'-aY] - \bbH[X']$$
  which (when combined with \Cref{neg-ent}) gives the second inequality.
\end{proof}

\begin{lemma}[Sums of dilates II]\label{sum-dilate-II}\lean{ent_sub_zsmul_sub_ent_le}\leanok  Let $X,Y$ be independent $G$-valued random variables, and let $a$ be an integer.  Then
  $$\bbH[X-aY] - \bbH[X] \leq 4 |a| d[X;Y].$$
\end{lemma}

\begin{proof}\uses{kv, ruz-indep, sign-flip, sum-dilate-I}\leanok From \Cref{kv} one has
  $$\bbH[Y-X+X'] - \bbH[Y-X] \leq \bbH[Y+X'] - \bbH[Y] = \bbH[Y+X] - \bbH[Y]$$
  which by \Cref{ruz-indep} gives
  $$\bbH[X-Y-X'] -\bbH[X] \leq d[X;Y] + d[X;-Y]$$
  and hence by \Cref{sign-flip}
  $$\bbH[X-Y-X'] - \bbH[X] \leq 4d[X;Y].$$
  From \Cref{sum-dilate-I} we then have
  $$\bbH[X-(a\pm 1)Y] \leq \bbH[X-aY] + 4 d[X;Y]$$
and the claim now follows by an induction on $|a|$.
\end{proof}

We remark that in the paper [GGMT2024] the variant estimate
$$\bbH[X-aY] - \bbH[X] \leq (4 + 10 \lfloor \log_2 |a| \rfloor) d[X;Y]$$
is also proven by a similar method.  This variant is superior for $|a| \geq 9$ (or $|a|=7$); but we will not need this estimate here.



\section{Multidistance}

We continue to let $G$ be an abelian group.

\begin{definition}[Multidistance]\label{multidist-def}\lean{multiDist}\leanok  Let $m$ be a positive integer, and let $X_{[m]} = (X_i)_{1 \leq i \leq m}$ be an $m$-tuple of $G$-valued random variables $X_i$. Then we define
\[
  D[X_{[m]}] := \bbH[\sum_{i=1}^m \tilde X_i] - \frac{1}{m} \sum_{i=1}^m \bbH[\tilde X_i],
\]
where the $\tilde X_i$ are independent copies of the $X_i$.
\end{definition}

\begin{lemma}[Multidistance of copy]\label{multidist-copy}\lean{multiDist_copy}\leanok  If $X_{[m]} = (X_i)_{1 \leq i \leq m}$ and $Y_{[m]} = (Y_i)_{1 \leq i \leq m}$ are such that $X_i$ and $Y_i$ have the same distribution for each $i$, then $D[X_{[m]}] = D[Y_{[m]}]$.
\end{lemma}

\begin{proof}\uses{multidist-def}\leanok Clear from Lemma \ref{copy-ent}.
\end{proof}

\begin{lemma}[Multidistance of independent variables]\label{multidist-indep}\lean{multiDist_indep}\leanok  If $X_{[m]} = (X_i)_{1 \leq i \leq m}$ are jointly independent, then $D[X_{[m]}] = \bbH[\sum_{i=1}^m X_i] - \frac{1}{m} \sum_{i=1}^m \bbH[X_i]$.
\end{lemma}

\begin{proof}\uses{multidist-def}\leanok  Clear from definition.
\end{proof}

\begin{lemma}[Nonnegativity]\label{multidist-nonneg}\lean{multiDist_nonneg}\uses{multidist-def}\leanok  For any such tuple, we have $D[X_{[m]}] \geq 0$.
\end{lemma}

\begin{proof}\uses{sumset-lower}\leanok  From \Cref{sumset-lower} one has
$$ \bbH[\sum_{i =1}^m \tilde X_i] \geq \bbH[\tilde X_i]$$
for each $1 \leq i \leq m$.  Averaging over $i$, we obtain the claim.
\end{proof}

\begin{lemma}[Relabeling]\label{multidist-perm}\lean{multiDist_of_perm}\uses{multidist-def}\leanok If $\phi: \{1,\dots,m\} \to \{1,\dots,m\}$ is a bijection, then $D[X_{[m]}] = D[(X_{\phi(j)})_{1 \leq j \leq m}]$.
\end{lemma}

\begin{proof}\leanok Trivial.
\end{proof}

\begin{lemma}[Multidistance and Ruzsa distance, I]\label{multidist-ruzsa-I}\lean{multidist_ruzsa_I}\uses{multidist-def}\leanok
  Let $m \ge 2$, and let $X_{[m]}$ be a tuple of $G$-valued random variables. Then
  $$\sum_{1 \leq j,k \leq m: j \neq k} d[X_j; -X_k] \leq m(m-1) D[X_{[m]}].$$
\end{lemma}

\begin{proof}\uses{ruz-copy, sumset-lower, ruz-indep, multidist-indep, multidist-copy}\leanok
By \Cref{ruz-copy}, \Cref{multidist-copy} we may take the $X_i$ to be jointly independent.  From \Cref{sumset-lower}, we see that for any distinct $1 \leq j,k \leq m$, we have
  \[
    \bbH[X_j+X_k] \leq \bbH[\sum_{i=1}^m X_i],
  \]
  and hence by \Cref{ruz-indep}
  \[
    d[X_j;-X_k] \leq \bbH[\sum_{i=1}^m X_i] - \frac{1}{2} \bbH[X_j] - \frac{1}{2} \bbH[X_k].
  \]
  Summing this over all pairs $(j,k)$, $j \neq k$ and using Lemma \ref{multidist-indep}, we obtain the claim.
\end{proof}

\begin{lemma}[Multidistance and Ruzsa distance, II]\label{multidist-ruzsa-II}\lean{multidist_ruzsa_II}\uses{multidist-def}\leanok
  Let $m \ge 2$, and let $X_{[m]}$ be a tuple of $G$-valued random variables. Then
  $$\sum_{j=1}^m d[X_j;X_j] \leq 2 m D[X_{[m]}].$$
\end{lemma}

\begin{proof}\uses{ruzsa-triangle,multidist-ruzsa-I}\leanok
From \Cref{ruzsa-triangle} we have $d[X_j;X_j] \leq 2 d[X_j;-X_k]$, and applying this to every summand in \Cref{multidist-ruzsa-I}, we obtain the claim.
\end{proof}

\begin{lemma}[Multidistance and Ruzsa distance, III]\label{multidist-ruzsa-III}\lean{multidist_ruzsa_III}\uses{multidist-def}\leanok
  Let $m \ge 2$, and let $X_{[m]}$ be a tuple of $G$-valued random variables. If the $X_i$ all have the same distribution, then $D[X_{[m]}] \leq m d[X_i;X_i]$ for any $1 \leq i \leq m$.
\end{lemma}

\begin{proof}\uses{klm-1, neg-ent, ruz-indep, sumset-lower, ruz-copy, multidist-copy, multidist-indep}\leanok
  By \Cref{ruz-copy}, \Cref{multidist-copy} we may take the $X_i$ to be jointly independent.    Let $X_0$ be a further independent copy of the $X_i$.
From \Cref{klm-1}, we have
$$ \bbH[-X_0 + \sum_{i=1}^m X_i] - \bbH[-X_0] \leq \sum_{i=1}^m \bbH[X_0 - X_i] - \bbH[-X_0]$$
and hence by \Cref{neg-ent} and \Cref{ruz-indep}
$$ \bbH[-X_0 + \sum_{i=1}^m X_i] - \bbH[X_0] \leq m d[X_i,X_i].$$
On the other hand, by \Cref{sumset-lower} we have
$$ \bbH[\sum_{i=1}^m X_i] \leq \bbH[-X_0 + \sum_{i=1}^m X_i]$$
and the claim follows.
\end{proof}

\begin{lemma}[Multidistance and Ruzsa distance, IV]\label{multidist-ruzsa-IV}\lean{multidist_ruzsa_IV}\uses{multidist-def}\leanok
  Let $m \ge 2$, and let $X_{[m]}$ be a tuple of independent $G$-valued random variables.  Let $W := \sum_{i=1}^m X_i$. Then
  $$ d[W;-W] \leq 2 D[X_i].$$
\end{lemma}

\begin{proof}\uses{independent-exist, kv, sumset-lower, ruz-indep, multidist-indep}\leanok
  Take $(X'_i)_{1 \leq i \leq m}$ to be further independent copies of $(X_i)_{1 \leq i \leq m}$ (which exist by \Cref{independent-exist}), and write $W' := \sum_{i=1}^m X'_i$.
  Fix any distinct $a,b \in I$.

  From \Cref{kv} one has
  \begin{equation}\label{7922}
    \bbH[W + W'] \leq  \bbH[W] + \bbH[X_{a} + W'] - \bbH[X_{a}]
   \end{equation}
   and also
   \[ \bbH[X_a + W'] \leq \bbH[X_a + X_b] + \bbH[W'] - \bbH[X'_b].\]
   Combining this with~\eqref{7922} and then applying \Cref{sumset-lower} we have
   \begin{align*}  \bbH[W + W']  & \leq    2\bbH[W] + \bbH[X_a + X_b]  - \bbH[X_a] - \bbH[X_b] \\ & \leq
    3 \bbH[W] - \bbH[X_a] - \bbH[X_b].
  \end{align*}
  Averaging this over all choices of $(a,b)$ gives $\bbH[W] + 2 D[X_{[m]}]$, and the claim follows from \Cref{ruz-indep}.
\end{proof}

\begin{proposition}[Vanishing]\label{multi-zero}\lean{multidist_eq_zero}\uses{multidist-def}\leanok
If $D[X_{[m]}]=0$, then for each $1 \leq i \leq m$ there is a finite subgroup $H_i \leq G$ such that $d[X_i; U_{H_i}] = 0$.
\end{proposition}

\begin{proof}\uses{multidist-ruzsa-I, ruzsa-nonneg, lem:100pc}\leanok  From \Cref{multidist-ruzsa-II} and \Cref{ruzsa-nonneg} we have $d[X_j; X_j]=0$ for all $1 \leq j \leq m$.  The claim now follows from \Cref{lem:100pc}.
\end{proof}

With more effort one can show that $H_i$ is independent of $i$, but we will not need to do so here.

\section{The tau functional}

Fix $m \geq 2$, and a reference variable $X^0$ in $G$.

\begin{definition}[$\eta$]\label{eta-def-multi}\lean{multiRefPackage}\leanok
We set $\eta := \frac{1}{32m^3}$.
\end{definition}

\begin{definition}[$\tau$-functional]\label{tau-def-multi}\uses{eta-def-multi}\lean{multiTau}\uses{multidist-def}\leanok  If $(X_i)_{1 \leq i \leq m}$ is a tuple, we define its $\tau$-functional
$$ \tau[ (X_i)_{1 \leq i \leq m}] := D[(X_i)_{1 \leq i \leq m}] + \eta \sum_{i=1}^m d[X_i; X^0].$$
\end{definition}

\begin{definition}[$\tau$-minimizer]
  \label{tau-min-multi}
  \uses{tau-def-multi}
  \lean{multiTauMinimizes}
  \leanok

  A $\tau$-minimizer is a tuple $(X_i)_{1 \leq i \leq m}$ that minimizes the $\tau$-functional
  among all tuples of $G$-valued random variables.
\end{definition}

\begin{proposition}[Existence of $\tau$-minimizer]\label{tau-min-exist-multi}\lean{multiTau_continuous, multiTau_min_exists}\leanok  If $G$ is finite, then a $\tau$-minimizer exists.
\end{proposition}

\begin{proof}\uses{tau-def-multi}\leanok This is similar to the proof of \Cref{tau-min}.
\end{proof}

\begin{proposition}[Minimizer close to reference variables]\label{tau-ref}\lean{multiTau_min_sum_le}\leanok  If $(X_i)_{1 \leq i \leq m}$ is a $\tau$-minimizer, then $\sum_{i=1}^m d[X_i; X^0] \leq \frac{2m}{\eta} d[X^0; X^0]$.
\end{proposition}

\begin{proof}\uses{tau-min-multi, multidist-nonneg, multidist-ruzsa-III}\leanok  By \Cref{tau-min-multi} we have
  $$ \tau[ (X_i)_{1 \leq i \leq m}] \leq \tau[ (X^0)_{1 \leq i \leq m}]$$
and hence by \Cref{tau-def-multi} and \Cref{multidist-nonneg}
$$ \eta \sum_{i=1}^m d[X_i; X^0] \leq D[(X^0)_{1 \leq i \leq m}] + m d[X^0; X^0].$$
The claim now follows from \Cref{multidist-ruzsa-III}.
\end{proof}

\begin{lemma}[Lower bound on multidistance]\label{multidist-lower}\lean{sub_multiDistance_le}\leanok  If  $(X_i)_{1 \leq i \leq m}$ is a $\tau$-minimizer, and $k := D[(X_i)_{1 \leq i \leq m}]$, then for any other tuple $(X'_i)_{1 \leq i \leq m}$, one has
  $$ k - D[(X'_i)_{1 \leq i \leq m}] \leq \eta \sum_{i=1}^m d[X_i; X'_i].$$
\end{lemma}

\begin{proof}\uses{tau-min-multi, ruzsa-triangle}\leanok
  By \Cref{tau-min-multi} we have
  $$ \tau[ (X_i)_{1 \leq i \leq m}] \leq \tau[ (X'_i)_{1 \leq i \leq m}]$$
  and hence by \Cref{tau-def-multi}
  $$ k + \eta \sum_{i=1}^m d[X_i; X^0] \leq D[(X'_i)_{1 \leq i \leq m}] + \eta \sum_{i=1}^m d[X'_i; X^0].$$
  On the other hand, by \Cref{ruzsa-triangle} we have
  $$ d[X'_i; X^0] \leq d[X_i; X^0] + d[X_i; X'_i].$$
  The claim follows.
\end{proof}

\begin{definition}[Conditional multidistance]\label{cond-multidist-def}\uses{multidist-def}\lean{condMultiDist}
  \leanok If $X_{[m]} = (X_i)_{1 \leq i \leq m}$ and $Y_{[m]} = (Y_i)_{1 \leq i \leq m}$ are tuples of random variables, with the $X_i$ being $G$-valued (but the $Y_i$ need not be), then we define
  \begin{equation}\label{multi-def-cond-alt}
    D[ X_{[m]} | Y_{[m]} ] = \sum_{(y_i)_{1 \leq i \leq m}} \biggl(\prod_{1 \leq i \leq m} p_{Y_i}(y_i)\biggr) D[ (X_i \,|\, Y_i \mathop{=}y_i)_{1 \leq i \leq m}]
  \end{equation}
  where each $y_i$ ranges over the support of $p_{Y_i}$ for $1 \leq i \leq m$.
\end{definition}

\begin{lemma}[Alternate form of conditional multidistance]\label{cond-multidist-alt}\lean{condMultiDist_eq}\leanok
If the $(X_i,Y_i)$ are independent,
  \begin{equation}\label{multi-def-cond}
  D[ X_{[m]} | Y_{[m]}] := \bbH[\sum_{i=1}^m X_i \big| (Y_j)_{1 \leq j \leq m} ] - \frac{1}{m} \sum_{i=1}^m \bbH[ X_i | Y_i].
    \end{equation}
\end{lemma}

\begin{proof}\uses{conditional-entropy-def, multidist-def, cond-multidist-def}\leanok
  This is routine from \Cref{conditional-entropy-def} and Definitions \ref{multidist-def} and \ref{cond-multidist-def}.
\end{proof}

\begin{lemma}[Conditional multidistance nonnegative]\label{cond-multidist-nonneg}\uses{cond-multidist-def}\lean{condMultiDist_nonneg}\leanok If $X_{[m]} = (X_i)_{1 \leq i \leq m}$ and $Y_{[m]} = (Y_i)_{1 \leq i \leq m}$ are tuples of random variables, then $D[ X_{[m]} | Y_{[m]} ] \geq 0$.
\end{lemma}

\begin{proof}\uses{multidist-nonneg}\leanok Clear from \Cref{multidist-nonneg} and \Cref{cond-multidist-def}, except that some care may need to be taken to deal with the $y_i$ where $p_{Y_i}$ vanish.
\end{proof}


\begin{lemma}[Lower bound on conditional multidistance]\label{cond-multidist-lower}\lean{sub_condMultiDistance_le}\leanok  If  $(X_i)_{1 \leq i \leq m}$ is a $\tau$-minimizer, and $k := D[(X_i)_{1 \leq i \leq m}]$, then for any other tuples $(X'_i)_{1 \leq i \leq m}$ and $(Y_i)_{1 \leq i \leq m}$ with the $X'_i$ $G$-valued, one has
  $$ k - D[(X'_i)_{1 \leq i \leq m} | (Y_i)_{1 \leq i \leq m}] \leq \eta \sum_{i=1}^m d[X_i; X'_i|Y_i].$$
\end{lemma}

\begin{proof}\uses{multidist-lower, cond-multidist-def, cond-multidist-alt}\leanok
  Immediate from \Cref{multidist-lower}, \Cref{cond-multidist-alt}, and \Cref{cond-dist-def}.
\end{proof}

\begin{corollary}[Lower bound on conditional multidistance, II]\label{cond-multidist-lower-II}\lean{sub_condMultiDistance_le'}\leanok With the notation of the previous lemma, we have
  \begin{equation}\label{5.3-conv}
    k - D[ X'_{[m]} | Y_{[m]} ] \leq \eta \sum_{i=1}^m d[X_{\sigma(i)};X'_i|Y_i]
  \end{equation}
for any permutation $\sigma : \{1,\dots,m\} \rightarrow \{1,\dots,m\}$.
\end{corollary}

\begin{proof}\uses{cond-multidist-lower, multidist-perm}\leanok  This follows from \Cref{cond-multidist-lower} and \Cref{multidist-perm}.
\end{proof}

\section{The multidistance chain rule}

\begin{lemma}[Multidistance chain rule]\label{multidist-chain-rule}\lean{multiDist_chainRule}\uses{multidist-def}\leanok  Let $\pi \colon G \to H$ be a homomorphism of abelian groups and let $X_{[m]}$ be a tuple of jointly independent $G$-valued random variables.  Then $D[X_{[m]}]$ is equal to
  \begin{equation}
      D[ X_{[m]} | \pi(X_{[m]}) ]  +D[ \pi(X_{[m]}) ]  + \bbI[ \sum_{i=1}^m X_i  : \pi(X_{[m]}) \; \big| \; \pi\bigl(\sum_{i=1}^m X_i\bigr) ]
    \label{chain-eq}
  \end{equation}
  where $\pi(X_{[m]}) := (\pi(X_i))_{1 \leq i \leq m}$.
  \end{lemma}

  \begin{proof}\uses{conditional-mutual-alt, relabeled-entropy, chain-rule}\leanok For notational brevity during this proof, write $S := \sum_{i=1}^m X_i$.

    From \Cref{conditional-mutual-alt} and \Cref{relabeled-entropy}, noting that $\pi(S)$ is determined both by $S$ and by $\pi(X_{[m]})$, we have
  \begin{equation*}
   \bbI[S:\pi(X_{[m]})|\pi(S)] = \bbH[S]+\bbH[\pi(X_{[m]})]-\bbH[S,\pi(X_{[m]})]-\bbH[\pi(S)],
  \end{equation*}
  and by \Cref{chain-rule} the right-hand side is equal to
  \begin{equation*}
  \bbH[S]-\bbH[S|\pi(X_{[m]})]-\bbH[\pi(S)].
  \end{equation*}
  Therefore,
  \begin{equation}\label{chain-1}
  \bbH[S]=\bbH[S|\pi(X_{[m]})]+\bbH[\pi(S)]+\bbI[S:\pi(X_{[m]})|\pi(S)]. \end{equation}
  From a further application of \Cref{chain-rule} and \Cref{relabeled-entropy} we have
  \begin{equation}\label{chain-2}
    \bbH[X_i] = \bbH[X_i \, | \, \pi(X_i)] + \bbH[\pi(X_i)]
  \end{equation}
  for all $1 \leq i \leq m$.  Averaging~\eqref{chain-2} in $i$ and subtracting this from~\eqref{chain-1}, we obtain the claim from \Cref{multidist-def}.
  \end{proof}

  We will need to iterate the multidistance chain rule, so it is convenient to observe a conditional version of this rule, as follows.

  \begin{lemma}[Conditional multidistance chain rule]\label{multidist-chain-rule-cond}\lean{cond_multiDist_chainRule}\leanok
    Let $\pi \colon G \to H$ be a homomorphism of abelian groups.
    Let $I$ be a finite index set and let $X_{[m]}$ be a tuple of $G$-valued random variables.
    Let $Y_{[m]}$ be another tuple of random variables (not necessarily $G$-valued).
    Suppose that the pairs $(X_i, Y_i)$ are jointly independent of one another (but $X_i$ need not be independent of $Y_i$).
    Then
    \begin{align}\nonumber
        D[ X_{[m]} | Y_{[m]} ] &=  D[ X_{[m]} \,|\, \pi(X_{[m]}), Y_{[m]}] + D[ \pi(X_{[m]}) \,|\, Y_{[m]}] \\
         &\quad\qquad + \bbI[ \sum_{i=1}^m X_i : \pi(X_{[m]}) \; \big| \;  \pi\bigl(\sum_{i=1}^m X_i \bigr), Y_{[m]} ].\label{chain-eq-cond}
    \end{align}
  \end{lemma}

  \begin{proof}\uses{multidist-chain-rule}\leanok
  For each $y_i$ in the support of $p_{Y_i}$, apply \Cref{multidist-chain-rule} with $X_i$ replaced by the conditioned random variable $(X_i|Y_i=y_i)$, and the claim~\eqref{chain-eq-cond} follows by averaging~\eqref{chain-eq} in the $y_i$ using the weights $p_{Y_i}$.
  \end{proof}

  We can iterate the above lemma as follows.

  \begin{lemma}\label{multidist-chain-rule-iter}\lean{iter_multiDist_chainRule,iter_multiDist_chainRule'}\leanok  Let $m$ be a positive integer.
    Suppose one has a sequence
    \begin{equation}\label{g-seq}
      G_m \to G_{m-1} \to \dots \to G_1 \to G_0 = \{0\}
    \end{equation}
    of homomorphisms between abelian groups $G_0,\dots,G_m$, and for each $d=0,\dots,m$, let $\pi_d : G_m \to G_d$ be the homomorphism from $G_m$ to $G_d$ arising from this sequence by composition (so for instance $\pi_m$ is the identity homomorphism and $\pi_0$ is the zero homomorphism).
    Let $X_{[m]} = (X_i)_{1 \leq i \leq m}$ be a jointly independent tuple of $G_m$-valued random variables.
    Then
    \begin{equation}
      \begin{split}
        D[ X_{[m]} ] &=  \sum_{d=1}^m D[ \pi_d(X_{[m]}) \,|\, \pi_{d-1}(X_{[m]})] \\
         &\quad + \sum_{d=1}^{m-1} \bbI[ \sum_i X_i : \pi_d(X_{[m]}) \; \big| \; \pi_d\big(\sum_i X_i\big), \pi_{d-1}(X_{[m]}) ].
      \end{split}\label{chain-eq-cond'}
    \end{equation}
    In particular, by \Cref{conditional-nonneg},
    \begin{align}\nonumber
        D[ X_{[m]} ] \geq  & \sum_{d=1}^m D[ \pi_d(X_{[m]})|\pi_{d-1}(X_{[m]}) ] \\
         & + \bbI[ \sum_i X_i : \pi_1(X_{[m]}) \; \big| \; \pi_1\bigl(\sum_i X_i\bigr) ].\label{chain-eq-cond''}
    \end{align}
  \end{lemma}

  \begin{proof}\uses{multidist-chain-rule-cond, conditional-nonneg}\leanok
  From \Cref{multidist-chain-rule-cond} (taking $Y_{[m]} = \pi_{d-1}(X_{[m]})$ and $\pi = \pi_d$ there, and noting that $\pi_d(X_{[m]})$ determines $Y_{[m]}$) we have
  \begin{align*}
    D[ X_{[m]} \,|\, \pi_{d-1}(X_{[m]}) ] &=  D[ X_{[m]} \,|\, \pi_d(X_{[m]}) ] + D[  \pi_d(X_{[m]})\,|\,\pi_{d-1}(X_{[m]}) ] \\
                                             &\quad + \bbI[ \sum_{i=1}^m X_i : \pi_d(X_{[m]}) \; \big| \; \pi_d\bigl(\sum_{i=1}^m X_i\bigr), \pi_{d-1}(X_{[m]}) ]
  \end{align*}
  for $d=1,\dots,m$. The claim follows by telescoping series, noting that $D[X_{[m]} | \pi_0(X_{[m]})] = D[X_{[m]}]$ and that $\pi_m(X_{[m]})=X_{[m]}$ (and also $\pi_m( \sum_i X_i ) = \sum_i X_i$).
  \end{proof}

  In our application we will need the following special case of the above lemma.

  \begin{corollary}\label{cor-multid}\lean{cor_multiDist_chainRule}\leanok Let $G$ be an abelian group and let $m \geq 2$.  Suppose that $X_{i,j}$, $1 \leq i, j \leq m$, are independent $G$-valued random variables.
    Then
    \begin{align*}
      &\bbI[ \bigl(\sum_{i=1}^m X_{i,j}\bigr)_{j =1}^{m} : \bigl(\sum_{j=1}^m X_{i,j}\bigr)_{i = 1}^m \; \big| \; \sum_{i=1}^m \sum_{j = 1}^m  X_{i,j} ] \\
      &\quad \leq \sum_{j=1}^{m-1} \Bigl(D[(X_{i, j})_{i = 1}^m] - D[ (X_{i, j})_{i = 1}^m  \; \big| \; (X_{i,j} + \cdots + X_{i,m})_{i =1}^m ]\Bigr) \\ & \qquad\qquad\qquad\qquad +  D[(X_{i,m})_{i=1}^m] - D[ \bigl(\sum_{j=1}^m X_{i,j}\bigr)_{i=1}^m ],
    \end{align*}
  where all the multidistances here involve the indexing set $\{1,\dots, m\}$.
  \end{corollary}

  \begin{proof}\uses{multidist-chain-rule-iter, add-entropy}\leanok
    In \Cref{multidist-chain-rule-iter} we take $G_d := G^d$ with the maps $\pi_d \colon G^m \to G^d$ for $d=1,\dots,m$ defined by
  \[
    \pi_d(x_1,\dots,x_m) := (x_1,\dots,x_{d-1}, x_d + \cdots + x_m)
  \]
  with $\pi_0=0$. Since $\pi_{d-1}(x)$ can be obtained from $\pi_{d}(x)$ by applying a homomorphism, we obtain a sequence of the form~\eqref{g-seq}.

  Now we apply \Cref{multidist-chain-rule-iter} with $I = \{1,\dots, m\}$ and $X_i :=  (X_{i,j})_{j = 1}^m$.  Using joint independence and \Cref{add-entropy}, we find that
  \[
    D[ X_{[m]} ] = \sum_{j=1}^m D[ (X_{i,j})_{1 \leq i \leq m} ].
  \]
  On the other hand, for $1 \leq j \leq m-1$, we see that once $\pi_{j}(X_i)$ is fixed, $\pi_{j+1}(X_i)$ is determined by $X_{i, j}$ and vice versa, so
  \[
    D[ \pi_{j+1}(X_{[m]}) \; | \; \pi_{j}(X_{[m]}) ] = D[ (X_{i, j})_{1 \leq i \leq m} \; | \; \pi_{j}(X_{[m]} )].
  \]
  Since the $X_{i,j}$ are jointly independent, we may further simplify:
  \[
    D[ (X_{i, j})_{1 \leq i \leq m} \; | \; \pi_{j}(X_{[m]})] = D[ (X_{i,j})_{1 \leq i \leq m} \; | \; ( X_{i, j} + \cdots + X_{i, m})_{1 \leq i \leq m} ].
  \]
  Putting all this into the conclusion of \Cref{multidist-chain-rule-iter}, we obtain
  \[
    \sum_{j=1}^{m} D[ (X_{i,j})_{1 \leq i \leq m} ]
    \geq
    \begin{aligned}[t]
    &\sum_{j=1}^{m-1} D[ (X_{i,j})_{1 \leq i \leq m} \; | \; (X_{i,j} + \cdots + X_{i,m})_{1 \leq i \leq m} ] \\
    &\!\!\!+
    D[ \bigl(\sum_{j=1}^m X_{i,j}\bigr)_{1 \leq i \leq m}] \\
    &\!\!\!+\bbI[  \bigl(\sum_{i=1}^m X_{i,j}\bigr)_{j =1}^{m} : \bigl(\sum_{j=1}^m X_{i,j}\bigr)_{i = 1}^m \; \big| \; \sum_{i=1}^m \sum_{j = 1}^m  X_{i,j} ]
    \end{aligned}
  \]
  and the claim follows by rearranging.
\end{proof}

\section{Bounding the mutual information}

As before, $G$ is an abelian group, and $m\geq 2$.  We let $X_{[m]} = (X_i)_{i =1}^m$ be a $\tau$-minimizer.

\begin{proposition}[Bounding mutual information]\label{key}\lean{mutual_information_le}\leanok
Suppose that $X_{i,j}$, $1 \leq i,j \leq m$, are jointly independent $G$-valued random variables, such that for each $j = 1,\dots,m$, the random variables $(X_{i,j})_{i = 1}^m$ coincide in distribution with some permutation of $X_{[m]}$.
  Write
  \[
    {\mathcal I} := \bbI[ \bigl(\sum_{i=1}^m X_{i,j}\bigr)_{j =1}^{m} : \bigl(\sum_{j=1}^m X_{i,j}\bigr)_{i = 1}^m \; \big| \; \sum_{i=1}^m \sum_{j = 1}^m  X_{i,j} ].
  \]
Then
  \begin{equation}\label{I-ineq}
    {\mathcal I} \leq m(4m+1) \eta k.
  \end{equation}
\end{proposition}


\begin{proof}\uses{cor-multid, multidist-perm, cond-multidist-lower-II, first-useful, klm-3, sumset-lower, ruzsa-triangle, compare-sums, multidist-ruzsa-II}\leanok
  For each $j \in \{1,\dots,m\}$ we call the tuple $(X_{i,j})_{i = 1}^m$ a \emph{column} and for each $i \in \{1,\dots, m\}$ we call the tuple $(X_{i,j})_{j = 1}^m$ a \emph{row}. Hence, by hypothesis, each column is a permutation of $X_{[m]} = (X_i)_{i=1}^m$.

From \Cref{cor-multid} we have
\begin{equation}\label{441} {\mathcal I} \leq \sum_{j=1}^{m-1} A_j + B,\end{equation}
where
\[
  A_j := D[ (X_{i, j})_{i = 1}^m] - D[ (X_{i, j})_{i = 1}^m  \; \big| \; (X_{i,j} + \cdots + X_{i,m})_{i =1}^m  ]
\]
and
\[
  B := D[ (X_{i,m})_{i=1}^m ] - D[ \bigl(\sum_{j=1}^m X_{i,j}\bigr)_{i=1}^m ].
\]
We first consider the $A_j$, for fixed $j \in \{1,\dots, m-1\}$.
By \Cref{multidist-perm} and our hypothesis on columns, we have
\[
  D[ (X_{i, j})_{i = 1}^m ]= D[ (X_i)_{i=1}^m ] = k.
\]
Let $\sigma = \sigma_j \colon I \to I$ be a permutation such that $X_{i,j} \equiv X_{\sigma(i)}$, and write $X'_i := X_{i,j}$ and $Y_i := X_{i,j} + \cdots + X_{i,m}$.
By \Cref{cond-multidist-lower-II}, we have
\begin{align}
  A_j & \leq \eta (\sum_{i = 1}^{m} d[X_{i,j}; X_{i, j}|X_{i, j} + \cdots + X_{i,m}]).\label{54a}
\end{align}
We similarly consider $B$.  By \Cref{multidist-perm} applied to the $m$-th column,
\[
  D[ (X_{i, m})_{i = 1}^m ] = D[X_{[m]}] = k.
\]
For $1 \leq i \leq m$, denote the sum of row $i$ by
\[
  V_i := \sum_{j=1}^m X_{i,j};
\]
if we apply \Cref{cond-multidist-lower-II} again, now with $X_{\sigma(i)} = X_{i,m}$, $X'_i := V_i$, and with the variable $Y_i$ being trivial, we obtain
\begin{equation}\label{55a}
  B \leq \eta \sum_{i = 1}^m d[X_{i,m}; V_i].
\end{equation}

It remains to bound the distances appearing in~\eqref{54a} and~\eqref{55a} further using Ruzsa calculus.
For $1 \leq j \leq m-1$ and $1 \leq i \leq m$, by \Cref{first-useful} we have
\begin{align*} &d[ X_{i,j}; X_{i,j}| X_{i,j}+\cdots+X_{i,m}]
\leq d[X_{i,j}; X_{i,j}] \\
&\quad + \tfrac{1}{2} \bigl(\bbH[X_{i,j}+\cdots+X_{i,m}] - \bbH[X_{i,{j+1}}+\cdots+X_{i,m}]\bigr).
\end{align*}
For each $i$, summing over $j = 1,\dots, m-1$ gives
\begin{align}
  \nonumber
  &\sum_{j=1}^{m-1} d[X_{i,j}; X_{i,j}| X_{i,j}+\cdots+X_{i,m}] \\
  &\qquad \leq \sum_{j=1}^{m-1} d[X_{i,j}; X_{i,j}] + \frac12 \bigl( \bbH[V_i] - \bbH[X_{i,m}] \bigr).
  \label{eq:distbnd1}
\end{align}
On the other hand, by \Cref{klm-3} (since $X_{i,m}$ appears in the sum $V_i$) we have
\begin{align}
  d[X_{i,m}; V_i]
  &\leq d[X_{i,m}; X_{i,m}] + \frac12 \bigl( \bbH[V_i] - \bbH[X_{i,m}] \bigr).
  \label{eq:distbnd2}
\end{align}
Combining~\eqref{441},~\eqref{54a} and~\eqref{55a} with~\eqref{eq:distbnd1} and~\eqref{eq:distbnd2} (the latter two summed over $i$), we get
\begin{align}
  \nonumber
  \frac1{\eta} {\mathcal I} &\leq \sum_{i,j=1}^m d[X_{i,j};X_{i,j}] + \sum_{i=1}^m (\bbH[V_i] - \bbH[X_{i,m}]) \\
      &= m \sum_{i=1}^m d[X_i; X_i] + \sum_{i=1}^m \bbH[V_i] - \sum_{i=1}^m \bbH[X_i].
      \label{eq:distbnd3}
\end{align}
By \Cref{compare-sums} (with $f$ taking each $j$ to the index $j'$ such that $X_{i,j}$ is a copy of $X_{j'}$) we obtain the bound
\[
  \bbH[V_i] \leq \bbH[\sum_{j=1}^m X_j] + \sum_{j=1}^m d[X_{i,j}; X_{i,j}].
\]
Finally, summing over $i$ and using $D[X_{[m]}] = k$ gives
\begin{align*}
  \sum_{i=1}^m \bbH[V_i] - \sum_{i=1}^m \bbH[X_i] & \leq \sum_{i,j=1}^m d[X_{i,j}; X_{i,j}] + m k \\ & = m\sum_{i = 1}^m d[X_i; X_i] + mk,
\end{align*}
where in the second step we used the permutation hypothesis. Combining this with~\eqref{eq:distbnd3} gives the
$$
{\mathcal I} \leq 2\eta m \biggl( \sum_{i=1}^m d[X_i;X_i] \biggr) + mk.$$
The claim \eqref{I-ineq} is now immediate from \Cref{multidist-ruzsa-II}.
\end{proof}

\section{Endgame}

Now let $m \geq 2$, let $G$ be an $m$-torsion abelian group, and let $(X_i)_{1 \leq i \leq m}$ be a $\tau$-minimizer.

\begin{definition}[Additional random variables]\label{more-random}\leanok  By a slight abuse of notation, we identify $\Z/m\Z$ and $\{1,\dots,m\}$ in the obvious way, and let $Y_{i,j}$ be an independent copy of $X_i$ for $i,j \in \Z/m\Z$.  Then also define:
  \[
    W := \sum_{i,j \in \Z/m\Z} Y_{i,j}
  \]
  and
  \[
    Z_1 := \sum_{i,j \in \Z/m\Z} i Y_{i,j},\ \ \
    Z_2 := \sum_{i,j \in \Z/m\Z} j Y_{i,j},\ \ \
    Z_3 := \sum_{i,j \in \Z/m\Z} (-i-j) Y_{i,j}.
  \]
  The addition $(-i-j)$ takes place over $\Z/m\Z$.
  Note that, because we are assuming $G$ is $m$-torsion, it is well-defined to multiply elements of $G$ by elements of $\Z/m\Z$.
  We will also define for $i,j,r \in \Z/m\Z$ the variables
  \begin{equation}\label{pqr-defs}
    P_i := \sum_{j \in \Z/m\Z} Y_{i,j} , \quad
    Q_j := \sum_{i \in \Z/m\Z} Y_{i,j} , \quad
    R_r := \sum_{\substack{i,j \in \Z/m\Z \\ i+j=-r}} Y_{i,j} .\end{equation}
\end{definition}

\begin{lemma}[Zero-sum]\label{Zero-sum}\lean{sum_of_z_eq_zero}\leanok We have
  \begin{equation}%
    \label{eq:sum-zero}
    Z_1+Z_2+Z_3= 0
  \end{equation}
\end{lemma}

\begin{proof}\uses{more-random}\leanok  Clear from definition.
\end{proof}

\begin{proposition}[Mutual information bound]\label{prop:52}\lean{mutual_information_le_t_12, mutual_information_le_t_13, mutual_information_le_t_23}\uses{more-random}\leanok
  We have
  \[
    \bbI[Z_1 : Z_2\, |\, W],\
    \bbI[Z_2 : Z_3\, |\, W],\
    \bbI[Z_1 : Z_3\, |\, W] \leq t
  \]
  where
  \begin{equation}\label{t-def}
    t :=  m(4m+1) \eta k.
  \end{equation}
\end{proposition}

\begin{proof}\uses{key, data-process}\leanok
  We analyze these variables by \Cref{key} in several different ways.
  In the first application, take $X_{i,j}=Y_{i,j}$.
  Note that each column $(X_{i,j})_{i=1}^m$ is indeed a permutation of $X_1,\dots,X_m$; in fact, the trivial permutation.
  Note also that for each $i \in \Z/m\Z$, the row sum is
  \[
    \sum_{j=1}^m X_{i,j} = \sum_{j \in \Z/m\Z} Y_{i,j} = P_i
  \]
  and for each $j \in \Z/m\Z$, the column sum is
  \[
    \sum_{i=1}^m X_{i,j} = \sum_{i \in \Z/m\Z} Y_{i,j} = Q_j.
  \]
  Finally note that $\sum_{i,j=1}^m X_{i,j} = W$.
  From \Cref{key} we then have
  \[
    \bbI[ (P_i)_{i \in \Z/m\Z} : (Q_j)_{j \in \Z/m\Z } \,\big|\, W ] \leq t,
  \]
  with $t$ as in~\eqref{t-def}.
  Since $Z_1$ is a function of $(P_i)_{i \in \Z/m\Z}$ by~\eqref{pqr-defs}, and similarly $Z_2$ is a function of $(Q_j)_{j \in \Z/m\Z}$, it follows immediately from \Cref{data-process} that
  \[
    \bbI[ Z_1 : Z_2 \,|\, W ] \leq t.
  \]

  In the second application of \Cref{key}, we instead consider $X'_{i,j} = Y_{i-j,j}$.
  Again, for each fixed $j$, the tuple $(X'_{i,j})_{i=1}^m$ is a permutation of $X_1,\dots,X_m$.
  This time the row sums for $i \in \{1,\dots, m\}$ are
  \[
    \sum_{j=1}^m X'_{i,j} = \sum_{j \in \Z/m\Z} Y_{i-j,j} = R_{-i}.
  \]
 Similarly, the column sums for $j \in \{1,\dots, m\}$ are
  \[
    \sum_{i=1}^m X'_{i,j} = \sum_{i \in \Z/m\Z} Y_{i-j,j} = Q_j.
  \]
  As before, $\sum_{i,j=1}^m X'_{i,j} = W$.
  Hence, using~\eqref{pqr-defs} and \Cref{data-process} again, \Cref{key} tells us
  \[
    \bbI[ Z_3 :  Z_2 \,|\, W] \leq \bbI[ (R_i)_{i \in \Z/m\Z} : (Q_j)_{j \in \Z/m\Z } \,\big|\, W ] \leq t.
  \]

  In the third application\footnote{In fact, by permuting the variables $(Y_{i,j})_{i,j \in \Z/m\Z}$, one can see that the random variables $(W, Z_1, Z_2)$ and $(W, Z_1, Z_3)$ have the same distribution, so this is in some sense identical to -- and can be deduced from -- the first application.} of \Cref{key}, take $X''_{i,j} = Y_{i,j-i}$.
  The column and row sums are respectively
  \[
    \sum_{j=1}^m X''_{i,j} = \sum_{j \in \Z/m\Z} Y_{i,j-i} = P_i \] and
\[     \sum_{i=1}^m X''_{i,j} = \sum_{i \in \Z/m\Z} Y_{i,j-i} = R_{-j}.
  \]
  Hence, \Cref{key} and \Cref{data-process} give
  \[
    \bbI[ Z_1 : Z_3 \,|\, W] \leq \bbI[ (P_i)_{i \in \Z/m\Z} : (R_j)_{j \in \Z/m\Z } \,\big|\, W ] \leq t,
  \]
  which completes the proof.
\end{proof}

\begin{lemma}[Entropy of $W$]\label{ent-w}\lean{entropy_of_W_le}\uses{more-random}\leanok  We have $\bbH[W] \leq (2m-1)k + \frac1m \sum_{i=1}^m \bbH[X_i]$.
\end{lemma}

\begin{proof}\uses{multidist-def, multidist-ruzsa-IV, klm-1}\leanok Without loss of generality, we may take $X_1,\dots,X_m$ to be independent. Write $S = \sum_{i=1}^m X_i$.
  Note that for each $j \in \Z/m\Z$, the sum $Q_j$ from~\eqref{pqr-defs} above has the same distribution as $S$.
  By \Cref{klm-1} we have
  \begin{align*}
    \bbH[W] = \bbH[\sum_{j \in \Z/m\Z} Q_j]  & \leq \bbH[S] + \sum_{j=2}^m (\bbH[Q_1+Q_j] - \bbH[S]) \\ & = \bbH[S] + (m-1) d[S;-S].
  \end{align*}
  By \Cref{multidist-ruzsa-IV}, we have
  \begin{equation}
    \label{eq:s-bound}
    d[S; -S] \leq 2 k
  \end{equation}
  and hence
  \[
    \bbH[W] \leq 2 k (m-1) + \bbH[S].
  \]
  From \Cref{multidist-def} we have
  \begin{equation}
    \label{eq:ent-s}
    \bbH[S] = k + \frac1m \sum_{i=1}^m \bbH[X_i],
  \end{equation}
  and the claim follows.
\end{proof}

\begin{lemma}[Entropy of $Z_2$]\label{ent-z2}\lean{entropy_of_Z_two_le}\uses{more-random}\leanok  We have $\bbH[Z_2] \leq (8m^2-16m+1) k + \frac{1}{m} \sum_{i=1}^m \bbH[X_i]$.
\end{lemma}

\begin{proof}\uses{sum-dilate-II, klm-1}\leanok
  We observe
  \[
    \bbH[Z_2] = \bbH[\sum_{j \in \Z/m\Z} j Q_j].
  \]
  Applying \Cref{klm-1} one has
  \begin{align*}
    \bbH[Z_2] &\leq \sum_{i=2}^{m-1} \bbH[Q_1 + i Q_i]  - (m-2) \bbH[S].
  \end{align*}
  Using \Cref{sum-dilate-II} and~\eqref{eq:s-bound} we get
  \begin{align*}
    \bbH[Z_2]
              &\leq \bbH[S] + 4m (m-2) d[S;-S] \\
              &\leq \bbH[S] + 8m (m-2) k.
  \end{align*}
  Applying~\eqref{eq:ent-s} gives the claim.
\end{proof}

\begin{lemma}[Mutual information bound]\label{mutual-w-z2}\lean{mutual_of_W_Z_two_le}\leanok  We have  $\bbI[W : Z_2] \leq 2 (m-1) k$.
\end{lemma}

\begin{proof}\uses{alternative-mutual, ent-w}\leanok
  From \Cref{alternative-mutual} we have $\bbI[W : Z_2] = \bbH[W] - \bbH[W | Z_2]$, and since $Z_2 = \sum_{j=1}^{m-1} j Q_j$ and $W = \sum_{j=1}^m Q_j$,
  \[
    \bbH[W | Z_2] \geq \bbH[W \,|\, Q_1,\dots,Q_{m-1}] = \bbH[Q_m] = \bbH[S].
  \]
  Hence, by \Cref{ent-w},
  \[
    \bbI[W : Z_2] \leq \bbH[W] - \bbH[S] \leq 2 (m-1) k,
  \]
  as claimed.
\end{proof}

\begin{lemma}[Distance bound]\label{xi-z2-w-dist}\lean{sum_of_conditional_distance_le}\leanok We have $\sum_{i=1}^m d[X_i;Z_2|W] \leq 4(m^3-m^2) k$.
\end{lemma}

\begin{proof}\uses{klm-3, cond-dist-fact, mutual-w-z2, ent-z2, multidist-ruzsa-II}\leanok
  For each $i \in \{1,\dots, m\}$, using \Cref{klm-3} (noting the sum $Z_2$ contains $X_i$ as a summand) we have
  \begin{equation}\label{in-a-bit-6}
    d[X_i;Z_2] \leq d[X_i;X_i] + \tfrac12 (\bbH[Z_2] - \bbH[X_i])
  \end{equation}
  and using \Cref{cond-dist-fact} we have
  \[
    d[X_i;Z_2 | W] \leq d[X_i;Z_2] + \tfrac12 \bbI[W : Z_2].
  \]
 Combining with~\eqref{in-a-bit-6} and \Cref{mutual-w-z2} gives
 \[ d[X_i;Z_2 | W] \leq d[X_i;X_i] + \tfrac12 (\bbH[Z_2] - \bbH[X_i]) + (m-1)k.\]
 Summing over $i$ and applying \Cref{ent-z2} gives
 \[ \sum_{i = 1}^m d[X_i;Z_2 | W] \leq \sum_{i = 1}^m d[X_i;X_i] + m(8m^2-16m+1) k/2 + m(m-1) k.\]
Finally, applying \Cref{multidist-ruzsa-II} (and dropping some lower order terms) gives the claim.
\end{proof}

\begin{lemma}[Application of BSG]
  \label{lem:get-better}\lean{dist_of_U_add_le}\leanok
  Let $G$ be an abelian group, let $(T_1,T_2,T_3)$ be a $G^3$-valued random variable such that $T_1+T_2+T_3=0$ holds identically, and write
  \[
    \delta := \bbI[T_1 : T_2] + \bbI[T_1 : T_3] + \bbI[T_2 : T_3].
  \]
  Let $Y_1,\dots,Y_n$ be some further $G$-valued random variables and let $\alpha>0$ be a constant.
  Then there exists a random variable $U$ such that
  \begin{equation}
    \label{eq:get-better}
    d[U;U] + \alpha \sum_{i=1}^n d[Y_i;U] \leq \Bigl(2 + \frac{\alpha n}{2} \Bigr) \delta + \alpha \sum_{i=1}^n d[Y_i;T_2].
  \end{equation}
\end{lemma}

\begin{proof}\uses{entropic-bsg, relabeled-entropy, first-useful}\leanok
  We apply \Cref{entropic-bsg} with $X=T_1$ and $Y=T_2$.
  Since $T_1+T_2=-T_3$, we find that
  \begin{equation}\label{514a}
  \begin{split}
    \sum_{z} p_{T_3}(z) & d[T_1 \,|\, T_3 \mathop{=} z;T_2 \,|\, T_3 \mathop{=} z] \\
    &\leq  3 \bbI[T_1 : T_2] + 2 \bbH[T_3] - \bbH[T_1] - \bbH[T_2] \\
    &\ = \bbI[T_1 : T_2] + \bbI[T_1 : T_3] + \bbI[T_2 : T_3]
    = \delta,
  \end{split}
\end{equation}
  where the last line follows from \Cref{relabeled-entropy} by observing
  \[
    \bbH[T_1,T_2] = \bbH[T_1,T_3] = \bbH[T_2,T_3] = \bbH[T_1,T_2,T_3]
  \]
  since any two of $T_1,T_2,T_3$ determine the third.

  By~\eqref{514a} and the triangle inequality,
  \[
    \sum_z p_{T_3}(z) d[T_2 \,|\, T_3 \mathop{=} z; T_2 \,|\, T_3\mathop{=}z] \leq 2 \delta
  \]
  and by \Cref{first-useful}, for each $Y_i$,
  \begin{align*}
    &\sum_z p_{T_3}(z) d[Y_i; T_2 \,|\, T_3 \mathop{=} z] \\
    &\qquad= d[Y_i; T_2 \,|\, T_3]
                                    \leq d[Y_i;T_2] + \frac12 \bbI[T_2 : T_3]
                                    \leq d[Y_i;T_2] + \frac{\delta}{2}.
  \end{align*}
  Hence,
  \begin{align*}
    &\sum_z p_{T_3}(z) \bigg( d[T_2 \,|\, T_3 \mathop{=} z; T_2 \,|\, T_3 \mathop{=} z] + \alpha \sum_{i=1}^n d[Y_i;T_2 \,|\, T_3 \mathop{=} z] \bigg) \\
    &\qquad \leq
    \Bigl(2 + \frac{\alpha n}{2} \Bigr) \delta + \alpha \sum_{i=1}^n d[Y_i; T_2],
  \end{align*}
  and the result follows by setting $U=(T_2 \,|\, T_3 \mathop{=} z)$ for some $z$ such that the quantity in parentheses on the left-hand side is at most the weighted average value.
\end{proof}

\begin{proposition}[Vanishing entropy]\label{k-vanish}\lean{k_eq_zero}\leanok  We have $k = 0$.
\end{proposition}

\begin{proof}\uses{lem:get-better, multidist-ruzsa-III, multidist-lower, xi-z2-w-dist, eta-def-multi, multidist-nonneg, Zero-sum, prop:52}\leanok For each value $W=w$, apply \Cref{lem:get-better} (and \Cref{Zero-sum}) to
  \[
    T_1 = (Z_1 \,|\, W \mathop{=} w),\
    T_2 = (Z_2 \,|\, W \mathop{=} w),\
    T_3 = (Z_3 \,|\, W \mathop{=} w)
  \]
  with $Y_i=X_i$ and $\alpha=\eta/m$.  Write
  \[
    \delta_w :=  \bbI[T_1 : T_2] + \bbI[T_1 : T_3] + \bbI[T_2 : T_3]
  \]
  for this choice, and note that
  \begin{align}
    \label{eq:delta-w}
    \nonumber
    \delta_\ast := \sum_w p_W(w) \delta_w &= \bbI[Z_1 : Z_2 \,|\, W] + \bbI[Z_1 : Z_3 \,|\, W] + \bbI[Z_2 : Z_3 \,|\, W] \\
                           & \leq 3 m (4m+1) \eta k
  \end{align}
  by \Cref{prop:52}.
  Write $U_w$ for the random variable guaranteed to exist by \Cref{lem:get-better},
  so that~\eqref{eq:get-better} gives \begin{equation}
    \label{eq:uw1}
    d[U_w; U_w] \leq \Bigl( 2 + \frac{\alpha m}{2} \Bigr) \delta_w + \alpha \sum_{i=1}^m \bigl( d[X_i; T_2] - d[X_i; U_w] \bigr).
  \end{equation}
  Let $(U_w)_I$ denote the tuple consisting of the same variable $U_w$ repeated $m$ times.
  By \Cref{multidist-ruzsa-III}
  \begin{equation}
    \label{eq:uw2}
    D[(U_w)_I] \leq m d[U_w; U_w].
  \end{equation}
  On the other hand, from \Cref{multidist-lower} one has
  \begin{equation}
    \label{eq:uw3}
    D[(U_w)_I] \geq k - \eta \sum_{i=1}^m d[X_i;U_w].
  \end{equation}
  Combining~\eqref{eq:uw1},~\eqref{eq:uw2} and~\eqref{eq:uw3} and averaging over $w$ (with weight $p_W(w)$), and recalling the value $\alpha=\eta/m$, gives
  \[
     m \Bigl( 2 + \frac{\eta}{2} \Bigr) \delta_\ast + \eta \sum_{i=1}^m d[X_i; Z_2 | W]
    \geq k
  \]
  since the terms $d[X_i; U_w]$ cancel by our choice of $\alpha$.
  Substituting in \Cref{xi-z2-w-dist} and~\eqref{eq:delta-w}, and using the fact that $2 + \frac{\eta}{2} < 3$, we have
  \[
    3 m^2 (4m+1) (2+\frac{\eta}{2}) \eta k + \eta 8(m^3-m^2) k \geq k.
  \]
  From \Cref{eta-def-multi} we have we have
  $$ 3 m^2 (4m+1) (2+\frac{\eta}{2}) \eta + \eta 8(m^3-m^2) < 1$$
  and hence $k \leq 0$.  The claim now follows from \Cref{multidist-nonneg}.
\end{proof}

\section{Wrapping up}

\begin{theorem}[Entropy form of PFR]\label{main-entropy}\lean{dist_of_X_U_H_le}\leanok Suppose that $G$ is a finite abelian group of torsion $m$.  Suppose that $X$ is a $G$-valued random variable. Then there exists a subgroup $H \leq G$ such that \[ d[X;U_H] \leq 64 m^3 d[X;X].\]
\end{theorem}

\begin{proof}\uses{k-vanish, ruzsa-triangle, tau-def, multi-zero, eta-def-multi, tau-ref, tau-min-exist-multi}\leanok  Set $X^0 := X$. By \Cref{tau-min-exist-multi}, there exists a $\tau$-minimizer $X_{[m]} = (X_i)_{1 \leq i \leq m}$.  By \Cref{k-vanish}, we have $D[X_{[m]}]=0$. By \Cref{tau-ref} and the pigeonhole principle, there exists $1 \leq i \leq m$ such that $d[X_i; X] \leq \frac{2}{\eta} d[X;X]$.  By \Cref{multi-zero}, we have $d[X_i;U_H]=0$ for some subgroup $H \leq G$, hence by \Cref{ruzsa-triangle} we have $d[U_H; X] \leq \frac{2}{\eta} d[X;X]$. The claim then follows from \Cref{eta-def-multi}.
\end{proof}

\begin{lemma}\label{pfr_aux_torsion}\lean{torsion_PFR_conjecture_aux}\leanok  Suppose that $G$ is a finite abelian group of torsion $m$.  If $A \subset G$ is non-empty and
  $|A+A| \leq K|A|$, then $A$ can be covered by at most $K ^
  {128m^3+1}|A|^{1/2}/|H|^{1/2}$ translates of a subspace $H$ of $G$ with
  \begin{equation}
    \label{ah2}
    |H|/|A| \in [K^{-256m^3}, K^{256m^3}].
  \end{equation}
  \end{lemma}

\begin{proof}\uses{main-entropy, unif-exist, uniform-entropy-II, jensen-bound,ruz-dist-def,ruzsa-diff,bound-conc,ruz-cov}\leanok  Repeat the proof of \Cref{pfr_aux}, but with \Cref{main-entropy} in place of \Cref{entropy-pfr}.  Because of the lack of $2$-torsion, one has to use the Ruzsa triangle inequality to bound $d[U,U]$ by $2d[U,-U]$, costing an additional factor of $2$ in the estimates.
\end{proof}

\begin{theorem}[PFR]\label{pfr-torsion}\lean{torsion_PFR}\leanok  Suppose that $G$ is an abelian group of torsion $m$ (not necessarily finite).
  If $A \subset G$ is a finite non-empty set with $|A+A| \leq K|A|$, then $A$ can be covered by at most $mK^{256m^3+1}$ translates of a finite subspace $H$ of $G$ with $|H| \leq |A|$.
  \end{theorem}

  \begin{proof}\uses{pfr_aux_torsion}\leanok If $G$ is finite, repeat the proof of \Cref{pfr}, but with \Cref{pfr_aux_torsion} in place of \Cref{pfr_aux}. In general, reduce to the finite span of $A$ over $\mathrm{Z}/m\mathrm{Z}$.
  \end{proof}


%%FILE blueprint/src/chapter/weak_pfr.tex
\chapter{Weak PFR over the integers}

\begin{lemma}\label{torsion-free-doubling}\lean{torsion_free_doubling}\leanok
If $G$ is torsion-free and $X,Y$ are $G$-valued random variables then $d[X;2Y]\leq 5d[X;Y]$.
\end{lemma}
\begin{proof}
\uses{alt-submodularity,ruzsa-triangle,ruzsa-diff}\leanok
Let $Y_1,Y_2$ be independent copies of $Y$ (also independent of $X$). Since $G$ is torsion-free we know $X,Y_1-Y_2,X-2Y_1$ uniquely determine $X,Y_1,Y_2$ and so
\[\mathbb{H}(X,Y_1,Y_2,X-2Y_1)=\mathbb{H}(X,Y_1,Y_2)=\mathbb{H}(X)+2\mathbb{H}(Y).\]
Similarly
\[\mathbb{H}(X,X-2Y_1)=\mathbb{H}(X)+\mathbb{H}(2Y_1)=\mathbb{H}(X)+\mathbb{H}(Y).\]
Furthermore
\[\mathbb{H}(Y_1-Y_2,X-2Y_1)=\mathbb{H}(Y_1-Y_2,X-Y_1-Y_2)\leq \mathbb{H}(Y_1-Y_2)+\mathbb{H}(X-Y_1-Y_2).\]
By submodularity (\Cref{alt-submodularity})
\[\mathbb{H}(X,Y_1,Y_2,X-2Y_1)+\mathbb{H}(X-2Y_1)\leq \mathbb{H}(X,X-2Y_1)+\mathbb{H}(Y_1-Y_2,X-2Y_1).\]
Combining these inequalities
\[\mathbb{H}(X-2Y_1)\leq \mathbb{H}(Y_1-Y_2)+\mathbb{H}(X-Y_1-Y_2)-\mathbb{H}(Y).\]
Similarly we have
\[\mathbb{H}(Y_1,Y_2,X-Y_1-Y_2)=\mathbb{H}(X)+2\mathbb{H}(Y),\]
\[\mathbb{H}(Y_1,X-Y_1-Y_2)=\mathbb{H}(Y)+\mathbb{H}(X-Y_2),\]
and
\[\mathbb{H}(Y_2,X-Y_1-Y_2)=\mathbb{H}(Y)+\mathbb{H}(X-Y_1)\]
and by submodularity (\Cref{alt-submodularity}) again
\[\mathbb{H}(Y_1,Y_2,X-Y_1-Y_2)+ \mathbb{H}(X-Y_1-Y_2)\leq \mathbb{H}(Y_1,X-Y_1-Y_2)+\mathbb{H}(Y_2,X-Y_1-Y_2).\]
Combining these inequalities (and recalling the definition of Ruzsa distance) gives
\[\mathbb{H}(X-Y_1-Y_2)\leq \mathbb{H}(X-Y_1)+\mathbb{H}(X-Y_2)-\mathbb{H}(X)=2d[X;Y]+\mathbb{H}(Y).\]
It follows that
\[\mathbb{H}(X-2Y_1)\leq \mathbb{H}(Y_1-Y_2)+2d[X;Y]\]
and so (using $\mathbb{H}(2Y)=\mathbb{H}(Y)$)
\begin{align*}
d[X;2Y]
&=\mathbb{H}(X-2Y_1)-\mathbb{H}(X)/2-\mathbb{H}(2Y)/2\\
&\leq \mathbb{H}(Y_1-Y_2)+2d[X;Y]-\mathbb{H}(X)/2-\mathbb{H}(Y)/2\\
&= d[Y_1;Y_2]+\frac{\mathbb{H}(Y)-\mathbb{H}(X)}{2}+2d[X;Y].
\end{align*}
Finally note that by the triangle inequality (\Cref{ruzsa-triangle}) we have
\[d[Y_1;Y_2]\leq d[Y_1;X]+d[X;Y_2]=2d[X;Y].\]
The result follows from $(\mathbb{H}(Y)-\mathbb{H}(X))/2\leq d[X;Y]$ (\Cref{ruzsa-diff}).
\end{proof}

\begin{lemma}\label{torsion-dist-shrinking}\lean{torsion_dist_shrinking}\leanok
If $G$ is a torsion-free group and $X,Y$ are $G$-valued random variables and $\phi:G\to \mathbb{F}_2^d$ is a homomorphism then
\[\mathbb{H}(\phi(X))\leq 10d[X;Y].\]
\end{lemma}
\begin{proof}
\uses{fibring-ineq, torsion-free-doubling, dist-zero}\leanok
By \Cref{fibring-ineq} and \Cref{torsion-free-doubling} we have
\[d[\phi(X);\phi(2Y)]\leq d[X;2Y]\leq 5d[X;Y]\]
and $\phi(2Y)=2\phi(Y)\equiv 0$ so the left-hand side is equal to $d[\phi(X);0]=\mathbb{H}(\phi(X))/2$ (using \Cref{dist-zero}).
\end{proof}

\begin{lemma}\label{app-ent-pfr}\lean{app_ent_PFR}\leanok
Let $G=\mathbb{F}_2^n$ and $\alpha\in (0,1)$ and let $X,Y$ be $G$-valued
random variables such that
\[\mathbb{H}(X)+\mathbb{H}(Y)> \frac{20}{\alpha} d[X;Y].\]
There is a non-trivial subgroup $H\leq G$ such that
\[\log \lvert H\rvert <\frac{1+\alpha}{2}(\mathbb{H}(X)+\mathbb{H}(Y))\] and
\[\mathbb{H}(\psi(X))+\mathbb{H}(\psi(Y))< \alpha (\mathbb{H}(X)+\mathbb{H}(Y))\]
where $\psi:G\to G/H$ is the natural projection homomorphism.
\end{lemma}
\begin{proof}
\uses{entropy-pfr-improv, ruzsa-diff, dist-projection, ruzsa-nonneg}\leanok
By \Cref{entropy-pfr-improv} there exists a subgroup $H$ such that
$d[X;U_H] + d[Y;U_H] \leq 10 d[X;Y]$. Using \Cref{dist-projection} we
deduce that $\mathbb{H}(\psi(X)) + \mathbb{H}(\psi(X)) \leq 20 d[X;Y]$. The
second claim follows adding these inequalities and using the assumption on
$\mathbb{H}(X)+\mathbb{H}(Y)$.

Furthermore we have by \Cref{ruzsa-diff}
\[\log \lvert H \rvert-\mathbb{H}(X)\leq 2d[X;U_H]\]
and similarly for $Y$ and thus
\begin{align*}
\log \lvert  H\rvert
&\leq
\frac{\mathbb{H}(X)+\mathbb{H}(Y)}{2}+d[X;U_H] + d[Y;U_H] \leq
\frac{\mathbb{H}(X)+\mathbb{H}(Y)}{2}+ 10d[X;Y]
\\& <
\frac{1+\alpha}{2}(\mathbb{H}(X)+\mathbb{H}(Y)).
\end{align*}
Finally note that if $H$
were trivial then $\psi(X)=X$ and $\psi(Y)=Y$ and hence
$\mathbb{H}(X)+\mathbb{H}(Y)=0$, which contradicts \Cref{ruzsa-nonneg}.
\end{proof}


\begin{lemma}\label{pfr-projection'}\lean{PFR_projection'}\leanok
If $G=\mathbb{F}_2^d$ and $\alpha\in (0,1)$ and $X,Y$ are $G$-valued random
variables then there is a subgroup $H\leq \mathbb{F}_2^d$ such that
\[\log \lvert H\rvert \leq \frac{1+\alpha}{2(1-\alpha)} (\mathbb{H}(X)+\mathbb{H}(Y))\]
and if $\psi:G \to G/H$ is the natural projection then
\[\mathbb{H}(\psi(X))+\mathbb{H}(\psi(Y))\leq \frac{20}{\alpha} d[\psi(X);\psi(Y)].\]
\end{lemma}
\begin{proof}
\uses{app-ent-pfr}\leanok
Let $H\leq \mathbb{F}_2^d$ be a maximal subgroup such that
\[\mathbb{H}(\psi(X))+\mathbb{H}(\psi(Y))> \frac{20}{\alpha} d[\psi(X);\psi(Y)]\]
and such that there exists $c \ge 0$ with
\[\log \lvert H\rvert \leq \frac{1+\alpha}{2(1-\alpha)}(1-c)(\mathbb{H}(X)+\mathbb{H}(Y))\]
and
\[\mathbb{H}(\psi(X))+\mathbb{H}(\psi(Y))\leq c (\mathbb{H}(X)+\mathbb{H}(Y)).\]
Note that this exists since $H=\{0\}$ is an example of such a subgroup or we are done with this choice of $H$.

We know that $G/H$ is a $2$-elementary group and so by Lemma
\ref{app-ent-pfr} there exists some non-trivial subgroup $H'\leq G/H$ such
that
\[\log \lvert H'\rvert < \frac{1+\alpha}{2}(\mathbb{H}(\psi(X))+\mathbb{H}(\psi(Y))\]
and
\[\mathbb{H}(\psi' \circ\psi(X))+\mathbb{H}(\psi' \circ \psi(Y))< \alpha(\mathbb{H}(\psi(X))+\mathbb{H}(\psi(Y)))\]
where $\psi':G/H\to (G/H)/H'$. By group isomorphism theorems we know that
there exists some $H''$ with $H\leq H''\leq G$ such that $H'\cong H''/H$ and
$\psi' \circ \psi(X)=\psi''(X)$ where $\psi'':G\to G/H''$ is the projection
homomorphism.

Since $H'$ is non-trivial we know that $H$ is a proper subgroup of $H''$. On the other hand we know that
\[\log \lvert H''\rvert=\log \lvert H'\rvert+\log \lvert H\rvert< \frac{1+\alpha}{2(1-\alpha)}(1-\alpha c)(\mathbb{H}(X)+\mathbb{H}(Y))\]
and
\[\mathbb{H}(\psi''(X))+\mathbb{H}(\psi''(Y))< \alpha (\mathbb{H}(\psi(X))+\mathbb{H}(\psi(Y)))\leq \alpha c (\mathbb{H}(X)+\mathbb{H}(Y)).\]
Therefore (using the maximality of $H$) it must be the first condition that fails, whence
\[\mathbb{H}(\psi''(X))+\mathbb{H}(\psi''(Y))\leq \frac{20}{\alpha}d[\psi''(X);\psi''(Y)].\]
\end{proof}

We could use the previous lemma for any value of $\alpha \in (0,1)$, which
would give a whole range of estimates in Theorem~\ref{weak-pfr-int}. For
definiteness, we specialize only to $\alpha=3/5$, which gives a constant $2$
in the first bound below.

\begin{lemma}\label{pfr-projection}\lean{PFR_projection}\leanok
If $G=\mathbb{F}_2^d$ and $\alpha\in (0,1)$ and $X,Y$ are $G$-valued random
variables then there is a subgroup $H\leq \mathbb{F}_2^d$ such that
\[\log \lvert H\rvert \leq 2 (\mathbb{H}(X)+\mathbb{H}(Y))\]
and if $\psi:G \to G/H$ is the natural projection then
\[\mathbb{H}(\psi(X))+\mathbb{H}(\psi(Y))\leq 34 d[\psi(X);\psi(Y)].\]
\end{lemma}
\begin{proof}
\uses{pfr-projection'}\leanok
Specialize \Cref{pfr-projection'} to $\alpha=3/5$. In the second
inequality, it gives a bound $100/3 < 34$.
\end{proof}



\begin{lemma}\label{single-fibres}\lean{single_fibres}\leanok
Let $\phi:G\to H$ be a homomorphism and $A,B\subseteq G$ be finite subsets. If $x,y\in H$ then let $A_x=A\cap \phi^{-1}(x)$ and $B_y=B\cap \phi^{-1}(y)$. There exist $x,y\in H$ such that $A_x,B_y$ are both non-empty and
\[d[\phi(U_A);\phi(U_B)]\log \frac{\lvert A\rvert\lvert B\rvert}{\lvert A_x\rvert\lvert B_y\rvert}\leq (\mathbb{H}(\phi(U_A))+\mathbb{H}(\phi(U_B)))(d(U_A,U_B)-d(U_{A_x},U_{B_y})).\]
\end{lemma}
\begin{proof}
\uses{fibring-ident}\leanok
The random variables $(U_A\mid \phi(U_A)=x)$ and $(U_B\mid \phi(U_B)=y)$ are equal in distribution to $U_{A_x}$ and $U_{B_y}$ respectively (both are uniformly distributed over their respective fibres). It follows from \Cref{fibring-ident} that
\begin{align*}
\sum_{x,y\in H}\frac{\lvert A_x\rvert\lvert B_y\rvert}{\lvert A\rvert\lvert B\rvert}d[U_{A_x};U_{B_y}]
&=d[U_A\mid \phi(U_A); U_B\mid \phi(U_B)]\\
&\leq d[U_A;U_B]-d[\phi(U_A);\phi(U_B)].
\end{align*}
Therefore with $M:=\mathbb{H}(\phi(U_A))+\mathbb{H}(\phi(U_B))$ we have
\[\left(\sum_{x,y\in H}\frac{\lvert A_x\rvert\lvert B_y\rvert}{\lvert A\rvert\lvert B\rvert}Md[U_{A_x};U_{B_y}]\right)+Md[\phi(U_A);\phi(U_B)]\leq Md[U_A;U_B].\]
Since
\[M=\sum_{x,y\in H}\frac{\lvert A_x\rvert\lvert B_y\rvert}{\lvert A\rvert\lvert B\rvert}\log \frac{\lvert A\rvert\lvert B\rvert}{\lvert A_x\rvert\lvert B_y\rvert}\]
we have
\[\sum_{x,y\in H} \frac{\lvert A_x\rvert\lvert B_y\rvert}{\lvert A\rvert\lvert B\rvert}\left(Md[U_{A_x};U_{B_y}]+d[\phi(U_A);\phi(U_B)]\log \frac{\lvert A\rvert\lvert B\rvert}{\lvert A_x\rvert\lvert B_y\rvert}\right)\leq  Md[U_A;U_B].\]
It follows that there exists some $x,y\in H$ such that $\lvert A_x\rvert,\lvert B_y\rvert\neq 0$ and
\[Md[U_{A_x};U_{B_y}]+d[\phi(U_A);\phi(U_B)]\log \frac{\lvert A\rvert\lvert B\rvert}{\lvert A_x\rvert\lvert B_y\rvert}\leq  Md[U_A;U_B].\]
\end{proof}


\begin{definition}
  \label{dimension-def}
  \lean{AffineSpace.finrank}
  \leanok

  If $A\subseteq \mathbb{Z}^{d}$ then by $\dim(A)$ we mean the dimension of the span of $A-A$
  over the reals -- equivalently, the smallest $d'$ such that $A$ lies in a coset of a subgroup
  isomorphic to $\mathbb{Z}^{d'}$.
\end{definition}


\begin{theorem}\label{weak-pfr-asymm}\lean{weak_PFR_asymm}\leanok
If $A,B\subseteq \mathbb{Z}^d$ are finite non-empty sets then there exist non-empty $A'\subseteq A$ and $B'\subseteq B$ such that
\[\log\frac{\lvert A\rvert\lvert B\rvert}{\lvert A'\rvert\lvert B'\rvert}\leq 34d[U_A;U_B]\]
such that $\max(\dim A',\dim B')\leq \frac{40}{\log 2} d[U_A;U_B]$.
\end{theorem}
\begin{proof}
\uses{torsion-dist-shrinking, pfr-projection, single-fibres, dimension-def}\leanok
Without loss of generality we can assume that $A$ and $B$ are not both inside (possibly distinct) cosets of the same subgroup of $\mathbb{Z}^d$, or we just replace $\mathbb{Z}^d$ with that subgroup. We prove the result by induction on $\lvert A\rvert+\lvert B\rvert$.

Let $\phi:\mathbb{Z}^d\to \mathbb{F}_2^d$ be the natural mod-2 homomorphism. By \Cref{torsion-dist-shrinking}
\[\max(\mathbb{H}(\phi(U_A)),\mathbb{H}(\phi(U_B)))\leq 10d[U_A;U_B].\]
We now apply \Cref{pfr-projection}, obtaining some subgroup $H\leq \mathbb{F}_2^d$ such that
\[\log \lvert H\rvert \leq 40d[U_A;U_B]\]
and
\[\mathbb{H}(\tilde{\phi}(U_A))+\mathbb{H}(\tilde{\phi}(U_B))\leq 34 d[\tilde{\phi}(U_A);\tilde{\phi}(U_B)]\]
where $\tilde{\phi}:\mathbb{Z}^d\to \mathbb{F}_2^d/H$ is $\phi$ composed with the projection onto $\mathbb{F}_2^d/H$.


By \Cref{single-fibres} there exist $x,y\in \mathbb{F}_2^d/H$ such that, with $A_x=A\cap \tilde{\phi}^{-1}(x)$ and similarly for $B_y$,
\[\log \frac{\lvert A\rvert\lvert B\rvert}{\lvert A_x\rvert\lvert B_y\rvert}\leq 34(d[U_A;U_B]-d[U_{A_x};U_{B_y}]).\]
Suppose first that $\lvert A_x\rvert+\lvert B_y\rvert=\lvert A\rvert+\lvert B\rvert$. This means that $\tilde{\phi}(A)=\{x\}$ and $\tilde{\phi}(B)=\{y\}$, and hence both $A$ and $B$ are in cosets of $\ker \tilde{\phi}$. Since by assumption $A,B$ are not in cosets of a proper subgroup of $\mathbb{Z}^d$ this means $\ker\tilde{\phi}=\mathbb{Z}^d$, and so (examining the definition of $\tilde{\phi}$) we must have $H=\mathbb{F}_2^d$. Then our bound on $\log\lvert H\rvert$ forces $d\leq \frac{40}{\log 2}d[U_A;U_B]$ and we are done with $A'=A$ and $B'=B$.

Otherwise,
\[\lvert A_x\rvert+\lvert B_y\rvert <\lvert A\rvert+\lvert B\rvert.\]
By induction we can find some $A'\subseteq A_x$ and $B'\subseteq B_y$ such that $\dim A',\dim B'\leq \frac{40}{\log 2} d[U_{A_x};U_{B_y}]\leq \frac{40}{\log 2}d[U_A;U_B]$ and

\[\log \frac{\lvert A_x\rvert\lvert B_y\rvert}{\lvert A'\rvert\lvert B'\rvert}\leq 34d[U_{A_x};U_{B_y}].\]
Adding these inequalities implies

\[\log\frac{\lvert A\rvert\lvert B\rvert}{\lvert A'\rvert\lvert B'\rvert}\leq 34d[U_A;U_B]\]
as required.
\end{proof}

\begin{theorem}
  \label{weak-pfr-symm}
  \lean{weak_PFR}
  \leanok

  If $A\subseteq \mathbb{Z}^d$ is a finite non-empty set with $d[U_A;U_A]\leq \log K$
  then there exists a non-empty $A'\subseteq A$ such that
  \[\lvert A'\rvert\geq K^{-17}\lvert A\rvert\]
  and $\dim A'\leq \frac{40}{\log 2} \log K$.
\end{theorem}
\begin{proof}
\uses{weak-pfr-asymm,dimension-def}\leanok
Immediate from \Cref{weak-pfr-asymm} and rearranging.
\end{proof}

\begin{theorem}\label{weak-pfr-int}\lean{weak_PFR_int}\leanok
Let $A\subseteq \mathbb{Z}^d$ and $\lvert A-A\rvert\leq K\lvert A\rvert$.
There exists $A'\subseteq A$ such that $\lvert A'\rvert \geq K^{-17}\lvert
A\rvert$ and $\dim A' \leq \frac{40}{\log 2}\log K$.
\end{theorem}
\begin{proof}\leanok
\uses{weak-pfr-symm,dimension-def}
As in the beginning of \Cref{pfr} the doubling condition forces $d[U_A;U_A]\leq \log K$, and then we apply \Cref{weak-pfr-symm}.
\end{proof}

\begin{corollary}\label{weak-pfr-int-sumset}\lean{weak_PFR_int_sumset}\leanok
Let $A\subseteq \mathbb{Z}^d$ and $\lvert A+A\rvert\leq K\lvert A\rvert$.
There exists $A'\subseteq A$ such that $\lvert A'\rvert \geq K^{-34}\lvert
A\rvert$ and $\dim A' \leq \frac{80}{\log 2}\log K$.
\end{corollary}
\begin{proof}\leanok
\uses{weak-pfr-int}
By the Ruzsa triangle inequality one has $\lvert A-A\rvert \leq K^2 \lvert A\rvert$, and then we apply \Cref{weak-pfr-int}.
\end{proof}

